\documentclass[11pt]{article}
\usepackage[a4paper,margin=27mm]{geometry}
\usepackage[T1]{fontenc}
\usepackage[utf8]{inputenc}
\usepackage{lmodern}
\usepackage{microtype}
\usepackage{amsmath,amssymb,amsthm,mathtools,mathrsfs}
\usepackage{booktabs,longtable,array}
\usepackage{enumitem}
\usepackage{natbib}
\usepackage{xurl}
\usepackage{xcolor}
\usepackage[colorlinks=true,linkcolor=blue!45!black,citecolor=blue!45!black,urlcolor=blue!45!black]{hyperref}
\usepackage[nameinlink,noabbrev]{cleveref}

\setlength{\parindent}{0pt}
\setlength{\parskip}{0.52em}
\numberwithin{equation}{section}
\newtheorem{theorem}{Theorem}[section]
\newtheorem{proposition}[theorem]{Proposition}
\newtheorem{lemma}[theorem]{Lemma}
\newtheorem{corollary}[theorem]{Corollary}
\theoremstyle{definition}
\newtheorem{definition}[theorem]{Definition}
\newtheorem{conjecture}[theorem]{Conjecture}
\newtheorem{test}[theorem]{Open problem}
\theoremstyle{remark}
\newtheorem{remark}[theorem]{Remark}

\newcommand{\R}{\mathbb R}
\newcommand{\Q}{\mathbb Q}
\newcommand{\Z}{\mathbb Z}
\newcommand{\A}{\mathbb A}
\newcommand{\Spec}{\operatorname{Spec}}
\newcommand{\Nm}{\operatorname{Nm}}
\newcommand{\diag}{\operatorname{diag}}
\newcommand{\dd}{\mathrm d}
\newcommand{\Open}{\textbf{OPEN.}\ }

\title{The Koide Quadratic Cone as a Lorentzian Determinant\\
\large Vieta dynamics and a ternary defect-one counting problem}
\author{}
\date{Specialist-verification draft, 4 September 2026}

\begin{document}
\maketitle

\begin{abstract}
Writing the charged-lepton Koide relation in square-root mass coordinates produces a quadratic cone of signature $(1,2)$.  We give an explicit orthogonal identification of this cone with the determinant-zero locus in $\operatorname{Sym}_2(\R)$.  In the exact period convention $\Pi=[Z\mid I_2]$, $Z=\left(\begin{smallmatrix}6\mu&\tau\\ \beta&\mu\end{smallmatrix}\right)$, used in a recent proposed complex-torus construction on $S^6$, the same map identifies the Koide cone with the zero-discriminant locus in the ambient imaginary-coefficient space.  This coefficient-space calculation neither assumes nor establishes the claimed global construction.  The arithmetic part is independent of it.  For
\[
 \Delta(x,y,z)=4(xy+xz+yz)-(x^2+y^2+z^2),
\]
we study ordered positive solutions of $\Delta=-1$ reduced by the Vieta wall $z\geq4(x+y)$.  An exact integral change of variables identifies them with a chamber on
\[
 a^2-3q^2-18n^2=-2,
 \qquad z=4a+17n.
\]
The three Vieta substitutions are integral reflections of the Lorentzian
form $\Delta$; on unordered positive solutions they form a forest in which
every component has one reduced root.  The two boundary faces of the ternary
chamber are completely described by negative-Pell recurrences and converge
to the two null directions bounding the archimedean sector.
The local densities, the archimedean chamber factor, the Gorodnik--Oh--Shah
sector hypotheses, and the normalization against Borovoi's Tamagawa measure
are all made explicit.  A $2$-adic sign-reversing, measure-preserving
bijection proves that Borovoi's adelic coefficient is one, yielding the
unconditional explicit main term.  An exact HNF parametrization by ideals of
$\mathbb Z[\sqrt3]$, unique unit-orbit descent, and a sharp floor formula
identify every counted point with multiplicity one.  Direct verification of
Benoist--Oh effective well-roundedness then gives
$R(X)=C_0X+O(X^{1-\eta})$ for some unspecified $\eta>0$.  The same shell is
identified exactly with a finite union of primitive hyperbolic conjugacy
classes in a cocompact quaternionic lattice, but its standard radial height
is $7a+15\sqrt2\,n$, not the sharp height $4a+17n$.  Exact enumeration through
height $10^8$ gives $R(10^8)=3{,}536{,}002$.  The primary-source audit finds
no theorem proving the pointwise square-root discrepancy for the exact
five-wall cutoff; that stronger estimate remains open.
\end{abstract}

\paragraph{Status.}
This is a mathematical draft for independent checking.  It makes no claim that the Koide relation is a law of particle physics, no claim about the existence of the proposed complex structure on $S^6$, and no Yang--Mills claim.  The exact algebraic identities, parametrization, theorem audit, and bounded census are stated with their verification data in the canonical source modules included below.

\section{The Koide form and its Lorentzian model}

For a triple of nonnegative masses $m_1,m_2,m_3$, put $r_i=\sqrt{m_i}$.  The Koide relation may be written
\begin{equation}\label{eq:koide-relation}
 m_1+m_2+m_3=\frac23\left(\sqrt{m_1}+\sqrt{m_2}+\sqrt{m_3}\right)^2,
\end{equation}
which was introduced in the charged-lepton context by \citet{Koide_1983}.  The circle interpretation of \citet{Kocik_2012} supplies a separate Euclidean realization.  Here the relevant object is the homogeneous quadratic form
\begin{equation}\label{eq:QK}
 Q_K(r_1,r_2,r_3)
 :=\frac23(r_1+r_2+r_3)^2-(r_1^2+r_2^2+r_3^2).
\end{equation}
Thus \eqref{eq:koide-relation} is $Q_K(r)=0$.

Define the orthogonal Fourier coordinates
\begin{equation}\label{eq:Fourier}
 \mathcal F(r_1,r_2,r_3)
 =\left(
 \frac{r_1+r_2+r_3}{\sqrt3},
 \frac{r_1-r_3}{\sqrt2},
 \frac{r_1+r_3-2r_2}{\sqrt6}
 \right)=:(h,x,y),
\end{equation}
and the symmetric-matrix embedding
\begin{equation}\label{eq:iota}
 \iota(h,x,y)=
 \begin{pmatrix}h+y&x\\x&h-y\end{pmatrix}.
\end{equation}

\begin{theorem}[Orthogonal determinant model]\label{thm:orthogonal-model}
The map $\mathcal F$ is orthogonal and
\begin{equation}\label{eq:lorentz-identity}
 Q_K(r)=h^2-x^2-y^2=\det\iota(h,x,y).
\end{equation}
Its inverse is
\begin{align}
 r_1&=\frac h{\sqrt3}+\frac x{\sqrt2}+\frac y{\sqrt6},\label{eq:F-inv1}\\
 r_2&=\frac h{\sqrt3}-\frac{2y}{\sqrt6},\label{eq:F-inv2}\\
 r_3&=\frac h{\sqrt3}-\frac x{\sqrt2}+\frac y{\sqrt6}.\label{eq:F-inv3}
\end{align}
For $r\in\R_{\geq0}^3\setminus\{0\}$, the Koide equation is equivalent to the assertion that $\iota(\mathcal F(r))$ is positive semidefinite of rank one.  Equivalently, the Euclidean angle between $r$ and $(1,1,1)/\sqrt3$ is $\pi/4$.
\end{theorem}

\begin{proof}
The three rows of the matrix in \eqref{eq:Fourier} form an orthonormal basis.  Hence
\[
 r_1^2+r_2^2+r_3^2=h^2+x^2+y^2,
 \qquad (r_1+r_2+r_3)^2=3h^2,
\]
and substitution in \eqref{eq:QK} gives \eqref{eq:lorentz-identity}.  Transposition of the orthogonal matrix gives \eqref{eq:F-inv1}--\eqref{eq:F-inv3}.  If $r\in\R_{\geq0}^3\setminus\{0\}$, then $h>0$.  On the null cone, $h^2=x^2+y^2$, so $\iota$ has trace $2h>0$, determinant zero, and eigenvalues $2h,0$.  Finally,
\[
 \cos\angle\left(r,\frac{(1,1,1)}{\sqrt3}\right)
 =\frac{h}{\sqrt{h^2+x^2+y^2}}=\frac1{\sqrt2}.
\]
All implications reverse.
\end{proof}

\begin{corollary}[Exact phase sector]
\label{cor:exact-phase-sector}
Let $r\in\R^3$ satisfy $Q_K(r)=0$, and suppose that the first Fourier
coordinate in \eqref{eq:Fourier} satisfies $h>0$.  There is a unique
$\theta\in\R/(2\pi\Z)$ such that
\begin{align}
 r_1&=\frac h{\sqrt3}
 \left(1+\sqrt2\cos\left(\theta-\frac\pi6\right)\right),
 \label{eq:phase-r1}\\
 r_2&=\frac h{\sqrt3}
 \left(1+\sqrt2\cos\left(\theta+\frac\pi2\right)\right),
 \label{eq:phase-r2}\\
 r_3&=\frac h{\sqrt3}
 \left(1+\sqrt2\cos\left(\theta-\frac{5\pi}6\right)\right).
 \label{eq:phase-r3}
\end{align}
Moreover, $r_1,r_2,r_3\geq0$ if and only if
\begin{equation}
 \theta\in
 \left[\frac\pi{12},\frac\pi4\right]
 \cup\left[\frac{3\pi}4,\frac{11\pi}{12}\right]
 \cup\left[\frac{17\pi}{12},\frac{19\pi}{12}\right]
 \pmod {2\pi}.
 \label{eq:nonnegative-phase-arcs}
\end{equation}
At each endpoint exactly one coordinate vanishes, and the ratio of the two
nonzero coordinates is $2+\sqrt3$ or $(2+\sqrt3)^{-1}=2-\sqrt3$.
\end{corollary}

\begin{proof}
The null equation and $h>0$ give
$x=h\cos\theta$, $y=h\sin\theta$ for a unique
$\theta\in\R/(2\pi\Z)$.  Substitution in
\eqref{eq:F-inv1}--\eqref{eq:F-inv3}, followed by the three elementary
angle-addition identities, gives
\eqref{eq:phase-r1}--\eqref{eq:phase-r3}.

The inequality $1+\sqrt2\cos\phi\geq0$ fails precisely for
$3\pi/4<\phi<5\pi/4$ modulo $2\pi$.  Applied to the three phases in
\eqref{eq:phase-r1}--\eqref{eq:phase-r3}, the three excluded open arcs are
\[
 \left(\frac{11\pi}{12},\frac{17\pi}{12}\right),\qquad
 \left(\frac\pi4,\frac{3\pi}4\right),\qquad
 \left(\frac{19\pi}{12},2\pi\right)
 \cup\left(0,\frac\pi{12}\right).
\]
Their complement is exactly \eqref{eq:nonnegative-phase-arcs}.  If, for
example, $r_3=0$, the unchanged equation $Q_K(r)=0$ becomes
\[
 r_1^2-4r_1r_2+r_2^2=0,
\]
so $r_1/r_2=2\pm\sqrt3$.  Cyclic symmetry proves the assertion at every
endpoint.
\end{proof}

\subsection{The canonical period discriminant as an ambient model}

The recent online manuscript \citet{AlpogeS6_2026} uses period functions
$\tau,\mu,\beta$ and the exact convention
\begin{equation}\label{eq:canonical-complex-period}
 Z=\begin{pmatrix}6\mu&\tau\\ \beta&\mu\end{pmatrix},
 \qquad \Pi=[Z\mid I_2].
\end{equation}
Write
\[
 t=\operatorname{Im}\tau,\qquad
 m=\operatorname{Im}\mu,\qquad
 \varrho=\operatorname{Im}\beta.
\]
Relative to the period-lattice basis
$(\widehat\gamma,\widehat u,\widehat w,\widehat\delta)$ and the interleaved
real codomain order
$(\operatorname{Re}z_1,\operatorname{Im}z_1,
  \operatorname{Re}z_2,\operatorname{Im}z_2)$,
the realification of \eqref{eq:canonical-complex-period} is
\begin{equation}\label{eq:canonical-real-period}
 P(\tau,\mu,\beta)=
 \begin{pmatrix}
  6\operatorname{Re}\mu&\operatorname{Re}\tau&1&0\\
  6m&t&0&0\\
  \operatorname{Re}\beta&\operatorname{Re}\mu&0&1\\
  \varrho&m&0&0
 \end{pmatrix}.
\end{equation}
For comparison with the Lorentz cone, define the auxiliary symmetric
coefficient matrix
\begin{equation}\label{eq:symmetric-coefficient-matrix}
 S(t,m,\varrho)
 :=\begin{pmatrix}t&\sqrt6\,m\\ \sqrt6\,m&\varrho\end{pmatrix}.
\end{equation}
The matrix $S$ represents the same quadratic discriminant, but it is not the
canonical $4\times4$ real period matrix $P$.

\begin{proposition}[Exact canonical coefficient dictionary]
\label{prop:S6-dictionary}
The real-linear map
\begin{equation}\label{eq:coefficient-map}
 \mathcal K(h,x,y)
 =\left(h+y,\frac{x}{\sqrt6},h-y\right)=(t,m,\varrho)
\end{equation}
is an isomorphism from $\R^3$ to the ambient imaginary-coefficient space,
with inverse
\begin{equation}\label{eq:coefficient-map-inverse}
 \mathcal K^{-1}(t,m,\varrho)
 =\left(\frac{t+\varrho}{2},\sqrt6\,m,
        \frac{t-\varrho}{2}\right).
\end{equation}
For every choice of the real parts of $\tau,\mu,\beta$,
\begin{equation}\label{eq:canonical-discriminant}
 \det P(\tau,\mu,\beta)
 =\det S(t,m,\varrho)
 =t\varrho-6m^2
 =-\det\operatorname{Im}Z.
\end{equation}
Consequently, if $(t,m,\varrho)=\mathcal K(\mathcal F(r))$, then
\begin{equation}\label{eq:period-koide-identity}
 \det P(\tau,\mu,\beta)
 =\det S\bigl(\mathcal K(\mathcal F(r))\bigr)
 =Q_K(r).
\end{equation}
For $r\in\R_{\geq0}^3\setminus\{0\}$ on the Koide cone, the symmetric
matrix $S\bigl(\mathcal K(\mathcal F(r))\bigr)=\iota(\mathcal F(r))$ is
positive semidefinite of rank one.
\end{proposition}

\begin{proof}
Adding and subtracting the first and third entries of
\eqref{eq:coefficient-map}, and multiplying the second by $\sqrt6$, gives
\eqref{eq:coefficient-map-inverse}.  Expanding
\eqref{eq:canonical-real-period} along its last two columns gives
\[
 \det P
 =-\det\begin{pmatrix}6m&t\\ \varrho&m\end{pmatrix}
 =t\varrho-6m^2.
\]
The same expression is the determinant of
\eqref{eq:symmetric-coefficient-matrix}, and substitution from
\eqref{eq:coefficient-map} gives
\[
 (h+y)(h-y)-6\left(\frac{x}{\sqrt6}\right)^2
 =h^2-x^2-y^2=Q_K(r).
\]
The final assertion is Theorem~\ref{thm:orthogonal-model} with
$S\bigl(\mathcal K(h,x,y)\bigr)=\iota(h,x,y)$.
\end{proof}

\begin{corollary}[Rank-one square parametrization]
\label{cor:rank-one-square}
The map $\mathcal K$ identifies the future Lorentz null cone
$\{h^2=x^2+y^2,\ h>0\}$ with the nonzero positive-semidefinite rank-one
locus in $\operatorname{Sym}_2(\R)$.  Equivalently, every point of that
coefficient locus has the form
\begin{equation}
 (t,m,\varrho)=
 \left(u^2,\frac{uv}{\sqrt6},v^2\right),
 \qquad (u,v)\in\R^2\setminus\{(0,0)\},
 \label{eq:rank-one-square-coefficients}
\end{equation}
and $(u,v)$ is unique up to simultaneous sign.  Under the inverse
coefficient map,
\begin{equation}
 (h,x,y)=
 \left(\frac{u^2+v^2}{2},uv,\frac{u^2-v^2}{2}\right).
 \label{eq:rank-one-square-lorentz}
\end{equation}
The sublocus arising from nonnegative $r_i$ is exactly the phase sector in
Corollary~\ref{cor:exact-phase-sector}; positive semidefiniteness alone does
not impose the three mass-coordinate inequalities.
\end{corollary}

\begin{proof}
A nonzero real symmetric matrix is positive semidefinite of rank one if and
only if it is $ww^{\mathsf T}$ for a nonzero column $w=(u,v)^{\mathsf T}$;
the factor $w$ is unique up to $w\mapsto-w$.  Comparing
$ww^{\mathsf T}$ with \eqref{eq:symmetric-coefficient-matrix} gives
\eqref{eq:rank-one-square-coefficients}, and
\eqref{eq:coefficient-map-inverse} gives
\eqref{eq:rank-one-square-lorentz}.  Conversely the latter triple has
$h>0$ and $h^2-x^2-y^2=0$.  The final statement follows from
Corollary~\ref{cor:exact-phase-sector}.
\end{proof}

\begin{corollary}[The Koide cone is not a family of regular tori]
\label{cor:ambient-only}
The identification in Proposition~\ref{prop:S6-dictionary} is only a
real-linear identification of ambient coefficient spaces and their quadratic
discriminants.  It supplies neither the real parts of the periods nor the
holomorphic and equivariance conditions of the proposed period family.  At
every point of the Koide cone, every matrix of the canonical form
\eqref{eq:canonical-real-period} with coefficient triple
\eqref{eq:coefficient-map} has determinant zero.  Its columns therefore do
not span $\R^4$, do not generate a full lattice, and do not define a regular
compact complex two-torus.
\end{corollary}

\begin{remark}[Separation from the positive radial cusp]
The audited period calculation associated with \citet{AlpogeS6_2026}
records holomorphic cusp germs $\mu,h_{\mathrm{per}},b$ at $t_c=0$, with
$s=iT$, $\tau=h_{\mathrm{per}}+s$, and $b=\beta+\tau$ on the positive
radial specialization $t_c=e^{-2\pi T}$.  Write
\[
 \mu_T=\alpha_T+i m_T,\qquad
 h_{\mathrm{per},T}=\eta_T+i H_T,\qquad
 b_T=\xi_T+i B_T,
\]
and set
\[
 c_T=\xi_T-\eta_T,\qquad
 L_T=T+H_T,\qquad q_T=L_T-B_T,\qquad
 D_T=L_Tq_T+6m_T^2.
\]
Then the imaginary-coefficient triple is exactly
\begin{equation}\label{eq:radial-coefficient-triple}
 (\operatorname{Im}\tau_T,\operatorname{Im}\mu_T,
   \operatorname{Im}\beta_T)=(L_T,m_T,-q_T).
\end{equation}
Its exact realification is
\begin{equation}\label{eq:radial-real-period}
 P_T=
 \begin{pmatrix}
  6\alpha_T&\eta_T&1&0\\
  6m_T&L_T&0&0\\
  c_T&\alpha_T&0&1\\
  -q_T&m_T&0&0
 \end{pmatrix},
 \qquad \det P_T=-D_T.
\end{equation}
Indeed, \eqref{eq:canonical-discriminant} and
\eqref{eq:radial-coefficient-triple} give
\[
 \det P_T=L_T(-q_T)-6m_T^2=-D_T.
\]
If
$m=\operatorname{Im}\mu(0)$, $H=\operatorname{Im}h_{\mathrm{per}}(0)$, and
$B=\operatorname{Im}b(0)$, then
\[
 m_T=m+O(e^{-2\pi T}),\quad
 H_T=H+O(e^{-2\pi T}),\quad
 B_T=B+O(e^{-2\pi T}).
\]
Direct multiplication, without suppressing the constant term, gives
\begin{equation}\label{eq:radial-determinant-asymptotic}
 D_T=T^2+(2H-B)T+H(H-B)+6m^2+O(Te^{-2\pi T}),
 \qquad \det P_T\sim-T^2.
\end{equation}
In particular $D_T>0$ for all sufficiently large $T$, so these radial
matrices are invertible and lie strictly off the Koide discriminant
hypersurface $\{\det P=0\}$; their period-lattice covolume is then $D_T$.
More precisely, their auxiliary symmetric matrix is
\[
 S_T=\begin{pmatrix}L_T&\sqrt6\,m_T\\
                     \sqrt6\,m_T&-q_T\end{pmatrix},
 \qquad \det S_T=-D_T<0,
\]
and therefore has signature $(1,1)$ for all sufficiently large $T$.  This
is different not merely from the zero determinant on the Koide locus but
also from its positive-semidefinite rank-one sign chamber.
This exact separation uses the recorded cusp formulas only; it is not an
independent proof of the proposed global $S^6$ construction.
\end{remark}

% Workbench-ready, self-contained mathematical section.
% Required bibliography keys: Borovoi_2001, GorodnikOhShah_2009,
% Voight_2021, BenoistOh_2012, Good_1983, and ChatzakosPetridis_2016.
% The exact parametrization, census, quaternion bridge, and primary-source
% theorem audit are included as namespaced modules at the end of this file.

\section{Exact arithmetic of the Koide defect-one shell}
\label{koide-count-sec-main}

This section isolates an arithmetic result that does not use the existence of
the proposed complex structure on \(S^6\), any period-map assertion, or any
Yang--Mills interpretation.  All coordinates, chambers, local factors, and
normalizations used in the final constant are fixed explicitly below.

For \(r=(r_1,r_2,r_3)\), define the Koide quadratic form
\begin{equation}
 Q_K(r_1,r_2,r_3)
 :=\frac23(r_1+r_2+r_3)^2-(r_1^2+r_2^2+r_3^2).
 \label{koide-count-eq-qk}
\end{equation}
For integral coordinates put
\begin{equation}
 \Delta(x,y,z):=4(xy+xz+yz)-(x^2+y^2+z^2).
 \label{koide-count-eq-defect}
\end{equation}
Expanding \eqref{koide-count-eq-qk} gives the exact identity
\begin{equation}
 \Delta(x,y,z)=3Q_K(x,y,z).
 \label{koide-count-eq-defect-qk}
\end{equation}
Thus the equation studied here,
\(\Delta=-1\), is the nonzero level \(Q_K=-1/3\); it is not an affine
translate of the null cone \(Q_K=0\).

\subsection{Vieta reduction and exact ternary coordinates}
\label{koide-count-subsec-vieta}

Regarded as a monic quadratic equation in \(z\), the defect equation is
\begin{equation}
 z^2-4(x+y)z+x^2+y^2-4xy-1=0.
 \label{koide-count-eq-monic-z}
\end{equation}
Consequently the other root is
\begin{equation}
 z'=4(x+y)-z.
 \label{koide-count-eq-vieta-root}
\end{equation}
The analogous substitutions in \(x\) and \(y\) are integral involutions
preserving \(\Delta\).

Let $e_1,e_2,e_3$ be the standard basis and let $B_K$ be the integral
bilinear form with Gram matrix
\begin{equation}
 A_K=
 \begin{pmatrix}
 -1&2&2\\
 2&-1&2\\
 2&2&-1
 \end{pmatrix},
 \qquad
 B_K(u,v)=u^{\mathsf T}A_Kv.
 \label{koide-count-eq-vieta-gram}
\end{equation}
Thus $\Delta(u)=B_K(u,u)$.  The eigenvalue of $A_K$ on
$\mathbb R(1,1,1)$ is $3$, and its eigenvalue on the orthogonal complement is
$-3$, so this integral presentation has signature $(1,2)$.

\begin{proposition}[Vieta moves are Lorentzian reflections]
\label{koide-count-prop-vieta-reflections}
For $i=1,2,3$, define
\begin{equation}
 s_i(u)=u-
 \frac{2B_K(u,e_i)}{B_K(e_i,e_i)}e_i.
 \label{koide-count-eq-reflection-definition}
\end{equation}
Then $s_i\in\operatorname O(B_K,\mathbb Z)$,
$s_i^2=1$, and $\det s_i=-1$.  Explicitly,
\begin{align}
 s_1(x,y,z)&=(-x+4y+4z,y,z),\label{koide-count-eq-vieta-s1}\\
 s_2(x,y,z)&=(x,-y+4x+4z,z),\label{koide-count-eq-vieta-s2}\\
 s_3(x,y,z)&=(x,y,-z+4x+4y).\label{koide-count-eq-vieta-s3}
\end{align}
Consequently these are exactly the three Vieta root substitutions for
$\Delta=-1$, now identified as reflections of the unchanged Lorentzian
lattice $(\mathbb Z^3,B_K)$.
\end{proposition}

\begin{proof}
Since $B_K(e_i,e_i)=-1$, formula
\eqref{koide-count-eq-reflection-definition} is integral and direct
evaluation gives \eqref{koide-count-eq-vieta-s1}--
\eqref{koide-count-eq-vieta-s3}.  The standard reflection calculation
gives
\[
 B_K(s_i(u),s_i(v))=B_K(u,v),\qquad s_i^2=1.
\]
The map fixes the hyperplane $e_i^\perp$ pointwise and sends $e_i$ to
$-e_i$, hence has determinant $-1$.  Finally, when the equation is viewed
as a monic quadratic in one coordinate, the displayed replacement is the
sum of its two roots minus the original root, proving the last assertion.
\end{proof}

\begin{proposition}[Strict Vieta reduction]
\label{koide-count-prop-vieta-reduction}
Let \(0<x\leq y\leq z\) be integers satisfying
\(\Delta(x,y,z)=-1\).  If \(z<4(x+y)\), then the other root in
\eqref{koide-count-eq-vieta-root} satisfies \(0<z'<y\).  If
\(z\geq4(x+y)\), then \(z'\leq0\).
\end{proposition}

\begin{proof}
In the first case \(z'>0\) follows immediately from
\eqref{koide-count-eq-vieta-root}.  Let \(P(Z)\) be the polynomial on the
left of \eqref{koide-count-eq-monic-z}.  Since \(x\leq y\),
\[
 P(y)=x^2-8xy-2y^2-1
 \leq xy-8xy-2y^2-1<0.
\]
But \(P(y)=(y-z)(y-z')\), and \(y-z\leq0\); hence \(y-z'>0\), so
\(z'<y\).  The second assertion again follows directly from
\eqref{koide-count-eq-vieta-root}.
\end{proof}

Define the reduced-root set and its height count by
\begin{align}
 \mathcal R
 &:={}
 \{(x,y,z)\in\mathbb Z_{>0}^3:
 x\leq y\leq z,\ \Delta(x,y,z)=-1,\ z\geq4(x+y)\},
 \label{koide-count-eq-reduced-set}\\
 R(X)&:=\#\{(x,y,z)\in\mathcal R:z\leq X\}.
 \label{koide-count-eq-rx-definition}
\end{align}
Proposition~\ref{koide-count-prop-vieta-reduction} shows that the wall
\(z\geq4(x+y)\) is exactly the terminal side of the decreasing Vieta move
after the coordinates have been ordered.

\begin{proposition}[The positive Vieta graph is a rooted forest]
\label{koide-count-prop-vieta-forest}
Let $\mathcal S^+$ be the set of unordered triples of positive integers
satisfying $\Delta=-1$.  Join two vertices in a simple graph when one is
obtained from the other by a Vieta reflection $s_i$ and the reflected
coordinate is positive.  Every connected component of this graph is a
tree with exactly one root, and the roots are precisely the unordered
triples whose ordered representatives lie in $\mathcal R$.  Equivalently,
every positive integral solution has a unique reduced representative, and
it is obtained by repeatedly ordering the coordinates and reflecting the
largest coordinate whenever it lies below the terminal wall.
\end{proposition}

\begin{proof}
First, no Vieta reflection fixes a point of $\Delta=-1$.  Indeed, a fixed
reflection in a coordinate $w$, with the other two coordinates denoted by
$u,v$, would give $w=2(u+v)$, and substitution in the unchanged form gives
\[
 \Delta(u,v,2(u+v))=3(u^2+4uv+v^2)>0.
\]
Thus the sum of the three coordinates changes strictly across every edge.

Now order a vertex as $0<x\leq y\leq z$.  In fact $y<z$, because $y=z$
would give
\[
 \Delta(x,y,y)=-x^2+8xy+2y^2>0.
\]
Reflecting either of the two smaller coordinates produces a new coordinate
strictly larger than $z$:
\[
 4(y+z)-x>z,\qquad 4(x+z)-y>z.
\]
If $z<4(x+y)$, Proposition~\ref{koide-count-prop-vieta-reduction} says
that reflection in the largest coordinate produces $0<z'<y$ and strictly
decreases the coordinate sum.  If $z\geq4(x+y)$, that reflection is
nonpositive and hence is not an edge of $\mathcal S^+$.  Therefore every
nonterminal vertex has exactly one neighbour of smaller coordinate sum,
whereas a terminal vertex has none.

Iteration of the unique decreasing move terminates because the coordinate
sum is a positive integer.  The endpoint is represented by an element of
$\mathcal R$.  The graph has no cycle: on a putative simple cycle choose a
vertex of largest coordinate sum; its two neighbours on the cycle would be
two distinct smaller neighbours, contrary to uniqueness.  Thus each
component is a tree.  A tree in which following the unique decreasing edge
always terminates has exactly one terminal vertex, proving both uniqueness
and the characterization of the roots.
\end{proof}

For an ordered triple define
\begin{equation}
 n=z-4(x+y),\qquad q=y-x,\qquad a=x+y-4n,
 \label{koide-count-eq-forward-coordinates}
\end{equation}
and let
\begin{equation}
 \mathcal T:=
 \left\{(a,q,n)\in\mathbb Z^3:
 \begin{array}{c}
 a>0,\ q\geq0,\ n\geq0,\ a+4n>q,\\
 a^2-3q^2-18n^2=-2
 \end{array}
 \right\}.
 \label{koide-count-eq-ternary-chamber}
\end{equation}

\begin{theorem}[Integral ternary-coordinate bijection]
\label{koide-count-thm-coordinate-bijection}
The map in \eqref{koide-count-eq-forward-coordinates} restricts to a
bijection \(\Phi:\mathcal R\to\mathcal T\).  Its inverse is
\begin{equation}
 (a,q,n)\longmapsto
 \left(
 \frac{a+4n-q}{2},
 \frac{a+4n+q}{2},
 4a+17n
 \right).
 \label{koide-count-eq-inverse-coordinates}
\end{equation}
For arbitrary integral \((x,y,z)\) and the associated \((a,q,n)\),
\begin{equation}
 2\Delta(x,y,z)=a^2-3q^2-18n^2,
 \qquad z=4a+17n.
 \label{koide-count-eq-polynomial-identity}
\end{equation}
In particular,
\begin{equation}
 R(X)=\#\{(a,q,n)\in\mathcal T:4a+17n\leq X\}.
 \label{koide-count-eq-exact-count}
\end{equation}
\end{theorem}

\begin{proof}
Put \(p=x+y\).  Since \(q=y-x\), one has
\(2x=p-q\) and \(2y=p+q\), and direct substitution into
\eqref{koide-count-eq-defect} gives
\[
 2\Delta=p^2-3q^2+8pz-2z^2.
\]
The definition \(n=z-4p\) gives \(z=4p+n\), whence
\[
 2\Delta=p^2-3q^2-8pn-2n^2.
\]
Finally \(a=p-4n\), equivalently \(p=a+4n\), gives both identities in
\eqref{koide-count-eq-polynomial-identity}.

Suppose first that \((x,y,z)\in\mathcal R\).  Then \(q,n\geq0\), while
\(a+4n=p>q\) because \(p-q=2x>0\).  The defect equation and
\eqref{koide-count-eq-polynomial-identity} give
\(a^2-3q^2-18n^2=-2\).  If \(n=0\), then \(a=p>0\).  If \(n\geq1\),
then
\[
 a^2=3q^2+18n^2-2\geq18n^2-2\geq16n^2.
\]
Were \(a\leq0\), this inequality would force \(a\leq-4n\), contrary to
\(a+4n>q\geq0\).  Hence \(a>0\), and the forward map lands in
\(\mathcal T\).

Conversely, reducing the ternary equation modulo \(2\) gives
\(a\equiv q\pmod2\), so all three coordinates in
\eqref{koide-count-eq-inverse-coordinates} are integral.  Put
\(p=a+4n\).  The chamber inequality gives
\[
 x=\frac{p-q}{2}>0,\qquad
 y=\frac{p+q}{2}\geq x.
\]
Moreover
\[
 z=4a+17n=4p+n\geq4p>p>y.
\]
The identities \eqref{koide-count-eq-polynomial-identity} now give
\(\Delta=-1\) and \(z-4(x+y)=n\geq0\).  Thus the inverse map lands in
\(\mathcal R\).  The displayed formulae show directly that the two
compositions are the identity.  The last assertion follows from the height
identity in \eqref{koide-count-eq-polynomial-identity}.
\end{proof}

\begin{proposition}[Complete Pell classification of the chamber faces]
\label{koide-count-prop-pell-faces}
The integral points of $\mathcal T$ on its two coordinate faces are the
following two disjoint families.

On the terminal Vieta wall $n=0$, define integers $a_k,q_k$ by
\begin{equation}
 a_k+q_k\sqrt3=(1+\sqrt3)(2+\sqrt3)^k
 \qquad(k\geq0).
 \label{koide-count-eq-wall-pell-sequence}
\end{equation}
Then the points of $\mathcal T$ with $n=0$ are exactly
$(a_k,q_k,0)$ for $k\geq1$, and their reduced roots are
\begin{equation}
 \left(
  \frac{a_k-q_k}{2},
  \frac{a_k+q_k}{2},
  4a_k
 \right).
 \label{koide-count-eq-wall-pell-roots}
\end{equation}
The first is $(1,4,20)$.

On the symmetry face $q=0$, define integers $A_k,N_k$ by
\begin{equation}
 A_k+N_k\sqrt{18}=(4+\sqrt{18})(17+4\sqrt{18})^k
 \qquad(k\geq0).
 \label{koide-count-eq-symmetric-pell-sequence}
\end{equation}
Then the points of $\mathcal T$ with $q=0$ are exactly
$(A_k,0,N_k)$ for $k\geq0$, and their reduced roots are
\begin{equation}
 \left(
  \frac{A_k+4N_k}{2},
  \frac{A_k+4N_k}{2},
  4A_k+17N_k
 \right).
 \label{koide-count-eq-symmetric-pell-roots}
\end{equation}
The first is $(4,4,33)$.

In the projectivized real quadric these families tend, respectively, to
the two chamber-boundary rays
\begin{equation}
 (a,q,n)=\mathbb R_{>0}(\sqrt3,1,0),
 \qquad
 (a,q,n)=\mathbb R_{>0}(3\sqrt2,0,1).
 \label{koide-count-eq-pell-limit-rays}
\end{equation}
Each family contributes only $O(\log X)$ points to $R(X)$.
\end{proposition}

\begin{proof}
On $n=0$, the equation of $\mathcal T$ is
$a^2-3q^2=-2$.  Here $q=0$ is impossible.  If $q=1$, positivity gives
$a=1$, which fails the strict chamber inequality $a>q$.  For any positive
solution with $q\geq2$, multiplication by the norm-one unit
$2-\sqrt3$ gives
\begin{equation}
 a'+q'\sqrt3=(a+q\sqrt3)(2-\sqrt3),\qquad
 a'=2a-3q,\quad q'=2q-a.
 \label{koide-count-eq-wall-pell-descent}
\end{equation}
The unchanged equation implies
\[
 a<2q,\qquad
 a>\frac32q,\qquad
 a>q
 \quad(q\geq2),
\]
because its squares differ from $4q^2$, $9q^2/4$, and $q^2$ by,
respectively, $q^2+2$, $3q^2/4-2>0$, and $2q^2-2>0$.
Thus $a',q'>0$ and $q'<q$.  Repeated descent reaches $(1,1)$, and reversing
the descent proves \eqref{koide-count-eq-wall-pell-sequence} is the complete
positive solution set.  The recurrence
\[
 a_{k+1}=2a_k+3q_k,\qquad q_{k+1}=a_k+2q_k
\]
shows inductively that $a_k,q_k$ are both odd.  The chamber excludes only
$k=0$, and Theorem~\ref{koide-count-thm-coordinate-bijection} now gives
\eqref{koide-count-eq-wall-pell-roots}.

On $q=0$, one has $a^2-18n^2=-2$.  Necessarily $n>0$.  Direct substitution
for $1\leq n\leq5$ gives the sole solution $(a,n)=(4,1)$.  If $n\geq6$,
multiplication by the norm-one unit $17-4\sqrt{18}$ gives
\begin{equation}
 a'+n'\sqrt{18}=(a+n\sqrt{18})(17-4\sqrt{18}),\qquad
 a'=17a-72n,\quad n'=17n-4a.
 \label{koide-count-eq-symmetric-pell-descent}
\end{equation}
Both new coordinates are positive.  Indeed,
\[
 \frac an=\sqrt{18-\frac2{n^2}}
 >\frac{72}{17}\quad(n\geq6),
 \qquad
 \frac an<\sqrt{18}<\frac{17}{4}.
\]
Moreover $a>4n$ and hence $0<n'<n$.  Descent therefore reaches $(4,1)$,
and reversal proves \eqref{koide-count-eq-symmetric-pell-sequence} is the
complete positive solution set.  Its recurrence,
\[
 A_{k+1}=17A_k+72N_k,\qquad
 N_{k+1}=4A_k+17N_k,
\]
keeps $A_k$ even and $N_k$ odd, so
\eqref{koide-count-eq-symmetric-pell-roots} follows from the inverse
coordinate formula.

Conjugating the two identities defining the sequences shows
\[
 \frac{a_k}{q_k}\longrightarrow\sqrt3,\qquad
 \frac{A_k}{N_k}\longrightarrow\sqrt{18}=3\sqrt2,
\]
which proves \eqref{koide-count-eq-pell-limit-rays}.  Finally, the dominant
factors $2+\sqrt3>1$ and $17+4\sqrt{18}>1$ give exponential growth in $k$,
whereas the heights in \eqref{koide-count-eq-wall-pell-roots} and
\eqref{koide-count-eq-symmetric-pell-roots} are positive linear forms in
the displayed coordinates.  Hence either face contains only $O(\log X)$
points of height at most $X$.
\end{proof}

\subsection{The affine quadric and its local densities}
\label{koide-count-subsec-local}

Put
\begin{equation}
 F(a,q,n):=a^2-3q^2-18n^2,
 \qquad
 \mathscr X:=\{F=-2\}\subset\mathbb A^3.
 \label{koide-count-eq-affine-quadric}
\end{equation}
The matrix of \(F\) is
\(A=\operatorname{diag}(1,-3,-18)\), so
\(\det A=54\), and
\(-(-2)\det A=108\) is not a square.  The integral point
\(x_0=(1,1,0)\) lies on \(\mathscr X\).

For every prime \(p\), define
\begin{equation}
 N_{p,k}:=
 \#\{(a,q,n)\bmod p^k:F(a,q,n)\equiv-2\pmod{p^k}\},
 \qquad
 \sigma_p:=\lim_{k\to\infty}p^{-2k}N_{p,k}.
 \label{koide-count-eq-local-density-definition}
\end{equation}

\begin{theorem}[Exact local densities]
\label{koide-count-thm-local-densities}
The limits in \eqref{koide-count-eq-local-density-definition} exist and
are
\begin{equation}
 \sigma_2=\frac32,
 \qquad
 \sigma_3=2,
 \qquad
 \sigma_p=
 \frac{1-p^{-2}}{1-\left(\frac3p\right)p^{-1}}
 \quad(p\geq5).
 \label{koide-count-eq-all-local-densities}
\end{equation}
Writing \(\chi_{12}\) for the primitive real character modulo \(12\), the
singular series, with the primes taken in increasing order (equivalently,
with Borovoi's convergence-factor prescription), is therefore
\begin{equation}
 \mathfrak S(\mathscr X):=\prod_p\sigma_p
 =\frac92\frac{L(1,\chi_{12})}{\zeta(2)}
 =\frac92\frac{\log(2+\sqrt3)}{\sqrt3\,\zeta(2)}.
 \label{koide-count-eq-singular-series}
\end{equation}
\end{theorem}

\begin{proof}
We compute the two bad primes first.

For \(p=3\), once \(q,n\bmod3^k\) have been chosen, the equation for
\(a\) is
\[
 a^2\equiv-2+3q^2+18n^2\equiv1\pmod3.
\]
The two roots modulo \(3\) are nonsingular, and each has a unique Hensel
lift modulo every \(3^k\).  Thus
\begin{equation}
 N_{3,k}=2\cdot3^{2k},
 \label{koide-count-eq-three-adic-count}
\end{equation}
which gives \(\sigma_3=2\).

For \(p=2\), reduction modulo \(4\) leaves exactly two parity branches:
\begin{equation}
 (a,q,n)\equiv(0,0,1)\pmod2,
 \qquad
 (a,q,n)\equiv(1,1,0)\pmod2.
 \label{koide-count-eq-two-adic-branches}
\end{equation}
Fix \(k\geq4\) and put \(j=k-1\).  On the first branch write
\(a=2A\), \(q=2Q\), and put \(r=3n\).  Then
\begin{equation}
 F(a,q,n)+2
 =-2\bigl(r^2-2A^2+6Q^2-1\bigr).
 \label{koide-count-eq-two-adic-even-branch}
\end{equation}
Hence we must solve
\[
 r^2\equiv1+2A^2-6Q^2\pmod{2^j}.
\]
An odd unit is a square modulo \(2^j\), \(j\geq3\), precisely when it is
\(1\pmod8\), in which case it has four roots.  Here that condition is
equivalent to \(A,Q\) both being even.  There are \(2^{2j-2}\) such pairs
\((A,Q)\bmod2^j\), four choices of \(r\bmod2^j\), and, because \(3\) is
a unit, two odd choices of \(n\bmod2^k\) above each such \(r\).  This
branch therefore contributes
\begin{equation}
 2^{2j-2}\cdot4\cdot2=2^{2k-1}
 \label{koide-count-eq-two-adic-first-branch-count}
\end{equation}
solutions modulo \(2^k\).

On the second branch write \(n=2m\) and set
\begin{equation}
 v=\frac{a-q}{2},
 \qquad
 r=\frac{-a+3q}{2}.
 \label{koide-count-eq-two-adic-divided-coordinates}
\end{equation}
Then \(r,v\) have opposite parity, the inverse formula is
\(a=r+3v\), \(q=r+v\), and
\begin{equation}
 F(a,q,n)+2
 =-2\bigl(r^2-3v^2+36m^2-1\bigr).
 \label{koide-count-eq-two-adic-odd-branch}
\end{equation}
First suppose \(v=2V\) and \(r\) is odd.  For fixed
\(V\bmod2^{j-1}\) and \(m\bmod2^j\), the equation for \(r\) is
\[
 r^2\equiv1+12V^2-36m^2\pmod{2^j}.
\]
Its right side is \(1\pmod8\) exactly when \(V,m\) have the same parity.
There are
\(2\cdot2^{j-2}2^{j-1}=2^{2j-2}\) such pairs, and four roots \(r\),
so this subbranch gives \(2^{2j}\) triples \((r,v,m)\bmod2^j\).

Now suppose \(r=2R\) and \(v\) is odd.  The congruence in
\eqref{koide-count-eq-two-adic-odd-branch}, solved for \(v^2\), has an
odd square on the left; its right side is \(1\pmod8\) exactly when
\(R,m\) have opposite parity.  Again there are \(2^{2j-2}\) admissible
pairs and four square roots, hence \(2^{2j}\) transformed triples.

It remains to account for the divided-coordinate map.  Each class
\((r,v)\bmod2^j\) of opposite parity has exactly two inverse classes
\((a,q)\bmod2^k\): shifting either \(r\) or \(v\) by \(2^j\) changes both
\(a=r+3v\) and \(q=r+v\) by \(2^j\), while shifting both changes neither
modulo \(2^k\).  The class \(m\bmod2^j\) gives one even class
\(n=2m\bmod2^k\).  Thus the second branch contributes
\begin{equation}
 2\bigl(2^{2j}+2^{2j}\bigr)=2^{2k}.
 \label{koide-count-eq-two-adic-second-branch-count}
\end{equation}
Combining \eqref{koide-count-eq-two-adic-first-branch-count} and
\eqref{koide-count-eq-two-adic-second-branch-count} gives
\begin{equation}
 N_{2,k}=3\cdot2^{2k-1}\qquad(k\geq4),
 \label{koide-count-eq-two-adic-total-count}
\end{equation}
and hence \(\sigma_2=3/2\).

For \(p\geq5\), one has \(p\nmid108\).  Borovoi's good-prime formula
\cite[Section~1.3]{Borovoi_2001} gives
\[
 \sigma_p=\frac{1-p^{-2}}{1-\chi(p)p^{-1}},
 \qquad
 \chi(p)=\left(\frac{-(-2)\det A}{p}\right)
 =\left(\frac{108}{p}\right)=\left(\frac3p\right),
\]
which proves the remaining formula in
\eqref{koide-count-eq-all-local-densities}.  Since
\(\chi(p)=\chi_{12}(p)\) for \(p\geq5\),
\[
 \prod_{p\geq5}\sigma_p
 =L(1,\chi_{12})\prod_{p\geq5}(1-p^{-2})
 =\frac32\frac{L(1,\chi_{12})}{\zeta(2)}.
\]
Multiplication by \(\sigma_2\sigma_3=3\) gives the first equality in
\eqref{koide-count-eq-singular-series}.  Finally, the residue classes on
which \(\chi_{12}\) is nonzero give
\begin{align*}
 L(1,\chi_{12})
 &=\sum_{m\geq0}
 \left(
 \frac1{12m+1}-\frac1{12m+5}
 -\frac1{12m+7}+\frac1{12m+11}
 \right)\\
 &=\int_0^1\frac{1-x^2}{x^4-x^2+1}\,\mathrm dx
 =\frac{\log(2+\sqrt3)}{\sqrt3},
\end{align*}
where the middle equality follows by summing the geometric series and the
last by an elementary partial-fraction integration.  This completes the
local computation.
\end{proof}

\subsection{The archimedean chamber factor}
\label{koide-count-subsec-archimedean}

The Poincar\'e-residue gauge form on the whole smooth affine quadric is the
unique regular \(2\)-form \(\omega_\infty\) whose local representatives
satisfy
\begin{equation}
 \omega_\infty\wedge\mathrm dF
 =\mathrm da\wedge\mathrm dq\wedge\mathrm dn,
 \label{koide-count-eq-gauge-form}
\end{equation}
Explicitly, on the three standard charts it is
\begin{equation}
 \omega_\infty
 =\frac{\mathrm dq\wedge\mathrm dn}{2a}\quad(a\neq0),
 \qquad
 \omega_\infty
 =\frac{\mathrm da\wedge\mathrm dn}{6q}\quad(q\neq0),
 \qquad
 \omega_\infty
 =-\frac{\mathrm da\wedge\mathrm dq}{36n}\quad(n\neq0).
 \label{koide-count-eq-global-gauge-charts}
\end{equation}
These expressions agree on overlaps after restriction to the tangent bundle
of \(F=-2\).  Since the gradient
\((2a,-6q,-36n)\) cannot vanish on that level, the three charts cover
\(\mathscr X\), so \(\omega_\infty\) is globally defined and nowhere
vanishing.  Let \(m_\infty:=|\omega_\infty|\) be its associated positive
measure.  On the sheet \(a>0\), used for the chamber count, the first
expression in \eqref{koide-count-eq-global-gauge-charts} applies.
Let
\begin{equation}
 \mathcal C:=
 \{(a,q,n)\in\mathbb R^3:a>0,\ q\geq0,\ n\geq0,\ a+4n>q\}.
 \label{koide-count-eq-real-chamber}
\end{equation}

\begin{theorem}[Exact archimedean factor]
\label{koide-count-thm-archimedean-factor}
As \(X\to\infty\),
\begin{equation}
 m_\infty\{u\in\mathscr X(\mathbb R)\cap\mathcal C:
              4a+17n\leq X\}
 \sim I_\infty X,
 \qquad
 I_\infty=\frac1{2\sqrt3}
 \log\!\left(\frac3{2\sqrt2}\right).
 \label{koide-count-eq-archimedean-asymptotic}
\end{equation}
\end{theorem}

\begin{proof}
First, the wall \(a+4n>q\) removes only a bounded part of the positive
orthant.  Indeed, if \(q\geq a+4n\), then \(a\leq q-4n\); using
\(a^2=3q^2+18n^2-2\) gives
\[
 3q^2+18n^2-2\leq(q-4n)^2,
 \qquad\text{hence}\qquad
 q^2+4qn+n^2\leq1.
\]
Thus the removed real region has bounded \(m_\infty\)-measure.  Its only
integral point on \(\mathscr X\) with \(a>0,q,n\geq0\) is
\((a,q,n)=(1,1,0)\).

Introduce polar coordinates in the negative plane by
\begin{equation}
 \sqrt3\,q=r\cos\theta,
 \qquad
 3\sqrt2\,n=r\sin\theta,
 \qquad
 0\leq\theta\leq\frac\pi2.
 \label{koide-count-eq-polar-coordinates}
\end{equation}
On \(F=-2\), one has \(a=\sqrt{r^2-2}\) and \(r\,\mathrm dr=a\,\mathrm da\).
The Jacobian in \eqref{koide-count-eq-polar-coordinates} therefore gives
\begin{equation}
 \mathrm dm_\infty
 =\frac1{6\sqrt6}\,\mathrm da\,\mathrm d\theta.
 \label{koide-count-eq-polar-measure}
\end{equation}
At fixed \(\theta\), the height boundary is
\begin{equation}
 4a+\frac{17}{3\sqrt2}\sqrt{a^2+2}\sin\theta=X.
 \label{koide-count-eq-height-boundary}
\end{equation}
Let \(A_X(\theta)\) be its nonnegative solution when one exists, and put
it equal to zero otherwise.  Then \(0\leq A_X(\theta)/X\leq1/4\), and
for every fixed \(\theta\),
\[
 \frac{A_X(\theta)}{X}
 \longrightarrow
 \frac1{4+\frac{17}{3\sqrt2}\sin\theta}.
\]
Dominated convergence and \eqref{koide-count-eq-polar-measure} give
\begin{equation}
 I_\infty
 =\frac1{6\sqrt6}\int_0^{\pi/2}
 \frac{\mathrm d\theta}
 {4+\frac{17}{3\sqrt2}\sin\theta}.
 \label{koide-count-eq-archimedean-integral}
\end{equation}
For completeness, set \(t=\tan(\theta/2)\).  If
\(\alpha=2\sqrt2/3\) and \(\beta=3/(2\sqrt2)\), then
\(\alpha\beta=1\), and
\[
 \int_0^{\pi/2}
 \frac{\mathrm d\theta}
 {4+\frac{17}{3\sqrt2}\sin\theta}
 =\frac12\int_0^1\frac{\mathrm dt}{(t+\alpha)(t+\beta)}
 =3\sqrt2\log\!\left(\frac3{2\sqrt2}\right).
\]
Substitution in \eqref{koide-count-eq-archimedean-integral} proves
\eqref{koide-count-eq-archimedean-asymptotic}.
\end{proof}

\subsection{A genuine norm and the homogeneous counting theorem}
\label{koide-count-subsec-gos}

\begin{lemma}[A genuine norm realizing the height]
\label{koide-count-lem-genuine-norm}
Define
\begin{equation}
 \lVert(a,q,n)\rVert_*
 :=\max\{|4a+17n|,|q|,|a|\}.
 \label{koide-count-eq-star-norm}
\end{equation}
Then \(\lVert\cdot\rVert_*\) is a norm on \(\mathbb R^3\), and on the
chamber \(\mathcal C\) it satisfies
\begin{equation}
 \lVert(a,q,n)\rVert_*=4a+17n.
 \label{koide-count-eq-norm-is-height}
\end{equation}
For every integer \(X\geq0\),
\begin{equation}
 R(X)=\#\{u\in\mathscr X(\mathbb Z)\cap\mathcal C:
                    \lVert u\rVert_*<X+1\}.
 \label{koide-count-eq-r-as-sector-count}
\end{equation}
\end{lemma}

\begin{proof}
The linear map
\[
 (a,q,n)\longmapsto(4a+17n,q,a)
\]
has determinant \(-17\).  Thus \eqref{koide-count-eq-star-norm} is the
pullback of the supremum norm by an invertible linear map and is a genuine
norm.  If \((a,q,n)\in\mathcal C\), then
\[
 0<a<4a+17n,
 \qquad
 0\leq q<a+4n<4a+17n,
\]
which proves \eqref{koide-count-eq-norm-is-height}.  Since \(4a+17n\)
is integral on \(\mathscr X(\mathbb Z)\), equation
\eqref{koide-count-eq-r-as-sector-count} follows from
\eqref{koide-count-eq-exact-count}.
\end{proof}

\begin{lemma}[Verification of the Gorodnik--Oh--Shah hypotheses]
\label{koide-count-lem-gos-hypotheses}
Let \(\mathbf G=\operatorname{Spin}(F)\), acting through its standard
three-dimensional orthogonal representation on \(W=\mathbb R^3\), and let
\(x_0=(1,1,0)\).  Then
\[
 \mathscr X(\mathbb R)=\mathbf G(\mathbb R)^\circ x_0.
\]
This orbit is a closed affine symmetric variety; \(\mathbf G\) is connected
and absolutely almost simple over \(\mathbb Q\), is isotropic over
\(\mathbb R\), and acts irreducibly on \(W\) over \(\mathbb R\); and the
stabilizer \(\mathbf H=\operatorname{Stab}_{\mathbf G}(x_0)\) has no
nontrivial \(\mathbb Q\)-characters.  The cone \(\mathcal C\) in
\eqref{koide-count-eq-real-chamber} is admissible and generic for
\cite[Theorem~1.6]{GorodnikOhShah_2009}.
\end{lemma}

\begin{proof}
The form \(F\) has signature \((1,2)\).  Hence \(\mathbf G\) is a
connected absolutely almost simple group of type \(A_1\), it is
\(\mathbb R\)-isotropic, and its standard three-dimensional representation
is \(\mathbb R\)-irreducible.  The nonzero level \(\mathscr X\) is
Zariski closed.  With
\[
 Q=\sqrt3\,q,
 \qquad
 N=3\sqrt2\,n,
\]
its real equation becomes
\[
 Q^2+N^2-a^2=2.
\]
The parametrization
\[
 (a,\theta)\longmapsto
 \bigl(a,\sqrt{a^2+2}\cos\theta,
          \sqrt{a^2+2}\sin\theta\bigr),
 \qquad (a,\theta)\in\mathbb R\times(\mathbb R/2\pi\mathbb Z),
\]
shows that \(\mathscr X(\mathbb R)\) is connected.  Witt extension first
produces an element \(t\in\operatorname O(F)(\mathbb R)\) taking any
prescribed point of the level set to any other.  If \(\det t=-1\), choose
an anisotropic vector in the orthogonal complement of the target point and
compose \(t\) on the left with its reflection.  That reflection fixes the
target point and has determinant \(-1\); hence the composite lies in
\(\operatorname{SO}(F)(\mathbb R)\).  Thus
\(\operatorname{SO}(F)(\mathbb R)\) is transitive on the level set.  The
identity-component orbits have the full dimension of the level surface and
are open; because their complements are unions of open orbits, each is also
closed.  Connectedness therefore leaves only one such orbit, proving
\(\mathscr X(\mathbb R)=\mathbf G(\mathbb R)^\circ x_0\).  The stabilizer
of a nonisotropic vector is the fixed subgroup of the involution induced by
reflection in that vector, so the orbit is affine symmetric.

The orthogonal complement of \(x_0\) is
\begin{equation}
 x_0^\perp=\{(3s,s,t):s,t\in\mathbb Q\},
 \qquad
 F(3s,s,t)=6(s^2-3t^2).
 \label{koide-count-eq-rational-complement}
\end{equation}
Thus \(\mathbf H\) is the norm-one torus of
\(K=\mathbb Q(\sqrt3)\).  Its character lattice is \(\mathbb Z\), with
the nontrivial element of \(\operatorname{Gal}(K/\mathbb Q)\) acting by
\(-1\); consequently it has no nonzero Galois-invariant character and no
nontrivial \(\mathbb Q\)-character.  The integral point \(x_0\) supplies
the required nonempty integral locus.

It remains to check the sector rather than merely the affine orbit.  For the
given representation and any Euclidean sphere in \(W\), the Satake boundary
is
\begin{equation}
 \mathscr X^\infty=\{F=0\}\cap S(W)
 =\mathcal O_+\sqcup\mathcal O_-,
 \qquad
 \mathcal O_\pm:=\{F=0,\ \pm a>0\}\cap S(W),
 \label{koide-count-eq-satake-boundary}
\end{equation}
and \(\mathbf G(\mathbb R)^\circ\) is transitive on each displayed circle.
For a nonzero null vector with \(a,q,n\geq0\),
\[
 a=\sqrt{3q^2+18n^2}>q-4n.
\]
Hence both the closure and the relative interior of \(\mathcal C\) meet
exactly \(\mathcal O_+\), not \(\mathcal O_-\).  The intersection of
\(\partial\mathcal C\) with the Satake boundary consists of the two radial
directions
\begin{equation}
 (a,q,n)=(\sqrt3,1,0),
 \qquad
 (a,q,n)=(3\sqrt2,0,1).
 \label{koide-count-eq-sector-boundary-rays}
\end{equation}
These are exactly the two projective limits of the integral Pell families
in Proposition~\ref{koide-count-prop-pell-faces}.
This two-point intersection has measure zero for the smooth measure class
on the boundary orbit.
This is precisely the admissibility and genericity condition required by
the cited sector theorem.
\end{proof}

\subsection{The 2-adic orbit sign and Borovoi's coefficient}
\label{koide-count-subsec-borovoi}

Write $\widehat{\mathbb Z}=\prod_p\mathbb Z_p$.  For the nonsquare
ternary quadric under consideration, Borovoi associates to every local
point an orbit sign $\nu_v\in\{\pm1\}$, constant on
$\operatorname{Spin}(F)(\mathbb Q_v)$-orbits and equal to $+1$ at almost
every finite place.  Thus
$\nu_f((x_p)_p)=\prod_p\nu_p(x_p)$ is defined on
$\mathscr X(\widehat{\mathbb Z})$.  Because $\mathscr X(\mathbb R)$ is
connected, his normalization reads
\begin{equation}
 c_{\mathscr X}
 =\frac{2m_f\{x\in\mathscr X(\widehat{\mathbb Z}):\nu_f(x)=+1\}}
 {m_f(\mathscr X(\widehat{\mathbb Z}))},
 \label{koide-count-eq-borovoi-sign-definition}
\end{equation}
where $m_f$ is the finite Tamagawa measure attached to the gauge form
\eqref{koide-count-eq-gauge-form} and Borovoi's convergence factors.

\begin{proposition}[Finite adelic sign balance]
\label{koide-count-prop-finite-sign-balance}
For the affine quadric \(\mathscr X=\{F=-2\}\), the coefficient
\(c_{\mathscr X}\) in Borovoi's counting formula
\cite[Theorem~1.4.1 and Proposition~1.5.1]{Borovoi_2001} is exactly
\begin{equation}
 c_{\mathscr X}=1.
 \label{koide-count-eq-borovoi-coefficient}
\end{equation}
\end{proposition}

\begin{proof}
Continue to write \(A=\operatorname{diag}(1,-3,-18)\).  Put
\(e=(1,0,0)\), \(v=(1,0,1)\), and for every anisotropic vector \(w\) let
\[
 r_w(u)=u-\frac{2u^{\mathsf T}Aw}{w^{\mathsf T}Aw}\,w
\]
be its orthogonal reflection.  The element \(g=r_vr_e\) is
\begin{equation}
 g=
 \begin{pmatrix}
  -19/17&0&-36/17\\
  0&1&0\\
  -2/17&0&-19/17
 \end{pmatrix}.
 \label{koide-count-eq-sign-toggle-matrix}
\end{equation}
Direct multiplication gives
\begin{equation}
 g^{\mathsf T}Ag=A,
 \qquad
 \det g=1.
 \label{koide-count-eq-sign-toggle-orthogonal}
\end{equation}
Because \(17\in\mathbb Z_2^\times\), both \(g\) and \(g^{-1}\) have entries
in \(\mathbb Z_2\), so \(g\in\operatorname{SO}(F,\mathbb Z_2)\).  Its
spinor norm is
\begin{equation}
 \vartheta(g)=F(v)F(e)=-17
 \quad\text{in }\mathbb Q_2^\times/\mathbb Q_2^{\times2}.
 \label{koide-count-eq-toggle-spinor-norm}
\end{equation}

We identify the stabilizer's spinor-norm image.  By
\eqref{koide-count-eq-rational-complement}, over \(\mathbb Q_2\) the
orthogonal complement of \(x_0\) is
\begin{equation}
 F|_{x_0^\perp}
 =6(s^2-3t^2)
 =6N_{K_2/\mathbb Q_2}(s+t\sqrt3),
 \qquad K_2=\mathbb Q_2(\sqrt3).
 \label{koide-count-eq-two-adic-complement}
\end{equation}
The odd unit \(3\) is not a square in \(\mathbb Q_2\), so \(K_2\) is a
quadratic field and this binary plane is anisotropic.  Put
\(u_0=(3,1,0)\), so \(F(u_0)=6\).  If
\(h\in\operatorname{SO}(F|_{x_0^\perp})(\mathbb Q_2)\), then
\(hr_{u_0}\) is a two-dimensional orthogonal transformation of determinant
\(-1\), hence is a reflection \(r_w\) in some nonisotropic vector
\(w=(3s,s,t)\).  Therefore \(h=r_wr_{u_0}\) and
\[
 \vartheta(h)=F(w)F(u_0)
 =36N_{K_2/\mathbb Q_2}(s+t\sqrt3)
 \quad\bmod\mathbb Q_2^{\times2}.
\]
Conversely these products realize every norm square class as \(w\) ranges
through the nonzero vectors of the anisotropic plane.  Hence
\begin{equation}
 \vartheta\!\left(
 \operatorname{Stab}_{\operatorname{SO}(F)}(x_0)(\mathbb Q_2)
 \right)
 =N_{K_2/\mathbb Q_2}(K_2^\times)
 \quad\bmod\mathbb Q_2^{\times2}.
 \label{koide-count-eq-stabilizer-spinor-image}
\end{equation}

For a quadratic extension, its norm group is the kernel of the Hilbert
symbol with the defining square class
\cite[Corollary~13.4.6]{Voight_2021}.  The odd-unit formula
\cite[Equation~12.4.14]{Voight_2021} gives
\begin{equation}
 (-17,3)_2
 =(-1)^{((-17-1)/2)((3-1)/2)}=-1.
 \label{koide-count-eq-minus17-hilbert-symbol}
\end{equation}
Thus \(-17\notin N_{K_2/\mathbb Q_2}(K_2^\times)\).

By Witt extension, every \(x\in\mathscr X(\mathbb Q_2)\) is \(tx_0\) for
some \(t\in\operatorname O(F)(\mathbb Q_2)\).  If \(\det t=-1\), choose
an anisotropic vector in \(x^\perp\) and compose on the left with its
reflection; this fixes \(x\) and changes the determinant.  Hence every
such \(x\) has the form \(x=sx_0\) with
\(s\in\operatorname{SO}(F)(\mathbb Q_2)\).  The image of
\(\operatorname{Spin}(F)(\mathbb Q_2)\) in
\(\operatorname{SO}(F)(\mathbb Q_2)\) is the kernel of the spinor norm,
as used in \cite[Lemma~3.1]{Borovoi_2001}.  Consequently the
\(\operatorname{Spin}(F)(\mathbb Q_2)\)-orbit of \(x=sx_0\) is determined
by
\begin{equation}
 \vartheta(s)\in
 \mathbb Q_2^\times/
 \bigl(\mathbb Q_2^{\times2}
 N_{K_2/\mathbb Q_2}(K_2^\times)\bigr),
 \label{koide-count-eq-two-adic-orbit-quotient}
\end{equation}
a quotient of order two.  Equations
\eqref{koide-count-eq-toggle-spinor-norm} and
\eqref{koide-count-eq-minus17-hilbert-symbol} show that multiplication by
\(g\) exchanges these two orbits.  In Borovoi's convention, the orbit of
\(x_0\) has sign \(\nu_2=+1\) and the other orbit has sign
\(\nu_2=-1\).  Hence
\begin{equation}
 \nu_2(gx)=-\nu_2(x)
 \qquad(x\in\mathscr X(\mathbb Q_2)).
 \label{koide-count-eq-local-sign-toggle}
\end{equation}

Apply \(g\) at the prime \(2\) and the identity at every other finite
prime.  This gives a bijection of \(\mathscr X(\widehat{\mathbb Z})\).  It
preserves the finite Tamagawa measure: at \(2\) this follows from invariance
of the local measure, equivalently from \(g^*F=F\) and \(\det g=1\), and at
all other primes the map is the identity.  Equation
\eqref{koide-count-eq-local-sign-toggle} shows that it exchanges the finite
sign loci \(\nu_f=+1\) and \(\nu_f=-1\).  They therefore have equal
measure.  Since \(\mathscr X(\mathbb R)\) is connected, Borovoi's formula
becomes
\[
 c_{\mathscr X}
 =\frac{2m_f(\mathscr X(\widehat{\mathbb Z})_+)}
        {m_f(\mathscr X(\widehat{\mathbb Z}))}=1.
\]
The argument uses a measure-preserving, sign-reversing bijection; it does
not require \(g\) to be an involution (indeed \(g^2\neq I\)).
\end{proof}

\subsection{Calibration and the counting constant}
\label{koide-count-subsec-calibration}

\begin{theorem}[Calibrated reduced-root asymptotic]
\label{koide-count-thm-linear-asymptotic}
The number of reduced positive defect-one roots satisfies
\begin{equation}
 R(X)\sim C_0X\qquad(X\to\infty),
 \label{koide-count-eq-final-asymptotic}
\end{equation}
where
\begin{equation}
 \begin{aligned}
 C_0
 &=\mathfrak S(\mathscr X)I_\infty\\
 &=\frac9{2\pi^2}\log(2+\sqrt3)
   \log\!\left(\frac3{2\sqrt2}\right)\\
 &=0.035362047799731853\ldots.
 \end{aligned}
 \label{koide-count-eq-final-constant}
\end{equation}
\end{theorem}

\begin{proof}
Let \(\mu_{\mathrm{GOS}}\) denote the invariant measure on
\(\mathscr X(\mathbb R)\) with the normalization appearing in the
Gorodnik--Oh--Shah sector theorem
\cite[Theorem~1.6]{GorodnikOhShah_2009}.  The gauge measure \(m_\infty\)
in \eqref{koide-count-eq-gauge-form} is also a nonzero
\(\mathbf G(\mathbb R)^\circ\)-invariant Radon measure.  Because
\(\mathscr X(\mathbb R)\) is a single homogeneous orbit, uniqueness of
invariant measure gives
\begin{equation}
 \mu_{\mathrm{GOS}}=\kappa m_\infty
 \label{koide-count-eq-measure-scalar}
\end{equation}
for some \(\kappa>0\).

The scalar is fixed by applying both counting theorems to the same Euclidean
balls and the full cone \(W\).  Gorodnik--Oh--Shah give
\[
 \#\{x\in\mathscr X(\mathbb Z):|x|_2<T\}
 \sim\kappa\,m_\infty
 \{x\in\mathscr X(\mathbb R):|x|_2<T\}.
\]
Borovoi's Theorem~1.4.1 and Proposition~1.5.1, together with
Proposition~\ref{koide-count-prop-finite-sign-balance}, give for the same
balls
\[
 \#\{x\in\mathscr X(\mathbb Z):|x|_2<T\}
 \sim\mathfrak S(\mathscr X)\,m_\infty
 \{x\in\mathscr X(\mathbb R):|x|_2<T\}.
\]
The real volume here is asymptotic to a positive constant times \(T\).
Comparison therefore gives the invariant-measure calibration
\begin{equation}
 \kappa=\mathfrak S(\mathscr X).
 \label{koide-count-eq-measure-calibration}
\end{equation}

Lemma~\ref{koide-count-lem-genuine-norm} and
Lemma~\ref{koide-count-lem-gos-hypotheses} now allow the sector theorem to
be applied with the actual norm and chamber of the reduced-root problem:
\begin{align*}
 R(X)
 &\sim\mu_{\mathrm{GOS}}
 \{u\in\mathscr X(\mathbb R)\cap\mathcal C:
       \lVert u\rVert_*\leq X\}\\
 &=\mathfrak S(\mathscr X)\,
 m_\infty\{u\in\mathscr X(\mathbb R)\cap\mathcal C:
       4a+17n\leq X\}\\
 &\sim\mathfrak S(\mathscr X)I_\infty X.
\end{align*}
The switch between strict and weak inequalities is harmless by a
\((1\pm\varepsilon)X\) squeeze; on integral points it also follows directly
from \eqref{koide-count-eq-r-as-sector-count}.  Since
\(R(X)=R(\lfloor X\rfloor)\), the integer-height asymptotic extends to all
real \(X\to\infty\).  Finally, substituting
\eqref{koide-count-eq-singular-series} and
\eqref{koide-count-eq-archimedean-asymptotic}, and using
\(\zeta(2)=\pi^2/6\), gives exactly
\eqref{koide-count-eq-final-constant}.
\end{proof}

\subsection{Exact computation and the remaining open problem}
\label{koide-count-subsec-computation}

The first five reduced roots are
\begin{equation}
 (1,4,20),\quad(4,4,33),\quad(4,8,49),\quad
 (4,15,76),\quad(8,24,129).
 \label{koide-count-eq-first-roots}
\end{equation}
The first root is the first admissible term of the $n=0$ Pell family and
the second is the first term of the $q=0$ Pell family in
Proposition~\ref{koide-count-prop-pell-faces}.
Exact integer enumeration in the ternary coordinates gives the following
values.  The residual column uses the full constant in
\eqref{koide-count-eq-final-constant}, not a fitted constant.
\begin{center}
\begin{tabular}{r@{\qquad}r@{\qquad}r}
\toprule
\(X\)&\(R(X)\)&\(R(X)-C_0X\)\\
\midrule
10,000&353&\(-0.62048\)\\
20,000&715&\(7.75904\)\\
50,000&1,755&\(-13.10239\)\\
75,000&2,654&\(1.84642\)\\
100,000&3,535&\(-1.20478\)\\
125,000&4,460&\(39.74403\)\\
150,000&5,294&\(-10.30717\)\\
175,000&6,195&\(6.64164\)\\
200,000&7,109&\(36.59044\)\\
\bottomrule
\end{tabular}
\end{center}
The enumeration loops over \(n\geq0\) and \(a>0\) with
\(4a+17n\leq X\), tests whether
\[
 q^2=\frac{a^2-18n^2+2}{3}
\]
is the square of a nonnegative integer satisfying \(q<a+4n\), and then
reconstructs \((x,y,z)\) with
\eqref{koide-count-eq-inverse-coordinates}.  Every retained triple is
rechecked in the original polynomial \eqref{koide-count-eq-defect} and
against all ordering and wall inequalities.  Separate exhaustive congruence
enumerations agree with the formulae above at
\(p=2,3,5,7,11,13\) through the computed powers.  These finite checks
verify the implementation; the asymptotic theorem rests on the proofs above,
not on numerical fitting.

The reproducibility script is
\begin{center}
 \path{supporting_materials/workbench/reproduction/t1_koide_exact/verify_koide_narain_root_density.py}.
\end{center}
Running it with \texttt{--max-z 200000} emits the coordinate, local-density,
$2$-adic isometry, and bounded-enumeration certificates recorded in
\begin{center}
 \path{supporting_materials/workbench/reproduction/t1_koide_exact/koide_verification_200000.json}.
\end{center}
The script uses vectorized square detection only to propose candidates; it
rechecks every retained candidate with integer arithmetic in both coordinate
systems and in the original polynomial $\Delta$.

\begin{test}[Square-root discrepancy]
\label{koide-count-open-square-root-discrepancy}
\Open Determine whether, for every \(\varepsilon>0\),
\begin{equation}
 R(X)-C_0X=O_\varepsilon(X^{1/2+\varepsilon}).
 \label{koide-count-eq-square-root-discrepancy}
\end{equation}
The homogeneous-dynamics argument above and the exact
effective-well-roundedness verification in
\cref{koide-param-thm:power-saving} below prove, for some fixed
\(\eta>0\),
\[
 R(X)=C_0X+O(X^{1-\eta}),
\]
but do not prove
\eqref{koide-count-eq-square-root-discrepancy}.  No unproved spectral gap,
GRH assumption, or implicit ``conditional theorem'' is being used here:
the sharper error term is an open problem in this workbench; see the
primary-source comparison in
\cref{koide-lit-determination}.
\end{test}

% Integrated 4 September 2026 from the isolated exact-parametrization audit.
% Every label is namespaced koide-param-*; the original variables, chamber,
% strict walls, lattice, and height are unchanged.
\providecommand{\OO}{\mathcal O}
\providecommand{\NN}{\mathbb N}
\providecommand{\ZZ}{\mathbb Z}
\providecommand{\RR}{\mathbb R}
\providecommand{\QQ}{\mathbb Q}
\providecommand{\Norm}{\operatorname N}
\providecommand{\Tr}{\operatorname {Tr}}
\providecommand{\e}{\varepsilon}

\section{Exact ideal and unit-orbit parametrization of the unchanged sharp Koide count}

Throughout this section the variables, equation, chamber, and height are
exactly
\begin{equation}
 a^2-3q^2-18n^2=-2,
 \qquad
 a>0,\quad q\geq0,\quad n\geq0,\quad a+4n>q,
 \qquad
 4a+17n\leq X.
 \label{koide-param-eq:original-problem}
\end{equation}
Every auxiliary symbol introduced below is connected to these same integer
points by an explicitly proved map and inverse.  No count for a normalized
form, a rescaled lattice, or a substituted cutoff is used in place of
\eqref{koide-param-eq:original-problem}.
Thus
\begin{equation}
 R(X)=\#\left\{(a,q,n)\in\ZZ^3:\text{all conditions in
 \eqref{koide-param-eq:original-problem} hold}\right\}.
 \label{koide-param-eq:R-definition}
\end{equation}
The already calibrated leading constant is
\begin{equation}
 C_0=
 \frac9{2\pi^2}\log(2+\sqrt3)
 \log\!\left(\frac3{2\sqrt2}\right).
 \label{koide-param-eq:C0}
\end{equation}

The exact parametrization proved below has four pieces of data for each
fixed $n\geq1$: a positive integer $C$, a positive integer $D$, a square
root $E$ of $3$ modulo $D$, and a nonnegative Pell exponent $k$.  There is
no quotient left implicit: equations \eqref{koide-param-eq:bounded-seed-search},
\eqref{koide-param-eq:HNF-to-seed}, and \eqref{koide-param-eq:unit-orbit-coordinates} construct
$(a,q,n)$ by finite integer operations, while
\eqref{koide-param-eq:forward-HNF} and the descent \eqref{koide-param-eq:inverse-descent} recover
$(C,D,E,k)$ from $(a,q,n)$.  The face $n=0$ is recorded separately and
exactly.

This parametrization gives the exact sharp-cutoff formula
\eqref{koide-param-eq:exact-floor-formula}.  It also separates the discrepancy into
the radial shifted-divisor sum and the unit-phase discrepancy in
\eqref{koide-param-eq:exact-discrepancy-splitting}.  The square-root estimate
\begin{equation}
 R(X)-C_0X=O_\e(X^{1/2+\e})
 \label{koide-param-eq:square-root-target}
\end{equation}
is not proved or refuted here.  What is proved unconditionally, by applying
the effective symmetric-variety theorem to the exact chamber, is that
there is a constant $\eta>0$ for which
\begin{equation}
 R(X)=C_0X+O(X^{1-\eta}).
 \label{koide-param-eq:unconditional-power-saving}
\end{equation}
The cited general theorem does not supply $\eta\geq1/2$, and its proof does
not justify replacing \eqref{koide-param-eq:unconditional-power-saving} by
\eqref{koide-param-eq:square-root-target}.

\section{Arithmetic of $\ZZ[\sqrt3]$}

Put
\begin{equation}
 K=\QQ(\sqrt3),\qquad \OO=\ZZ[\sqrt3],\qquad
 \epsilon=2+\sqrt3.
 \label{koide-param-eq:field-data}
\end{equation}
For $\alpha=u+v\sqrt3\in K$, write
\begin{equation}
 \alpha'=u-v\sqrt3,
 \qquad
 \Norm(\alpha)=\alpha\alpha'=u^2-3v^2.
 \label{koide-param-eq:conjugate-norm}
\end{equation}

\begin{lemma}[Ring, class number, and totally positive units]
\label{koide-param-lem:ring-units}
The ring of integers of $K$ is $\OO$.  Every nonzero integral ideal of
$\OO$ is principal.  Every unit has norm $+1$, and the totally positive
unit group is
\begin{equation}
 \OO^{\times,+}=\{\epsilon^j:j\in\ZZ\}.
 \label{koide-param-eq:positive-units}
\end{equation}
\end{lemma}

\begin{proof}
Since $3\equiv3\pmod4$, the standard integral-basis calculation for a
quadratic field gives $\OO_K=\ZZ[\sqrt3]$ and discriminant $12$.
Minkowski's bound says that every ideal class in a real quadratic field of
discriminant $12$ has an integral representative of norm at most
\(\frac{2!}{2^2}\sqrt{12}=\sqrt3<2\).  Its norm must therefore be $1$,
so it is $\OO$ and the ordinary class number is one.

The norm of a unit is $\pm1$.  The equation $u^2-3v^2=-1$ is impossible
modulo $3$, because it would give $u^2\equiv2\pmod3$.  Hence every unit
has norm $+1$.  Dirichlet's unit theorem makes the positive units an
infinite cyclic group.  To identify its least generator, suppose
$\gamma=u+v\sqrt3>1$ and $\gamma'>0$.  Replacing $\gamma$ by its inverse
if necessary makes $v\geq0$.  If $v=0$, then $\gamma=1$.  If $v\geq1$,
then $u^2-3v^2=1$ forces $u\geq2$, and consequently
$\gamma\geq2+\sqrt3=\epsilon$.  Since $\epsilon$ itself has norm one,
it is the least positive unit greater than one and generates the group.
\end{proof}

For $m\geq1$ define
\begin{equation}
 \chi_{12}(d)=\left(\frac{12}{d}\right),
 \qquad
 \rho(m)=\sum_{d\mid m}\chi_{12}(d).
 \label{koide-param-eq:rho-definition}
\end{equation}
Here the symbol is Kronecker's symbol; in particular
\begin{equation}
 \chi_{12}(d)=
 \begin{cases}
  0,&2\mid d\text{ or }3\mid d,\\
  1,&d\equiv1,11\pmod {12},\\
  -1,&d\equiv5,7\pmod {12}.
 \end{cases}
 \label{koide-param-eq:chi12-table}
\end{equation}

\begin{lemma}[Ideal number]
\label{koide-param-lem:ideal-number}
The number of integral ideals of $\OO$ having absolute norm $m$ is
$\rho(m)$.
\end{lemma}

\begin{proof}
The Dedekind zeta function factors as
$\zeta_K(s)=\zeta(s)L(s,\chi_{12})$.  Comparing Dirichlet coefficients
gives the assertion.  Equivalently, the equality follows prime by prime.
At $2$ and $3$ there is one ramified prime ideal, so there is one ideal of
norm $p^e$ for every $e\geq0$.  If $\chi_{12}(p)=1$, then $p$ splits and
there are $e+1$ ideals of norm $p^e$.  If $\chi_{12}(p)=-1$, then $p$ is
inert, so there is one such ideal when $e$ is even and none when $e$ is
odd.  These are respectively the local values of
$\sum_{j=0}^e\chi_{12}(p)^j$.
Consequently, if
$m=2^{e_2}3^{e_3}\prod_{p\geq5}p^{e_p}$, then
\begin{equation}
 \rho(m)=
 \begin{cases}
  \displaystyle\prod_{\substack{p\mid m,\ p\geq5\\p\equiv1,11\ (12)}}(e_p+1),
   &e_p\text{ is even for every }p\equiv5,7\pmod {12},\\[6pt]
  0,&\text{otherwise}.
 \end{cases}
 \label{koide-param-eq:rho-Euler-product}
\end{equation}
\end{proof}

\section{Hermite-normal-form parameters for every ideal}

For $m\geq1$ let
\begin{equation}
 \mathcal P(m)=
 \left\{(C,D,E):
 \begin{array}{l}
 C,D\in\ZZ_{>0},\quad C^2D=m,\\
 0\leq E<D,\quad E^2\equiv3\pmod D
 \end{array}
 \right\}.
 \label{koide-param-eq:P-m}
\end{equation}

\begin{proposition}[Exact HNF parametrization of ideals]
\label{koide-param-prop:HNF-ideals}
The map
\begin{equation}
 (C,D,E)\longmapsto
 \mathfrak a(C,D,E)
 :=C\bigl(D\ZZ+(E+\sqrt3)\ZZ\bigr)
 \label{koide-param-eq:param-ideal}
\end{equation}
is a bijection from $\mathcal P(m)$ to the integral ideals of $\OO$ of
norm $m$.  Consequently
\begin{equation}
 \#\mathcal P(m)
 =\sum_{C^2D=m}\#\{E\bmod D:E^2\equiv3\pmod D\}
 =\rho(m).
 \label{koide-param-eq:root-divisor-identity}
\end{equation}
\end{proposition}

\begin{proof}
Relative to the basis $1,\sqrt3$ of $\OO$, the lattice in
\eqref{koide-param-eq:param-ideal} has column basis
\[
 \begin{pmatrix}CD&CE\\0&C\end{pmatrix}.
\]
Its index is $C^2D=m$.  It is stable under multiplication by $\sqrt3$
because
\begin{align*}
 \sqrt3(CD)&=D\,C(E+\sqrt3)-E(CD),\\
 \sqrt3\,C(E+\sqrt3)
 &=E\,C(E+\sqrt3)+\frac{3-E^2}{D}(CD),
\end{align*}
and the last coefficient is integral.  It is therefore an integral ideal
of norm $m$.

Conversely, write the unique column Hermite normal form of an integral
ideal $\mathfrak a$ as
\[
 \mathfrak a=A\ZZ+(B+C\sqrt3)\ZZ,
 \qquad A,C>0,\quad0\leq B<A.
\]
Stability of $A$ under multiplication by $\sqrt3$ implies $C\mid A$ and
$C\mid B$.  Write $A=CD$ and $B=CE$.  Stability of $B+C\sqrt3$ then
implies $D\mid3-E^2$.  The HNF inequality gives $0\leq E<D$, and the
index is $AC=C^2D$.  HNF uniqueness proves injectivity and surjectivity.

For completeness, the last equality in \eqref{koide-param-eq:root-divisor-identity}
can be checked without zeta functions.  Let
$r_3(D)=\#\{E\bmod D:E^2\equiv3\pmod D\}$.  For $p\geq5$, Hensel's
lemma gives $r_3(p^j)=2$ for every $j\geq1$ if
$\chi_{12}(p)=1$, and $r_3(p^j)=0$ if $\chi_{12}(p)=-1$.  Also
$r_3(2)=r_3(3)=1$ and $r_3(2^j)=r_3(3^j)=0$ for $j\geq2$.
Thus, for every prime power $p^e$,
\[
 \sum_{2c+j=e}r_3(p^j)=\sum_{j=0}^e\chi_{12}(p)^j.
\]
Multiplicativity proves \eqref{koide-param-eq:root-divisor-identity}.
\end{proof}

The HNF parameters can be recovered directly from a generator.  Suppose
$\alpha=a+q\sqrt3$ has positive norm $m$, with $a>0$ and $q\geq0$.  Put
\begin{equation}
 C=\gcd(a,q),\qquad A=a/C,\qquad B=q/C,\qquad D=m/C^2.
 \label{koide-param-eq:forward-HNF-pre}
\end{equation}
If $D=1$, put $E=0$.  If $D>1$, then
\begin{equation}
 \gcd(B,D)=1,\qquad E\equiv A B^{-1}\pmod D,\qquad0\leq E<D.
 \label{koide-param-eq:forward-HNF}
\end{equation}
Indeed,
$\gcd(B,D)=\gcd(B,A^2-3B^2)=\gcd(B,A^2)=1$, and
$E^2\equiv3\pmod D$.  The principal ideal $(\alpha)$ is exactly
$\mathfrak a(C,D,E)$: the element $A+B\sqrt3$ lies in
$D\ZZ+(E+\sqrt3)\ZZ$, so the principal ideal is contained in the HNF
ideal, while both ideals have norm $D$ after division by $C$.

\section{Complete integral parametrization of the original equation}

For $n\geq1$ retain the exact integer
\begin{equation}
 m_n=18n^2-2.
 \label{koide-param-eq:mn}
\end{equation}
It satisfies $m_n>0$ and $m_n\equiv1\pmod3$.

\begin{lemma}[Every ideal of norm $m_n$ has positive-norm generators]
\label{koide-param-lem:positive-norm-generator}
If $\mathfrak a$ is an integral ideal of $\OO$ of norm $m_n$, then every
generator has norm $+m_n$, rather than $-m_n$.
\end{lemma}

\begin{proof}
Lemma~\ref{koide-param-lem:ring-units} makes $\mathfrak a$ principal.  A generator
$u+v\sqrt3$ has norm $\pm m_n$.  Negative norm would give
\[
 u^2-3v^2=-m_n,
\]
and reduction modulo $3$ would give
$u^2\equiv-m_n\equiv2\pmod3$, which is impossible.  Thus the norm is
$+m_n$.
\end{proof}

For each $\mathbf p=(C,D,E)\in\mathcal P(m_n)$, choose a generator of
$\mathfrak a(C,D,E)$.  By Lemma~\ref{koide-param-lem:positive-norm-generator} its
norm is $m_n$.  Change its sign so that both real embeddings are positive,
and multiply by a unique power of $\epsilon$ so as to obtain the unique
generator
\begin{equation}
 \alpha_{n,\mathbf p}=a_{n,\mathbf p,0}
   +q_{n,\mathbf p,0}\sqrt3,
 \qquad
 \sqrt{m_n}\leq\alpha_{n,\mathbf p}
       <\epsilon\sqrt{m_n}.
 \label{koide-param-eq:canonical-generator}
\end{equation}
Existence and uniqueness follow from
\eqref{koide-param-eq:positive-units}, because the half-open intervals
$[\epsilon^j\sqrt{m_n},\epsilon^{j+1}\sqrt{m_n})$ partition
$\RR_{>0}$.  Since
$\alpha_{n,\mathbf p}'=m_n/\alpha_{n,\mathbf p}$, the lower inequality
in \eqref{koide-param-eq:canonical-generator} implies
$q_{n,\mathbf p,0}\geq0$.

\begin{lemma}[Purely integral reduced seed and descent]
\label{koide-param-lem:integral-seed-descent}
For an integral pair $(a,q)$ satisfying
\[
 a^2-3q^2=m_n,\qquad a>0,\qquad q\geq0,
\]
the half-open inequalities
\[
 \sqrt{m_n}\leq a+q\sqrt3<\epsilon\sqrt{m_n}
\]
are equivalent to the purely integral inequalities
\begin{equation}
 a>2q\geq0.
 \label{koide-param-eq:integral-seed-condition}
\end{equation}
If $a\leq2q$, multiplication by $\epsilon^{-1}=2-\sqrt3$ is the
integral inverse step
\begin{equation}
 (a,q)\longmapsto(a^-,q^-)=(2a-3q,2q-a).
 \label{koide-param-eq:inverse-descent}
\end{equation}
It preserves the norm, has $a^->0$ and $0\leq q^-<q$, and can therefore
be iterated until the unique pair satisfying \eqref{koide-param-eq:integral-seed-condition}
is reached.
\end{lemma}

\begin{proof}
Put $u=(a+q\sqrt3)/\sqrt{m_n}$.  Norm positivity and the sign conditions
give $u\geq1$.  Moreover
\[
 \frac a{\sqrt{m_n}}=\frac{u+u^{-1}}2,
 \qquad
 \frac{q\sqrt3}{\sqrt{m_n}}=\frac{u-u^{-1}}2.
\]
For $u\geq1$, the ratio
\[
 \frac aq=\sqrt3\,\frac{u^2+1}{u^2-1}
\]
(with value $+\infty$ when $q=0$) is strictly decreasing.  At
$u=\epsilon$ it equals $2$, since
$\epsilon+\epsilon^{-1}=4$.  Thus $u<\epsilon$ is equivalent to $a>2q$;
the strict inequality exactly matches the half-open upper endpoint.

If $a\leq2q$, norm positivity gives $a>\sqrt3q$, and hence
\[
 a^-=2a-3q>(2\sqrt3-3)q>0,
 \qquad q^-=2q-a\geq0.
\]
Also $a>\sqrt3q>q$ gives $q^-=2q-a<q$.  Thus the positive integer $q$
strictly decreases at every inverse step, so the process terminates.  A
terminal pair has $a>2q$.  Conversely, for such a pair the second
coordinate after one inverse step is negative, so no further inverse step
stays in the nonnegative-$q$ region.  This proves uniqueness as well as
termination.
\end{proof}

The preceding seed can be reconstructed from the HNF data by a bounded
integer search; thus the parametrization does not require an unspecified
choice of a generator.  More precisely, for
$\mathbf p=(C,D,E)\in\mathcal P(m_n)$ there is a unique pair $(A,B)$ with
\begin{equation}
 \begin{gathered}
  A,B\in\ZZ,\qquad A>0,\qquad 0\leq B,\qquad B^2<D,\\
  A^2-3B^2=D,\qquad A\equiv EB\pmod D.
 \end{gathered}
 \label{koide-param-eq:bounded-seed-search}
\end{equation}
It gives the generator in \eqref{koide-param-eq:canonical-generator} by the exact
integer formula
\begin{equation}
 (a_{n,\mathbf p,0},q_{n,\mathbf p,0})=(CA,CB).
 \label{koide-param-eq:HNF-to-seed}
\end{equation}
Indeed, existence follows by writing the reduced generator as
$C(A+B\sqrt3)$.  The seed inequality $A>2B$ gives
$D=A^2-3B^2>B^2$.  Conversely, the conditions in
\eqref{koide-param-eq:bounded-seed-search} imply $A>2B$ and put
$A+B\sqrt3$ in $D\ZZ+(E+\sqrt3)\ZZ$.  Its principal ideal is contained
in that HNF ideal, and both have norm $D$, so they are equal.  Lemma
\ref{koide-param-lem:integral-seed-descent} and HNF uniqueness give uniqueness.  In
particular one needs to test only the integers
$0\leq B\leq\lfloor\sqrt{D-1}\rfloor$ and the positive square
$A=\sqrt{D+3B^2}$.
The phase introduced below is therefore also explicit in the same integer
data:
\begin{equation}
 \tau_{n,(C,D,E)}
 =\frac{\log((A+B\sqrt3)/\sqrt D)}{\log\epsilon}.
 \label{koide-param-eq:HNF-phase-explicit}
\end{equation}
Here the factor $C$ cancels because $m_n=C^2D$; no ideal generator remains
to be chosen.

For $k\geq0$ define integers $a_{n,\mathbf p,k}$ and
$q_{n,\mathbf p,k}$ by
\begin{equation}
 a_{n,\mathbf p,k}+q_{n,\mathbf p,k}\sqrt3
 =\alpha_{n,\mathbf p}\epsilon^k.
 \label{koide-param-eq:unit-orbit-coordinates}
\end{equation}
Equivalently,
\begin{equation}
 \begin{pmatrix}a_{n,\mathbf p,k+1}\\q_{n,\mathbf p,k+1}\end{pmatrix}
 =
 \begin{pmatrix}2&3\\1&2\end{pmatrix}
 \begin{pmatrix}a_{n,\mathbf p,k}\\q_{n,\mathbf p,k}\end{pmatrix}.
 \label{koide-param-eq:unit-recurrence}
\end{equation}

\begin{theorem}[Integral parametrization, with inverse and multiplicity]
\label{koide-param-thm:integral-parametrization}
The map
\begin{equation}
 (n,\mathbf p,k)
 \longmapsto
 (a_{n,\mathbf p,k},q_{n,\mathbf p,k},n),
 \qquad
 n\geq1,\quad\mathbf p\in\mathcal P(m_n),\quad k\geq0,
 \label{koide-param-eq:parameter-map}
\end{equation}
is a bijection onto all integral solutions of
\begin{equation}
 a^2-3q^2-18n^2=-2,\qquad a>0,\ q\geq0,\ n\geq1.
 \label{koide-param-eq:n-positive-solutions}
\end{equation}
Every such solution automatically satisfies $a+4n>q$.

The inverse is as follows.  From $(a,q,n)$ form
$\alpha=a+q\sqrt3$ and recover $(C,D,E)$ by
\eqref{koide-param-eq:forward-HNF-pre}--\eqref{koide-param-eq:forward-HNF}.  There is a unique
$k\geq0$ such that
\begin{equation}
 \alpha=\alpha_{n,(C,D,E)}\epsilon^k.
 \label{koide-param-eq:recover-k}
\end{equation}
Thus every point occurs with multiplicity exactly one.
\end{theorem}

\begin{proof}
Equation \eqref{koide-param-eq:unit-orbit-coordinates} preserves the norm, so its
image satisfies
\[
 a_{n,\mathbf p,k}^2-3q_{n,\mathbf p,k}^2=m_n,
\]
which is exactly the equation in \eqref{koide-param-eq:n-positive-solutions}.
Both embeddings of the generator are positive, hence $a>0$; and
$\alpha_{n,\mathbf p}\geq\sqrt{m_n}$ and $k\geq0$ imply $q\geq0$.

Conversely, a solution in \eqref{koide-param-eq:n-positive-solutions} has
\(\Norm(a+q\sqrt3)=m_n>0\).  Since $a>0$ and $q\geq0$, the inequality
$a^2-3q^2=m_n>0$ gives $a>\sqrt3q$; hence both embeddings are positive.
Its principal ideal has unique HNF data $(C,D,E)$.  The quotient of
 $a+q\sqrt3$ by the canonical generator of that ideal is a totally positive
 unit, hence is $\epsilon^k$ for a unique $k\in\ZZ$.  Lemma
 \ref{koide-param-lem:integral-seed-descent} supplies the completely integral recovery
 of $k$: iterate \eqref{koide-param-eq:inverse-descent} precisely while $a\leq2q$;
 the number of steps is $k$, and the terminal pair is
 $(a_{n,\mathbf p,0},q_{n,\mathbf p,0})$.  In particular $k\geq0$.  This
 proves surjectivity and the stated inverse.  HNF uniqueness, seed
 uniqueness, and unit-exponent uniqueness prove multiplicity one.

It remains to prove the chamber wall.  If a point with $n\geq1$ had
$q\geq a+4n$, then $a\leq q-4n$, and the unchanged equation would give
\[
 3q^2+18n^2-2=a^2\leq(q-4n)^2.
\]
After moving terms and dividing by $2$, this is
\[
 q^2+4qn+n^2\leq1.
\]
For nonnegative integers with $n\geq1$, equality could occur only at
$(q,n)=(0,1)$, but that pair does not satisfy $q\geq a+4n$ with $a>0$.
Thus the assumed wall inequality is impossible.
\end{proof}

The congruence needed by the original inverse map to the $(x,y,z)$
coordinates is retained automatically: reducing
$a^2-3q^2-18n^2=-2$ modulo $2$ gives
\begin{equation}
 a\equiv q\pmod2.
 \label{koide-param-eq:parity}
\end{equation}
No parity class has been lost or added in
Theorem~\ref{koide-param-thm:integral-parametrization}.

\begin{proposition}[The face $n=0$]
\label{koide-param-prop:n-zero}
All points satisfying the original equation and chamber with $n=0$ are
\begin{equation}
 a_k+q_k\sqrt3=(1+\sqrt3)\epsilon^k,\qquad k\geq1.
 \label{koide-param-eq:n-zero-family}
\end{equation}
They have height $4a_k$ and occur with multiplicity one.
\end{proposition}

\begin{proof}
When $n=0$, the norm is $-2$.  Its parity implies $a\equiv q\pmod2$, and
therefore
\[
 \frac{a+q\sqrt3}{1+\sqrt3}
 =\frac{3q-a}{2}+\frac{a-q}{2}\sqrt3\in\OO.
\]
The quotient has norm one, so it is $\pm\epsilon^k$ by
Lemma~\ref{koide-param-lem:ring-units}.  The conditions $a>0$ and $q\geq0$ select the
positive sign.  At $k=0$ one gets $(a,q)=(1,1)$, which is excluded by the
strict chamber wall $a>q$.  At every $k\geq1$, the recurrence
\eqref{koide-param-eq:unit-recurrence} gives
$a_{k+1}-q_{k+1}=a_k+q_k>0$.  For $k\leq-1$, direct use of
$a=(t-2/t)/2$ with $t=(1+\sqrt3)\epsilon^k\leq\sqrt3-1<\sqrt2$
gives $a<0$.  Hence precisely the exponents $k\geq1$ occur.
\end{proof}

\section{Exact divisor and sharp-height formulas}

The number of unit-orbits at a fixed positive $n$ has an especially simple
divisor formula.

\begin{proposition}[Orbit divisor sum and shifted factorization]
\label{koide-param-prop:orbit-divisor}
For $n\geq1$, the number of orbits under
\((a,q)\mapsto(2a+3q,a+2q)\) among the solutions with $a>0$ and
$q\geq0$ is
\begin{equation}
 \rho(m_n)=\sum_{d\mid18n^2-2}\chi_{12}(d)
 =\rho(3n-1)\rho(3n+1).
 \label{koide-param-eq:orbit-divisor-sum}
\end{equation}
\end{proposition}

\begin{proof}
The first equality is Theorem~\ref{koide-param-thm:integral-parametrization} and
Proposition~\ref{koide-param-prop:HNF-ideals}.  For the second, write
\[
 m_n=2(3n-1)(3n+1).
\]
The two displayed linear factors have greatest common divisor dividing
$2$.  The function $\rho=1*\chi_{12}$ is multiplicative and
$\rho(2^j)=1$ for all $j\geq0$.  Removing all powers of $2$ therefore
makes the two factors coprime without changing any of the three $\rho$
values, proving the product identity.
\end{proof}

For $\mathbf p\in\mathcal P(m_n)$ define its exact unit phase by
\begin{equation}
 \tau_{n,\mathbf p}
 =\frac{\log(\alpha_{n,\mathbf p}/\sqrt{m_n})}{\log\epsilon},
 \qquad0\leq\tau_{n,\mathbf p}<1.
 \label{koide-param-eq:unit-phase}
\end{equation}
For $n\geq1$ put
\begin{equation}
 Y_n(X)=\frac{X-17n}{4}.
 \label{koide-param-eq:Y-n}
\end{equation}
Whenever $Y_n(X)\geq\sqrt{m_n}$, put
\begin{equation}
 L_n(X)=
 \frac1{\log\epsilon}
 \log\!\left(
 \frac{Y_n(X)+\sqrt{Y_n(X)^2-m_n}}{\sqrt{m_n}}
 \right).
 \label{koide-param-eq:L-n}
\end{equation}
Let
\begin{equation}
 \mathcal N(X)=
 \{n\in\ZZ_{\geq1}:Y_n(X)\geq\sqrt{m_n}\}.
 \label{koide-param-eq:N-X}
\end{equation}
This is the exact finite range
\begin{equation}
 X\geq17n+4\sqrt{18n^2-2};
 \label{koide-param-eq:n-boundary}
\end{equation}
no asymptotic replacement of this boundary is made.

For the $n=0$ face define
\begin{equation}
 R_0(X)=
 \max\!\left\{0,
 \left\lfloor
 \frac{
 \log\!\left(
 \dfrac{X/4+\sqrt{(X/4)^2+2}}{1+\sqrt3}
 \right)}{\log\epsilon}
 \right\rfloor\right\}.
 \label{koide-param-eq:R0-exact}
\end{equation}

\begin{theorem}[Exact sharp-height floor formula]
\label{koide-param-thm:exact-floor}
For every real $X\geq0$,
\begin{equation}
 R(X)=R_0(X)+
 \sum_{n\in\mathcal N(X)}
 \sum_{\mathbf p\in\mathcal P(m_n)}
 \left(\left\lfloor L_n(X)-\tau_{n,\mathbf p}\right\rfloor+1\right).
 \label{koide-param-eq:exact-floor-formula}
\end{equation}
Every equality case in $4a+17n\leq X$ is included by the floor in
\eqref{koide-param-eq:exact-floor-formula}.
\end{theorem}

\begin{proof}
For a fixed positive $n$ and parameter $\mathbf p$, put
$t_k=\alpha_{n,\mathbf p}\epsilon^k$.  Since
$t_k\geq\sqrt{m_n}$,
\begin{equation}
 a_{n,\mathbf p,k}=\frac12\left(t_k+\frac{m_n}{t_k}\right)
 \label{koide-param-eq:a-from-t}
\end{equation}
is nondecreasing, and is strictly increasing after the possible initial
case $q=0$.  The unchanged height inequality is
$a_{n,\mathbf p,k}\leq Y_n(X)$.  It has no solution unless
$Y_n(X)\geq\sqrt{m_n}$.  In that range, solving the quadratic inequality
in $t_k$ gives
\begin{equation}
 t_k\leq Y_n(X)+\sqrt{Y_n(X)^2-m_n}.
 \label{koide-param-eq:t-height-bound}
\end{equation}
After division by $\sqrt{m_n}$ and logarithms, this is precisely
\begin{equation}
 k+\tau_{n,\mathbf p}\leq L_n(X).
 \label{koide-param-eq:k-phase-bound}
\end{equation}
Because $L_n(X)\geq0$ and $0\leq\tau_{n,\mathbf p}<1$, the number of
integers $k\geq0$ satisfying \eqref{koide-param-eq:k-phase-bound} is exactly
$\lfloor L_n(X)-\tau_{n,\mathbf p}\rfloor+1$; this expression is zero,
not negative, when $L_n(X)<\tau_{n,\mathbf p}$.  Equality is retained.
Theorem~\ref{koide-param-thm:integral-parametrization} now proves the positive-$n$
sum.

For $n=0$, put $t=(1+\sqrt3)\epsilon^k$.  Since the norm is $-2$,
$a_k=(t-2/t)/2$.  The inequality $4a_k\leq X$ is equivalent to
\[
 t\leq X/4+\sqrt{(X/4)^2+2}.
\]
Counting the integers $k\geq1$ gives \eqref{koide-param-eq:R0-exact} and completes
the proof.
\end{proof}

Write
\begin{equation}
 J_n(X)=\lfloor L_n(X)\rfloor,
 \qquad
 \lambda_n(X)=L_n(X)-J_n(X)\in[0,1),
 \label{koide-param-eq:J-lambda}
\end{equation}
and define the exact phase distribution
\begin{equation}
 A_n(t)=\#\{\mathbf p\in\mathcal P(m_n):
                     \tau_{n,\mathbf p}\leq t\},
 \qquad0\leq t<1,
 \label{koide-param-eq:A-n}
\end{equation}
with discrepancy
\begin{equation}
 \Delta_n(t)=A_n(t)-t\rho(m_n).
 \label{koide-param-eq:Delta-n}
\end{equation}
Since
\(\lfloor J+\lambda-\tau\rfloor+1
=J+\mathbf1_{\{\tau\leq\lambda\}}\), Theorem~\ref{koide-param-thm:exact-floor}
gives the exact identity
\begin{align}
 R(X)-C_0X
 &=R_0(X)
 +\left[
   \sum_{n\in\mathcal N(X)}
     \rho(m_n)L_n(X)-C_0X
  \right]
 +\sum_{n\in\mathcal N(X)}
    \Delta_n(\lambda_n(X)).
 \label{koide-param-eq:exact-discrepancy-splitting}
\end{align}
This identity includes the two unbounded coordinate faces and every sharp
height boundary point.  Formula \eqref{koide-param-eq:R0-exact} also gives
$R_0(X)=O(\log(2+X))$, so the $n=0$ face is below the requested
$X^{1/2+\varepsilon}$ scale; the two positive-$n$ brackets are the genuine
analytic terms.

\section{The phase term as exact Hecke-character data}

The last sum in \eqref{koide-param-eq:exact-discrepancy-splitting} is not determined
by the divisor count $\rho(m_n)$.  It records where the generator of each
ideal lies inside a fundamental interval for the positive units.

For $\ell\in\ZZ$ and a nonzero integral ideal $\mathfrak a=(\gamma)$,
define
\begin{equation}
 \xi_\ell(\mathfrak a)=
 \exp\!\left(
  \frac{\pi i\ell}{\log\epsilon}
  \log\left|\frac{\gamma}{\gamma'}\right|
 \right).
 \label{koide-param-eq:Hecke-character}
\end{equation}
This is independent of the generator: replacing $\gamma$ by
$\pm\epsilon^j\gamma$ changes the logarithm by $2j\log\epsilon$, and
the exponential by $e^{2\pi i\ell j}=1$.  It is multiplicative, hence is
the unramified ideal part of a unitary Hecke character of $K$.  For
$m=m_n$ and the canonical generator,
\begin{equation}
 \xi_\ell(\mathfrak a(C,D,E))
 =e^{2\pi i\ell\tau_{n,(C,D,E)}}.
 \label{koide-param-eq:character-phase}
\end{equation}
Define its norm-$m$ coefficient by
\begin{equation}
 S_\ell(m)=\sum_{\Norm\mathfrak a=m}\xi_\ell(\mathfrak a).
 \label{koide-param-eq:S-ell}
\end{equation}
Then $S_0(m)=\rho(m)$, whereas the nonzero $S_\ell(m_n)$ are the Fourier
coefficients of the phase multiset.

The coefficient $S_\ell(m)$ also has a completely explicit Euler-factor
description.  It is multiplicative in $m$.  If $p\geq5$ is split, choose
one prime ideal $\mathfrak p$ above $p$ and put
$\omega_{\ell,p}=\xi_\ell(\mathfrak p)$.  Since
$\mathfrak p\overline{\mathfrak p}=(p)$ and
$\xi_\ell((p))=1$, one has
\begin{equation}
 S_\ell(p^e)=\sum_{j=0}^{e}\omega_{\ell,p}^{2j-e}
 \qquad(\chi_{12}(p)=1).
 \label{koide-param-eq:split-Hecke-factor}
\end{equation}
If $p\geq5$ is inert, then
\begin{equation}
 S_\ell(p^e)=
 \begin{cases}1,&2\mid e,\\0,&2\nmid e.
 \end{cases}
 \qquad(\chi_{12}(p)=-1).
 \label{koide-param-eq:inert-Hecke-factor}
\end{equation}
Finally the unique primes above $2$ and $3$ are generated by
$1+\sqrt3$ and $\sqrt3$, respectively.  Since
\[
 \left|\frac{1+\sqrt3}{1-\sqrt3}\right|=\epsilon,
 \qquad
 \left|\frac{\sqrt3}{-\sqrt3}\right|=1,
\]
their factors are
\begin{equation}
 S_\ell(2^e)=(-1)^{\ell e},
 \qquad S_\ell(3^e)=1.
 \label{koide-param-eq:ramified-Hecke-factors}
\end{equation}
These formulas enumerate every ideal factor with multiplicity one.  At
$\ell=0$ they reduce prime by prime to the divisor sum $\rho(m)$; at
$\ell\ne0$ they show exactly what angular information the unweighted
divisor sum omits.

They also retain the shifted factorization of $m_n$ at every Fourier mode.
Writing $3n-1=2^r u$ and $3n+1=2^s v$ with $u,v$ odd, one has
$\gcd(u,v)=1$.  Multiplicativity and
\eqref{koide-param-eq:ramified-Hecke-factors} therefore give the exact identity
\begin{equation}
 S_\ell(18n^2-2)
 =(-1)^\ell S_\ell(3n-1)S_\ell(3n+1).
 \label{koide-param-eq:shifted-Hecke-factorization}
\end{equation}
Indeed, the exponent of $2$ on the left is $1+r+s$, while the product on
the right has total exponent $r+s$; all odd prime factors are disjoint.
Thus the angular obstruction is itself a shifted correlation, now of
nontrivial Hecke coefficients rather than merely of their dimensions.

Here is an endpoint-exact Fourier statement.  Suppose $0<\lambda<1$, and
let
\begin{equation}
 b_n(t)=\#\{\mathbf p\in\mathcal P(m_n):\tau_{n,\mathbf p}=t\}.
 \label{koide-param-eq:b-n}
\end{equation}
Fej\'er summation of the Fourier series of
$\mathbf1_{[0,\lambda]}-\lambda$ gives
\begin{align}
 \Delta_n(\lambda)
 &=\frac12\bigl(b_n(0)+b_n(\lambda)\bigr)\notag\\
 &\quad+
 \lim_{M\to\infty}
 \sum_{0<|\ell|\leq M}
 \left(1-\frac{|\ell|}{M+1}\right)
 \frac{1-e^{-2\pi i\ell\lambda}}{2\pi i\ell}
 S_\ell(m_n).
 \label{koide-param-eq:Fejer-exact}
\end{align}
The two half-weights correct the midpoint convention at the circular
endpoint $0$ and at $\lambda$.  When $\lambda=0$, the exact value is
simply $\Delta_n(0)=b_n(0)$.

Combining this with \eqref{koide-param-eq:shifted-Hecke-factorization} makes the
aggregate obstruction a literal shifted correlation.  Define the exact
endpoint contribution
\begin{equation}
 \mathcal E(X)=
 \sum_{\substack{n\in\mathcal N(X)\\\lambda_n(X)=0}}b_n(0)
 +\frac12
 \sum_{\substack{n\in\mathcal N(X)\\\lambda_n(X)>0}}
       \bigl(b_n(0)+b_n(\lambda_n(X))\bigr).
 \label{koide-param-eq:endpoint-aggregate}
\end{equation}
Since $\mathcal N(X)$ is finite, summation over $n$ commutes with the
Fej\'er limit and gives
\begin{align}
 \sum_{n\in\mathcal N(X)}\Delta_n(\lambda_n(X))
 &=\mathcal E(X)\notag\\
 &\quad+\lim_{M\to\infty}
 \sum_{0<|\ell|\leq M}
 \left(1-\frac{|\ell|}{M+1}\right)
 \frac{(-1)^\ell}{2\pi i\ell}\notag\\
 &\qquad\times
 \sum_{\substack{n\in\mathcal N(X)\\\lambda_n(X)>0}}
 \bigl(1-e^{-2\pi i\ell\lambda_n(X)}\bigr)
 S_\ell(3n-1)S_\ell(3n+1).
 \label{koide-param-eq:exact-shifted-Hecke-correlation}
\end{align}
Every coefficient and endpoint in this identity comes from the original
weak height inequality; no smoothing has been inserted.

Equations \eqref{koide-param-eq:exact-discrepancy-splitting} and
\eqref{koide-param-eq:Fejer-exact} identify the analytic problem without discarding
the original height.  Its radial part is the sharply weighted shifted
correlation
\begin{equation}
 \sum_{n\in\mathcal N(X)}
 \rho(3n-1)\rho(3n+1)L_n(X),
 \label{koide-param-eq:radial-shifted-correlation}
\end{equation}
and its angular part requires cancellation in the polynomially sampled
Hecke coefficients
\begin{equation}
 S_\ell(18n^2-2),
 \label{koide-param-eq:polynomial-Hecke-samples}
\end{equation}
uniformly over the Fourier modes needed to approximate the discontinuous
threshold $\tau\leq\lambda_n(X)$.  The endpoint
\eqref{koide-param-eq:n-boundary} gives $L_n(X)=0$ and has a square-root singularity
in its derivative; it is not a fixed smooth weight.  A bound for the
unweighted orbit total $\rho(m_n)$ therefore does not bound the actual
sharp-height count.

\section{An unconditional effective exponent for the exact sector}

This section records the power saving that is presently justified without
assuming the square-root estimate.

Let
\begin{equation}
 F(a,q,n)=a^2-3q^2-18n^2,
 \qquad
 \mathscr X=\{F=-2\},
 \qquad
 \mathbf G=\operatorname{Spin}(F).
 \label{koide-param-eq:group-variety}
\end{equation}
Let $x_0=(1,1,0)$.  Its orthogonal complement is
\begin{equation}
 x_0^\perp=\{(3s,s,t):s,t\in\QQ\},
 \qquad
 F(3s,s,t)=6(s^2-3t^2),
 \label{koide-param-eq:stabilizer-complement}
\end{equation}
so the identity component of its stabilizer is the norm-one torus of
$\QQ(\sqrt3)$.  Its rational character group is zero because Galois
conjugation acts by $-1$ on its geometric character lattice.

The form $F$ is anisotropic over $\QQ$.  Indeed, suppose a primitive
integral triple satisfies $F(a,q,n)=0$.  Modulo $3$ one first gets
$3\mid a$.  Writing $a=3a_1$ and dividing the equation by $3$ gives
$3a_1^2-q^2-6n^2=0$, hence $3\mid q$.  Write $q=3q_1$ and divide by
$3$ again to get
\begin{equation}
 a_1^2-3q_1^2-2n^2=0.
\end{equation}
Modulo $3$, the sum $a_1^2+n^2$ is zero, which forces
$3\mid a_1$ and $3\mid n$.  Thus $3$ divides $a,q,n$, contradicting
primitivity.  In particular the associated arithmetic Fuchsian quotient
is cocompact.

For $X>0$ let $B_X$ be the exact subset of $\mathscr X(\RR)$ cut out by
\begin{equation}
 a>0,\quad q\geq0,\quad n\geq0,\quad a+4n>q,
 \quad4a+17n\leq X.
 \label{koide-param-eq:B-X}
\end{equation}
With the gauge measure
\begin{equation}
 dm_\infty=\frac{dq\,dn}{2a}\qquad(a>0),
 \label{koide-param-eq:gauge-measure}
\end{equation}
put
\begin{equation}
 I_\infty=\frac1{2\sqrt3}
 \log\!\left(\frac3{2\sqrt2}\right).
 \label{koide-param-eq:I-infty}
\end{equation}

\begin{lemma}[Exact-sector volume with a power-saving remainder]
\label{koide-param-lem:volume}
As $X\to\infty$,
\begin{equation}
 m_\infty(B_X)=I_\infty X+O(1).
 \label{koide-param-eq:volume-linear-O1}
\end{equation}
Moreover, the family $B_X$ is effectively well-rounded in the sense of
Benoist--Oh, Definition~13.1~\cite[p.~1931]{BenoistOh_2012}.
\end{lemma}

\begin{proof}
Introduce exact polar coordinates on the negative plane by
\begin{equation}
 \sqrt3q=r\cos\theta,\qquad
 3\sqrt2n=r\sin\theta,\qquad
 0\leq\theta\leq\pi/2.
 \label{koide-param-eq:polar}
\end{equation}
On $F=-2$ one has $r=\sqrt{a^2+2}$ and
\begin{equation}
 dm_\infty=\frac1{6\sqrt6}\,da\,d\theta.
 \label{koide-param-eq:polar-measure}
\end{equation}
Write $b(\theta)=17\sin\theta/(3\sqrt2)$.  In the positive orthant the
outer boundary is the unique nonnegative solution $A_X(\theta)$ of
\begin{equation}
 4A_X(\theta)+b(\theta)\sqrt{A_X(\theta)^2+2}=X.
 \label{koide-param-eq:A-boundary}
\end{equation}
For all sufficiently large $X$, uniformly in $\theta$,
\begin{equation}
 A_X(\theta)=\frac{X}{4+b(\theta)}+O(X^{-1}).
 \label{koide-param-eq:A-asymptotic}
\end{equation}
To see this without a formal expansion, subtract
$(4+b)A_X$ from \eqref{koide-param-eq:A-boundary} and use
\[
 \sqrt{A_X^2+2}-A_X
 =\frac2{\sqrt{A_X^2+2}+A_X}=O(X^{-1}),
\]
where \eqref{koide-param-eq:A-boundary} uniformly gives $A_X\asymp X$.
Integration of \eqref{koide-param-eq:A-asymptotic} in
\eqref{koide-param-eq:polar-measure} gives
\[
 \frac{X}{6\sqrt6}\int_0^{\pi/2}
 \frac{d\theta}{4+\frac{17}{3\sqrt2}\sin\theta}+O(X^{-1})
 =I_\infty X+O(X^{-1}).
\]
For the asserted exact value, set
$c=17/(3\sqrt2)$, $d=\sqrt{c^2-16}=1/(3\sqrt2)$, and
$t=\tan(\theta/2)$.  Then
\begin{align*}
 \int_0^{\pi/2}\frac{d\theta}{4+c\sin\theta}
 &=\int_0^1\frac{dt}{2t^2+ct+2}\\
 &=\frac1d
   \left[\log\!\left(\frac{4t+c-d}{4t+c+d}\right)\right]_{t=0}^{t=1}\\
 &=3\sqrt2\log\!\left(\frac3{2\sqrt2}\right).
\end{align*}
Multiplication by $1/(6\sqrt6)$ is exactly the constant in
\eqref{koide-param-eq:I-infty}.
If the strict chamber wall fails, then the calculation in the proof of
Theorem~\ref{koide-param-thm:integral-parametrization}, now over $\RR$, gives
$q^2+4qn+n^2\leq1$.  The removed real subset is therefore compact.
For all large $X$ it contributes a fixed finite measure, proving
\eqref{koide-param-eq:volume-linear-O1}.

It remains to check effective well-roundedness for the actual sector.
Let $U_\delta$ be the image in the fixed vector representation of the
group ball used by Benoist--Oh~\cite[Definition~13.1]{BenoistOh_2012}.
For every linear functional $\ell$ on
$\RR^3$ there is a constant, independent of $x$ and $\delta$, such that
\begin{equation}
 |\ell(xg)-\ell(x)|\leq C_\ell\delta\|x\|
 \qquad(g\in U_\delta).
 \label{koide-param-eq:linear-functional-motion}
\end{equation}
Apply this to the five literal boundary functionals
\[
 a,\qquad q,\qquad n,\qquad a+4n-q,\qquad4a+17n.
\]
On the unbounded part of the positive sheet, $\|x\|\asymp a+1$, and the
last functional is also comparable to $\|x\|$.  Thus membership in the
outer height inequality can change only in a shell of radial width
$O(\delta X)$.  Membership in $q\geq0$ or $n\geq0$ can change only in
angular strips of width $O(\delta)$ about $\theta=\pi/2$ or $\theta=0$.
The boundary $a=0$ is compact.  The boundary $a+4n-q=0$ is compact by
the calculation preceding \eqref{koide-param-eq:volume-linear-O1}; moreover its
restriction to $\mathscr X$ is regular, because proportionality of
$(1,-1,4)$ to $(2a,-6q,-36n)$ is incompatible with $a,q,n\geq0$.
Consequently its group-action boundary tube has measure $O(\delta)$.
By \eqref{koide-param-eq:polar-measure}, all these regions together have measure
$O(\delta X)$.  The same analysis applies to the inner intersection that
defines $B^-_{X,\delta}$.  Hence
\begin{equation}
 m_\infty(B^+_{X,\delta}\setminus B^-_{X,\delta})
 \ll \delta\,m_\infty(B_X),
 \label{koide-param-eq:BO-boundary-condition}
\end{equation}
which is precisely Definition 13.1(2), with boundary exponent $1$, for
the group-action thickenings rather than merely Euclidean tubes.

For Definition~13.1(3) of \cite{BenoistOh_2012}, take the archimedean height to be the
Euclidean norm in the displayed integral basis.  It is bounded below on
$\mathscr X$.  Uniformly for $0<\delta<1$, the thickening
$B^+_{X,\delta}$ lies in a fixed finite enlargement of the polar sector;
there $H_\infty(x)\asymp1+a$ and the measure in a radial interval is
$O(da)$.  Hence for every $s>0$, uniformly in $\delta$,
\[
 \int_{B_X^+}\|u\|^{-s}\,dm_\infty(u)
 \ll
 \begin{cases}
 X^{1-s},&0<s<1,\\
 \log X,&s=1,\\
 1,&s>1.
 \end{cases}
\]
Each line is $O(m_\infty(B_X)^{1-\delta_s})$, for example with
$\delta_s=\min(s/2,1/2)>0$.  This is the remaining condition.  Since
$S=\{\infty\}$, Definition 13.1(1), concerning a finite-place compact
open subgroup, is void.  The choices of strict
or weak inequality on the finitely many piecewise smooth boundary arcs do
not change these neighborhood estimates.
\end{proof}

\begin{theorem}[Unconditional effective conclusion]
\label{koide-param-thm:power-saving}
There exists $\eta>0$, depending only on the fixed form, lattice, chamber,
and height in \eqref{koide-param-eq:original-problem}, such that
\begin{equation}
 R(X)=C_0X+O(X^{1-\eta}).
 \label{koide-param-eq:power-saving-again}
\end{equation}
\end{theorem}

\begin{proof}
Use the standard vector representation
$\rho:\operatorname{Spin}(F)\to\operatorname{SO}(F)\subset
\operatorname{GL}_3$, whose kernel is finite, and the displayed integral
basis of the exact lattice $\ZZ^3$.  Let
\[
 \mathbf G_{\ZZ}
 =\{g\in\mathbf G(\QQ):\rho(g)\ZZ^3=\ZZ^3\}.
\]
This is the integral group attached to the chosen representation and
lattice in the Benoist--Oh setup~\cite[pp.~1892--1893]{BenoistOh_2012};
take their congruence subgroup
$\Gamma_S$ to be the level-one group $\mathbf G_{\ZZ}$ itself.  The group
$\mathbf G$ is connected,
simply connected, and almost $\QQ$-simple.  It is noncompact over $\RR$
because $F$ has signature $(1,2)$.  Thus $S=\{\infty\}$ contains every
noncompact archimedean place and $G_S$ is noncompact.  The orbit
$x_0\mathbf G$ is the affine symmetric variety $\mathscr X$, and its
stabilizer has no nontrivial rational character by
\eqref{koide-param-eq:stabilizer-complement}.  These verify every standing hypothesis
on pp.~1892--1893 of \cite{BenoistOh_2012}.  Their Theorem~12.5 supplies the required
effective equidistribution of the symmetric-stabilizer translates.

The Borel--Harish-Chandra finiteness theorem quoted in \cite[p.~1892]{BenoistOh_2012}
p.~1892 says that $\mathscr X(\ZZ)$ is a finite union of
$\mathbf G_{\ZZ}$-orbits.  A rational representative of any such orbit
again has a symmetric stabilizer with no rational
character: its orthogonal complement is anisotropic over $\QQ$, since an
isotropic vector in it would be an isotropic vector for $F$.  Apply
Benoist--Oh, Theorem~13.2~\cite[pp.~1931--1932]{BenoistOh_2012}, to each orbit and to the exact effectively
well-rounded family $B_X$ of Lemma~\ref{koide-param-lem:volume}.  For each orbit it
gives its invariant-volume main term with error
$O(m_\infty(B_X)^{1-\eta_j})$ for some $\eta_j>0$.  There are finitely
many orbits, so their sum has the same form with
$\eta=\min_j\eta_j>0$.

All invariant measures on the same real homogeneous variety are
proportional.  Hence the sum of the finitely many orbit main terms is
$K\,m_\infty(B_X)$ for one fixed $K>0$.  Lemma~\ref{koide-param-lem:volume} and the
already established asymptotic $R(X)\sim C_0X$ force
$K I_\infty=C_0$.  Therefore
\[
 K\,m_\infty(B_X)=K(I_\infty X+O(1))=C_0X+O(1).
\]
Since $m_\infty(B_X)\asymp X$, summing the effective errors proves
\eqref{koide-param-eq:power-saving-again}; the $O(1)$ term is absorbed after decreasing
$\eta$ if necessary.  The theorem is applied first to the sequence of
positive integral $X$.  For real $X$, the original height $4a+17n$ is an
integer, so $R(X)=R(\lfloor X\rfloor)$; replacing $X$ by $\lfloor X\rfloor$
changes $C_0X$ by $O(1)$ and proves the same estimate.
\end{proof}

\begin{remark}[What exponent has actually been obtained]
Theorem~\ref{koide-param-thm:power-saving} is an unconditional exponent strictly
smaller than $1$: the error is $O(X^{1-\eta})$ for some $\eta>0$.
Neither the statement nor the proof of \cite[Theorem~13.2]{BenoistOh_2012} identifies an exponent
$1-\eta\leq1/2$ for this sharp sector.  Tracing their Sobolev and mixing
constants would make a positive $\eta$ effective in principle, but that
calculation has not been carried out in this section and cannot be replaced
by the assertion $\eta=1/2$.
\end{remark}

\section{Exact obstruction to the square-root claim}

There are two complementary descriptions of the remaining obstruction;
the first is an exact identity, while the second explains the relevant
spectral mechanisms without asserting an unproved intertwining formula.

First, the elementary ideal description gives the literal identity
\eqref{koide-param-eq:exact-discrepancy-splitting}.  Thus the square-root claim is
equivalent to a bound of that size for the displayed sum of the radial
bracket, the aggregate phase discrepancy, and $R_0(X)$.  Controlling the
radial and phase terms separately is a sufficient route, though the exact
identity also permits cancellation between them.  Formula
\eqref{koide-param-eq:Fejer-exact} shows that a separate estimate of the latter entails
uniform cancellation in the nonzero archimedean Hecke modes
\eqref{koide-param-eq:polynomial-Hecke-samples}.  Merely using
\(\#\mathcal P(m_n)=\rho(m_n)\), or replacing each phase distribution by
Lebesgue measure, deletes a term that is exactly present in $R(X)$.

Second, $F$ is rationally anisotropic, so $\mathbf G$ gives a cocompact
arithmetic Fuchsian quotient.  A negative-norm vector represents an
oriented geodesic, and its stabilizer is the split real torus appearing in
\eqref{koide-param-eq:stabilizer-complement}.  In the standard radial torus-orbit
counting kernel, the spectral expansion has two possible sources larger
than the square-root scale:
\begin{enumerate}
 \item an exceptional eigenvalue $\lambda=s(1-s)<1/4$ can contribute
 $c_sX^s$ with $s>1/2$, unless the exact torus period and exact sector
 transform make $c_s=0$;
 \item after exceptional terms are absent or removed, estimating the
 tempered spectral sum for a discontinuous cutoff at the conjectural
 $X^{1/2+\e}$ scale requires cancellation beyond the general
 effective-mixing argument.
\end{enumerate}
For the present nonradial chamber, one must first construct a
remainder-preserving sector spectral expansion before these candidate
terms can even be asserted with their exact coefficients.  An actual
refutation by a secondary term would then require proving both $s>1/2$
and $c_s\ne0$ for this fixed arithmetic quotient and exact chamber.  No
such construction or nonvanishing has been established here.  Conversely,
absence of exceptional terms alone would not estimate the full tempered
sharp-cutoff sum.  This is why the exact parametrization neither proves
nor refutes \eqref{koide-param-eq:square-root-target}, while it does isolate the data a
future proof must bound.

For comparison only, not as a theorem about the present nonradial sector,
the classical hyperbolic lattice-point problem for a conjugacy class has
an unconditional $O(X^{2/3})$ remainder due to Good
\cite[Theorem~4]{Good_1983}, with the modern proof of
Chatzakos--Petridis~\cite[Theorem~1.1 and Conjecture~5.7]{ChatzakosPetridis_2016},
and a conjectural $O_\e(X^{1/2+\e})$ scale.  That
published $O(X^{2/3})$ theorem counts a different radial kernel.  An exact
intertwining of its kernel with the simultaneous chamber and height in
\eqref{koide-param-eq:original-problem} has not been proved, so $O(X^{2/3})$ is not
claimed here for $R(X)$.

The independent finite replay used to check the HNF multiplicities, seed
reduction, inverse descent, and divisor factorization for \(1\leq n\leq200\)
is retained at
\path{supporting_materials/workbench/evidence/koide_exact_checkpoint_20260904/parametrization/verify_parametrization.py};
its SHA--256 digest is
\texttt{061c5e22258628169fd1ecb8bcc5e7cf\allowbreak%
8b8d19bfdb918650668c3e90f2f9730d}.
That replay is a finite consistency check.  The proofs above, rather than the
finite range, establish the parametrization for every admissible integer.

\section{Boundary ledger}

For clarity, every boundary used above has the following exact treatment.
\begin{itemize}
 \item $a>0$ is strict.  There is no omitted integral solution with
 $a=0$, because the original equation would give a multiple of $3$ equal
 to $2$.
 \item $q=0$ is included.  It corresponds to phase $\tau=0$, is included
 by $\tau\leq\lambda$, and is retained in the endpoint term
 $b_n(0)/2$ in \eqref{koide-param-eq:Fejer-exact}.
 \item $n=0$ is included and is counted exactly by $R_0(X)$; the sole
 solution $(a,q,n)=(1,1,0)$ on $a+4n=q$ is excluded.
 \item $a+4n=q$ is strict.  Its real intersection is bounded and its sole
 integral point in the nonnegative chamber is $(1,1,0)$.
 \item $4a+17n=X$ is included.  Equality is included in
 \eqref{koide-param-eq:t-height-bound}, \eqref{koide-param-eq:k-phase-bound}, and every floor in
 \eqref{koide-param-eq:exact-floor-formula}.
 \item The canonical unit interval is half-open,
 $[\sqrt{m_n},\epsilon\sqrt{m_n})$.  Consequently no generator is counted
 twice at a unit boundary.
 \item Conjugate ideals are distinct unless the ideal is fixed by
 conjugation.  They yield distinct canonical phases and distinct
 $q\geq0$ unit-orbits; no division by $2$ is made.  When $q=0$, the ideal
 is self-conjugate and still occurs once by HNF uniqueness.
\end{itemize}

% Integrated 4 September 2026 from the exact 10^8 census and independent replay.
% This section retains the literal ternary equation, chamber, strict wall, height,
% and original-coordinate inverse.

\section{Exact unit-orbit census through height \(10^8\)}
\label{koide-count-sec-census-1e8}

The census retains exactly the same equation, chamber, strict wall, and
height as the preceding theorems:
\begin{equation}
 (a,q,n)\in\mathbb Z^3,\qquad
 a^2-3q^2-18n^2=-2,\qquad
 a>0,\quad q,n\geq0,\quad a+4n>q,\qquad
 4a+17n\leq X.
 \label{koide-count-eq-census-domain}
\end{equation}
It does not infer an asymptotic from finite data.  Its exact fixed-\(n\)
equation and factorization are
\begin{equation}
 a^2-3q^2=m_n,\qquad
 m_n=18n^2-2=2(3n-1)(3n+1).
 \label{koide-count-eq-census-fixed-n}
\end{equation}
For each \(n\geq1\), every ideal of norm \(m_n\) is enumerated by its
Hermite-normal-form triple
\begin{equation}
 (C,D,E)\in\mathbb Z_{>0}^2\times\mathbb Z,\qquad
 C^2D=m_n,\qquad
 0\leq E<D,\qquad E^2\equiv3\pmod D.
 \label{koide-count-eq-census-hnf}
\end{equation}
The generator is reduced by exact multiplication by the fundamental unit
\(2+\sqrt3\) to the unique seed \(a_0>2q_0\geq0\).  Its forward unit
orbit is
\begin{equation}
 \binom{a_{k+1}}{q_{k+1}}
 =
 \begin{pmatrix}2&3\\1&2\end{pmatrix}
 \binom{a_k}{q_k},\qquad k\geq0.
 \label{koide-count-eq-census-unit-step}
\end{equation}
Each emitted point is reconstructed by
\begin{equation}
 x=\frac{a+4n-q}{2},\qquad
 y=\frac{a+4n+q}{2},\qquad
 z=4a+17n,
 \label{koide-count-eq-census-inverse}
\end{equation}
and is rechecked in both coordinate systems and in the original polynomial
\(\Delta(x,y,z)=-1\).  Thus the accelerated enumeration does not replace
the original objects by a floating-point or rescaled surrogate.

\subsection{Complete result and checkpoint table}

The complete integer-domain audit for
\(1\leq X\leq100{,}000{,}000\) gives
\begin{equation}
 \boxed{R(100{,}000{,}000)=3{,}536{,}002.}
 \label{koide-count-eq-census-final-count}
\end{equation}
The endpoint in \(n\) is \(2{,}943{,}725\); the enumeration contains
\(2{,}327{,}513\) canonical unit orbits.  The following checkpoints are
exact integer counts.  The residuals use the proved constant
\(C_0\) in \eqref{koide-count-eq-final-constant}.
\begin{center}
\begin{tabular}{r@{\qquad}r@{\qquad}r@{\qquad}r}
\toprule
\(X\)&\(R(X)\)&\(R(X)-C_0X\)&\((R(X)-C_0X)/\sqrt X\)\\
\midrule
200,000&7,109&\(36.590440053629\)&\(0.081818711287\)\\
250,000&8,854&\(13.488050067037\)&\(0.026976100134\)\\
500,000&17,695&\(13.976100134073\)&\(0.019765190359\)\\
1,000,000&35,381&\(18.952200268147\)&\(0.018952200268\)\\
2,000,000&70,637&\(-87.095599463706\)&\(-0.061585888992\)\\
5,000,000&176,943&\(132.761001340734\)&\(0.059372524752\)\\
10,000,000&353,598&\(-22.477997318532\)&\(-0.007108166877\)\\
20,000,000&707,458&\(217.044005362937\)&\(0.048532515010\)\\
50,000,000&1,768,230&\(127.610013407342\)&\(0.018046781166\)\\
100,000,000&3,536,002&\(-202.779973185317\)&\(-0.020277997319\)\\
\bottomrule
\end{tabular}
\end{center}

\subsection{Certified extrema over every integer}

Write \(E(X)=R(X)-C_0X\).  Between consecutive occupied heights,
\(R(X)\) is constant and both \(E(X)\) and
\(E(X)/\sqrt X\) are strictly decreasing, since
\begin{equation}
 \frac{d}{dX}\left(\frac{R}{\sqrt X}-C_0\sqrt X\right)
 =-\frac{R}{2X^{3/2}}-\frac{C_0}{2\sqrt X}<0.
 \label{koide-count-eq-census-monotonicity}
\end{equation}
Checking both integer endpoints of every constant-count interval therefore
certifies extrema over every integer in the audited range, rather than at
selected checkpoints.  On \(1\leq X\leq10^8\), the maximum signed residual
is
\begin{equation}
 \begin{aligned}
 E(77{,}078{,}404)
 &=721.7934249571297569656681553538691093379\ldots,\\
 R(77{,}078{,}404)&=2{,}726{,}372,
 \end{aligned}
 \label{koide-count-eq-census-max-positive}
\end{equation}
and the minimum signed, hence largest absolute, residual is
\begin{equation}
 \begin{aligned}
 E(32{,}461{,}294)
 &=-812.8300691488068897091966006797491655933\ldots,\\
 R(32{,}461{,}294)&=1{,}147{,}085.
 \end{aligned}
 \label{koide-count-eq-census-max-negative}
\end{equation}
The maximum of
\(\lvert E(X)\rvert/\sqrt X\) over the full audited range is attained at
\(X=160\), where
\[
 \frac{\lvert E(160)\rvert}{\sqrt{160}}
 =0.185157076934082141384010185515317233816\ldots.
\]
On \(10{,}000\leq X\leq10^8\), the maximum is
\begin{equation}
 \begin{aligned}
 \frac{\lvert E(867{,}240)\rvert}{\sqrt{867{,}240}}
 &=0.171400192808814559728949658032345552502\ldots,\\
 R(867{,}240)&=30{,}827,
 \end{aligned}
 \label{koide-count-eq-census-tail-extremum}
\end{equation}
and the root entering at that height is
\begin{equation}
 (a,q,n)=(199997,115061,3956),\qquad
 (x,y,z)=(50380,165441,867240).
 \label{koide-count-eq-census-tail-witness}
\end{equation}
The bound \(\lvert E(X)\rvert\leq0.171\sqrt X\) first fails on this tail at
\(X=867{,}216\), where \(R(X)=30{,}826\) and
\[
 E(X)=159.4663553077412220330449353234319817199\ldots.
\]
No counterexample to the finite statement with constant \(0.172\) occurs
through \(10^8\).  These are bounded computations, not a claim that either
constant persists.

The audit finds \(12{,}916\) sign changes.  The \(3{,}536{,}002\) roots
occupy \(2{,}472{,}015\) distinct heights; the maximum shell multiplicity
is \(35\), at \(z=92{,}923{,}120\), and the largest gap between consecutive
occupied heights is \(564\), from \(71{,}586{,}276\) to \(71{,}586{,}840\).

\subsection{Independent replay and exact algorithm}

An independent pass decompressed and rechecked all \(3{,}536{,}002\)
six-coordinate rows, recomputed every checkpoint, verified the coordinate
maps, equations, ordering, chamber, height, nonduplication, and gzip CRC,
and returned \texttt{PASS}.  It also compared the literal \((n,a)\) scan
through \(X=20{,}000\), every generalized-Pell seed for \(1\leq n\leq1000\),
and the full six-coordinate point set through \(X=200{,}000\).

The reproducibility sources and receipts are retained project-relatively under
\path{supporting_materials/workbench/evidence/koide_exact_checkpoint_20260904/census}.
The two \(10^8\)-range outputs and the full coordinate receipt are in its
\path{final_100000000} subdirectory.  The principal digests are
\begin{center}
\begin{tabular}{>{\raggedright\arraybackslash}p{0.43\linewidth}
                >{\raggedright\arraybackslash}p{0.51\linewidth}}
\toprule
artifact&SHA-256\\
\midrule
\texttt{koide\_fast\_enum.py}&
\texttt{dddb73b828b63be08e2ff0b5346edb32\allowbreak%
31662005f83d5b2448c06fef6ba31017}\\
\texttt{koide\_discrepancy\_100000000.json}&
\texttt{da92f0866074e530b8174a48915caa0f\allowbreak%
65cd4cd8957e209306adc200739d1961}\\
\texttt{koide\_roots\_through\_\allowbreak 100000000.csv.gz}&
\texttt{dfc313492aedb396dcfa10520bfefc699\allowbreak%
ba601357bc9933d921dfa3884310bd4}\\
\texttt{coordinate\_audit\_receipt.json}&
\texttt{3010363108e7b460b7331a92f3263bbd\allowbreak%
77992aedfc3064973750e16d707da80f}\\
\bottomrule
\end{tabular}
\end{center}

The exact algorithm first finds the largest admissible \(n\) by integer
binary search, then factors the unchanged integers \(3n-1\) and \(3n+1\)
with a smallest-prime-factor sieve.  It enumerates every HNF triple in
\eqref{koide-count-eq-census-hnf}, constructs all roots of \(3\) modulo
\(D\) by prime-power lifting and the Chinese remainder theorem, reduces
each generator to the unique seed, and iterates
\eqref{koide-count-eq-census-unit-step} until the literal height fails.
For each \(n\), the independently checked identity
\[
 \#\mathcal P(m_n)=\rho(m_n)=\rho(3n-1)\rho(3n+1)
\]
certifies the orbit multiplicity.  The work is
\[
 X^{1+o(1)}+O(R(X))
\]
integer operations before the optional height sort.  This complexity and
the finite census use no conjectural discrepancy estimate.

% Integrated 4 September 2026; exact in-project provenance is recorded below.
% The algebra, order, lattice, shell, centralizer, and cutoff are retained exactly.
\section{Exact quaternion bridge: order, shell, and hyperbolic conjugacy classes}\label{koide-quat-exact-quaternion-bridge-order-shell-and-hyperbolic-conjugacy-classes}

The proposed quaternionic realization is correct. The displayed lattice is an order in an indefinite rational quaternion division algebra of discriminant \(6\). The parity-restricted ternary shell is exactly the reduced-trace-\(4\), reduced-norm-\(1\) locus in that order. Every resulting element is primitive hyperbolic, with eigenvalues \(2\pm\sqrt3\) and translation length \(2\log(2+\sqrt3)\).

The centralizer statement is correct when written intrinsically as

\[
 \mathcal O\cap\mathbb Q(\gamma)
 =\mathbb Z[\gamma-2]\cong\mathbb Z[\sqrt3].
\]

The entire integral shell is a finite union of conjugacy classes under \(\Gamma=\mathcal O^1\). This gives an exact reduction of the Koide count to a finite sum of sharp, non-radial sector counts inside primitive hyperbolic conjugacy classes of a cocompact arithmetic Fuchsian group.

It does \textbf{not} prove the open estimate

\[
 R(X)-C_0X=O_\varepsilon(X^{1/2+\varepsilon}).
\]

Rather, it converts the count exactly into a sharp sector problem on a finite union of hyperbolic conjugacy classes. This places it beside, but does not prove it equivalent to, the standard hyperbolic conjugacy-class lattice-point problem. Even for the radial Huber cutoff, the general pointwise theorem gives an \(O(X^{2/3})\) remainder after all exceptional spectral terms are included, while the \(O(X^{1/2+\varepsilon})\) remainder is conjectural. The Koide cutoff is a different sharp polyhedral-sector height, so the standard radial theorem cannot be imported without a new transfer or trace-formula argument.

\subsection{Quaternion algebra and reduced invariants}\label{koide-quat-quaternion-algebra-and-reduced-invariants}

Let

\[
 B=(18,-6)_{\mathbb Q}
 =\mathbb Q\langle I,J\rangle/
 (I^2-18,\ J^2+6,\ IJ+JI),
 \qquad K=IJ.
\]

Then

\[
 K^2=-I^2J^2=108.
\]

For

\[
 x=x_0+x_1I+x_2J+x_3K,
\]

the standard involution, reduced trace, and reduced norm are

\[
 \bar x=x_0-x_1I-x_2J-x_3K,
 \qquad \operatorname{trd}(x)=x+\bar x=2x_0,
\]

and

\begin{equation}\label{koide-quat-eq-reduced-norm}
 \operatorname{nrd}(x)=x\bar x
 =x_0^2-18x_1^2+6x_2^2-108x_3^2.
\end{equation}

All signs in \eqref{koide-quat-eq-reduced-norm} are forced by \(\operatorname{nrd}(I)=-18\), \(\operatorname{nrd}(J)=6\), and \(\operatorname{nrd}(K)=-108\).

\subsection{The proposed lattice is an order}\label{koide-quat-the-proposed-lattice-is-an-order}

Put

\[
 E=\frac{I+J}{2},\qquad
 F_0=\frac{I-J}{2},\qquad
 G=\frac K2,
\]

and let

\begin{equation}\label{koide-quat-eq-order-lattice}
 \mathcal O=\mathbb Z1\oplus\mathbb ZE
             \oplus\mathbb ZF_0\oplus\mathbb ZG.
\end{equation}

The displayed elements are a rational basis of \(B\), because

\[
 I=E+F_0,\qquad J=E-F_0,\qquad K=2G.
\]

Direct multiplication gives the complete nontrivial table

\begin{equation}\label{koide-quat-eq-multiplication-table}
\begin{array}{c|ccc}
   &E&F_0&G\\ \hline
 E&3&6-G&6E-3F_0\\
 F_0&6+G&3&3E-6F_0\\
 G&-6E+3F_0&-3E+6F_0&27.
\end{array}
\end{equation}

Every structure constant in \eqref{koide-quat-eq-multiplication-table} is integral. Thus \eqref{koide-quat-eq-order-lattice} is a unital subring and a full rank-four \(\mathbb Z\)-lattice in \(B\), hence an order. In particular, the unital ring sometimes denoted \(\mathbb Z[1,E,F_0,G]\) is exactly the direct sum in \eqref{koide-quat-eq-order-lattice}.

An arbitrary order element is

\begin{equation}\label{koide-quat-eq-order-coordinates}
 x=r+sE+tF_0+uG
  =r+\frac{s+t}{2}I+\frac{s-t}{2}J+\frac u2K
 \qquad(r,s,t,u\in\mathbb Z),
\end{equation}

and \eqref{koide-quat-eq-reduced-norm} gives

\begin{equation}\label{koide-quat-eq-order-trace-norm}
 \operatorname{trd}(x)=2r,
 \qquad
 \operatorname{nrd}(x)
 =r^2-3s^2-12st-3t^2-27u^2.
\end{equation}

The reduced-trace pairing on the basis \((1,E,F_0,G)\) is

\[
 \bigl(\operatorname{trd}(b_ib_j)\bigr)_{i,j}
 =
 \begin{pmatrix}
  2&0&0&0\\
  0&6&12&0\\
  0&12&6&0\\
  0&0&0&54
 \end{pmatrix},
\]

with determinant

\begin{equation}\label{koide-quat-eq-order-discriminant-determinant}
 -11664=-108^2.
\end{equation}

Thus the reduced discriminant of this order is \(108\), not \(6\). The algebra has discriminant \(6\), proved below, so \(\mathcal O\) is nonmaximal and has index \(18\) in a maximal order containing it. This index is not an Eichler level: the order is nonmaximal at the ramified primes \(2\) and \(3\), so it is not an Eichler order in the standard sense. Any optimal-embedding or trace-formula result that assumes a maximal order must account for this exact nonmaximal order instead of replacing \(\mathcal O\).

\subsection{Trace-zero lattice and exact shell bijection}\label{koide-quat-trace-zero-lattice-and-exact-shell-bijection}

Equation~\eqref{koide-quat-eq-order-trace-norm} gives

\[
 \mathcal O^0
 :=\{x\in\mathcal O:\operatorname{trd}(x)=0\}
 =\mathbb ZE\oplus\mathbb ZF_0\oplus\mathbb ZG.
\]

Writing

\[
 q=s+t,\qquad a=s-t,\qquad n=u
\]

gives the exact description

\begin{equation}\label{koide-quat-eq-trace-zero-lattice}
 \boxed{
 \mathcal O^0=
 \left\{
 \frac{qI+aJ+nK}{2}:
 a,q,n\in\mathbb Z,\ a\equiv q\pmod2
 \right\}.}
\end{equation}

The inverse formulas are

\[
 s=\frac{q+a}{2},\qquad
 t=\frac{q-a}{2},\qquad
 u=n,
\]

which are integral exactly when \(a\equiv q\pmod2\).

Define

\[
 F(a,q,n)=a^2-3q^2-18n^2,
 \qquad
 v(a,q,n)=\frac{qI+aJ+nK}{2},
\]

and

\begin{equation}\label{koide-quat-eq-gamma-parametrization}
 \gamma(a,q,n)=2+v(a,q,n).
\end{equation}

Distinct pure quaternion basis elements anticommute, so

\begin{equation}\label{koide-quat-eq-pure-square}
 v(a,q,n)^2
 =\frac{18q^2-6a^2+108n^2}{4}
 =-\frac32F(a,q,n),
\end{equation}

and

\[
 \operatorname{nrd}(v)=\frac32F(a,q,n).
\]

Since \(v\) has reduced trace zero,

\begin{equation}\label{koide-quat-eq-gamma-trace-norm}
 \operatorname{trd}(\gamma)=4,
 \qquad
 \operatorname{nrd}(\gamma)
 =4+\operatorname{nrd}(v)
 =4+\frac32F(a,q,n).
\end{equation}

Therefore \eqref{koide-quat-eq-gamma-parametrization} defines a bijection

\begin{equation}\label{koide-quat-eq-shell-bijection}
 \boxed{
 \left\{(a,q,n)\in\mathbb Z^3:
 F(a,q,n)=-2,\ a\equiv q\pmod2\right\}
 \xrightarrow{\ \sim\ }
 \left\{\gamma\in\mathcal O:
 \operatorname{trd}(\gamma)=4,\
 \operatorname{nrd}(\gamma)=1\right\}.}
\end{equation}

The inverse sends

\[
 \gamma=2+sE+tF_0+uG
\]

to

\[
 (a,q,n)=(s-t,s+t,u).
\]

There is no hidden congruence loss. On the integral equation \(F=-2\), reduction modulo \(2\) gives

\[
 a^2-q^2\equiv0\pmod2,
\]

so \(a\equiv q\pmod2\) is automatic. The parity condition in \eqref{koide-quat-eq-shell-bijection} records the order lattice exactly but removes no integral shell point.

Finally,

\begin{equation}\label{koide-quat-eq-gamma-quadratic}
 (\gamma-2)^2=3,\qquad
 \gamma^2-4\gamma+1=0,\qquad
 \gamma^{-1}=4-\gamma=2-v\in\mathcal O.
\end{equation}

Thus every element on the right of \eqref{koide-quat-eq-shell-bijection} lies in \(\Gamma=\mathcal O^1\).

\subsection{\texorpdfstring{Division over \(\mathbb Q\) and splitting over \(\mathbb R\)}{Division over \textbackslash mathbb Q and splitting over \textbackslash mathbb R}}\label{koide-quat-division-over-mathbb-q-and-splitting-over-mathbb-r}

Square factors may be removed from a Hilbert-symbol presentation, so

\[
 (18,-6)_{\mathbb Q}\simeq(2,-6)_{\mathbb Q}.
\]

At \(p=3\), write

\[
 18=3^2\cdot2,\qquad -6=3\cdot(-2).
\]

The odd-prime Hilbert-symbol formula gives

\begin{equation}\label{koide-quat-eq-hilbert-three}
 (18,-6)_3
 =(-1)^{2\cdot1\cdot(3-1)/2}
   \left(\frac2{3}\right)
   \left(\frac{-2}{3}\right)^2
 =-1.
\end{equation}

Thus \(B\otimes\mathbb Q_3\) is division, and \(B\) is a division algebra over \(\mathbb Q\).

At \(p=2\), write

\[
 18=2\cdot9,\qquad -6=2\cdot(-3).
\]

The \(2\)-adic formula

\[
 (2^\alpha u,2^\beta w)_2
 =(-1)^{
 ((u-1)/2)((w-1)/2)
 +\alpha(w^2-1)/8+\beta(u^2-1)/8}
\]

gives

\begin{equation}\label{koide-quat-eq-hilbert-two}
 (18,-6)_2=-1.
\end{equation}

At every odd prime other than \(3\), both valuations are zero and the Hilbert symbol is \(+1\). At the real place, \(18>0\), so the real Hilbert symbol is \(+1\). Consequently

\begin{equation}\label{koide-quat-eq-ramification-discriminant}
 \operatorname{Ram}(B)=\{2,3\},\qquad
 \operatorname{disc}(B)=6,
\end{equation}

and

\[
 B\otimes_{\mathbb Q}\mathbb R\simeq M_2(\mathbb R).
\]

An explicit real splitting is

\begin{equation}\label{koide-quat-eq-real-splitting}
 \rho(I)=
 \begin{pmatrix}3\sqrt2&0\\0&-3\sqrt2\end{pmatrix},
 \qquad
 \rho(J)=
 \begin{pmatrix}0&6\\-1&0\end{pmatrix}.
\end{equation}

These matrices square to \(18\) and \(-6\) and anticommute. Therefore

\[
 \rho(K)=
 \begin{pmatrix}0&18\sqrt2\\3\sqrt2&0\end{pmatrix}.
\]

For the element associated with \((a,q,n)\),

\begin{equation}\label{koide-quat-eq-gamma-matrix}
 \rho(\gamma)=
 \begin{pmatrix}
 2+\frac{3\sqrt2}{2}q&3a+9\sqrt2\,n\\
 -\frac12a+\frac{3\sqrt2}{2}n&
 2-\frac{3\sqrt2}{2}q
 \end{pmatrix}.
\end{equation}

Its determinant is the reduced norm in \eqref{koide-quat-eq-gamma-trace-norm}.

Let \(\mathcal M\) be a maximal overorder containing \(\mathcal O\).
Since \([\mathcal M:\mathcal O]=18\), the finite abelian group
\(\mathcal M/\mathcal O\) has order \(18\).  Lagrange's theorem therefore
shows that multiplication by \(18\) annihilates every element of this
additive quotient; equivalently, the ideal \((18)=18\mathbb Z\) annihilates
\(\mathcal M/\mathcal O\), and \(18\mathcal M\subset\mathcal O\).  Hence the
principal congruence subgroup
\[
 \mathcal M^1(18\mathcal M)
 :=\{u\in\mathcal M^1:u-1\in18\mathcal M\}
\]
satisfies the exact inclusions
\[
 \mathcal M^1(18\mathcal M)\subset\mathcal O^1\subset\mathcal M^1.
\]
Thus \(\mathcal O^1\) has finite index in the norm-one group of a maximal
order and contains a principal congruence subgroup.  This is the required
finite-index and congruence bridge; it does not replace the local type of
\(\mathcal O\) by the local type of \(\mathcal M\).

The image

\[
 \bar\Gamma=\rho(\mathcal O^1)/\{\pm1\}
 \subset PSL_2(\mathbb R)
\]

is a discrete congruence arithmetic Fuchsian group. Since the defining
rational quaternion algebra is division,
\(\bar\Gamma\backslash\mathbb H\) is compact; see
\cite[Lemma~38.1.2, Proposition~38.1.3, and Definition~38.7.2]{Voight_2021}.

\subsection{Centralizer, optimal quadratic order, and primitivity}\label{koide-quat-centralizer-optimal-quadratic-order-and-primitivity}

Fix a shell point and set

\[
 v=\gamma-2.
\]

Equation~\eqref{koide-quat-eq-gamma-quadratic} gives \(v^2=3\), hence

\begin{equation}\label{koide-quat-eq-quadratic-field}
 \mathbb Q(\gamma)=\mathbb Q(v)\cong\mathbb Q(\sqrt3).
\end{equation}

Since \(\gamma\) is noncentral, the centralizer theorem for quaternion algebras gives

\begin{equation}\label{koide-quat-eq-algebra-centralizer}
 C_B(\gamma)=\mathbb Q(\gamma).
\end{equation}

Moreover,

\[
 \mathbb Z[v]\subseteq\mathcal O\cap\mathbb Q(v).
\]

Every element of \(\mathcal O\) is integral over \(\mathbb Z\), so every element of the intersection is an algebraic integer in \(\mathbb Q(\sqrt3)\). Because

\[
 \mathcal O_{\mathbb Q(\sqrt3)}=\mathbb Z[\sqrt3]
\]

and \(\mathbb Z[v]\) is already isomorphic to this maximal quadratic order, both inclusions force

\begin{equation}\label{koide-quat-eq-optimal-intersection}
 \boxed{
 \mathcal O\cap\mathbb Q(\gamma)
 =\mathbb Z[v]=\mathbb Z[\gamma-2]
 \cong\mathbb Z[\sqrt3].}
\end{equation}

Thus each shell point gives an optimal embedding

\begin{equation}\label{koide-quat-eq-optimal-embedding}
 \iota_\gamma:\mathbb Z[\sqrt3]\hookrightarrow\mathcal O,
 \qquad \sqrt3\longmapsto\gamma-2.
\end{equation}

Conversely, every embedding \eqref{koide-quat-eq-optimal-embedding} gives a trace-\(4\), norm-\(1\) element \(2+\iota(\sqrt3)\), hence one integral ternary point. The shell is exactly the set of these optimal embeddings before quotienting by conjugation.

The norm-one units of \(\mathbb Z[\sqrt3]\) are

\begin{equation}\label{koide-quat-eq-norm-one-units}
 \{\pm(2+\sqrt3)^k:k\in\mathbb Z\}.
\end{equation}

For completeness, \(2+\sqrt3\) is the smallest positive norm-one unit greater than \(1\). If \(x+y\sqrt3>1\) has norm \(1\) and \(y>0\), then \(y\ge1\), \(x^2=1+3y^2\), and \(x+y\sqrt3\ge2+\sqrt3\). Multiplication by a suitable power of \(2-\sqrt3\) reduces any positive norm-one unit to \([1,2+\sqrt3)\), proving \eqref{koide-quat-eq-norm-one-units}.

Equations~\eqref{koide-quat-eq-algebra-centralizer}--\eqref{koide-quat-eq-norm-one-units} imply

\begin{equation}\label{koide-quat-eq-group-centralizer}
 C_\Gamma(\gamma)=\{\pm\gamma^k:k\in\mathbb Z\}.
\end{equation}

Hence \(\gamma\) is primitive in \(\Gamma\) and its image is primitive in \(\bar\Gamma\). Indeed, a root of \(\gamma\), or a root of its class modulo the center, would commute with \(\gamma\) and therefore be a unit from \eqref{koide-quat-eq-norm-one-units}; the exponent equation \(rk=1\) has no solution with \(k\ge2\).

The reduced characteristic polynomial is

\[
 T^2-4T+1,
\]

so the real eigenvalues are

\begin{equation}\label{koide-quat-eq-eigenvalues}
 2+\sqrt3,\qquad 2-\sqrt3.
\end{equation}

For a hyperbolic element of \(PSL_2(\mathbb R)\),

\[
 2\cosh(\ell/2)=|\operatorname{tr}(\rho(\gamma))|=4.
\]

Thus every shell element has primitive translation length

\begin{equation}\label{koide-quat-eq-translation-length}
 \boxed{
 \ell_0=2\operatorname{arcosh}(2)
       =2\log(2+\sqrt3).}
\end{equation}

\subsection{Integral adjoint action and finitely many classes}\label{koide-quat-integral-adjoint-action-and-finitely-many-classes}

Let

\[
 \mathscr S
 =\{\gamma\in\Gamma:\operatorname{trd}(\gamma)=4\}.
\]

By \eqref{koide-quat-eq-shell-bijection}, \(\mathscr S\) is exactly the entire integral ternary shell. By \eqref{koide-quat-eq-translation-length}, every element of \(\mathscr S\) is hyperbolic with translation length \(\ell_0\).

Put

\[
 L=\{(a,q,n)\in\mathbb Z^3:a\equiv q\pmod2\}.
\]

Conjugation preserves \(\mathcal O\), \(\mathcal O^0\), and reduced norm. Because

\[
 \operatorname{nrd}(v(a,q,n))=\frac32F(a,q,n),
\]

it gives an exact orthogonal representation on the actual parity lattice

\begin{equation}\label{koide-quat-eq-adjoint-representation}
 \operatorname{Ad}:\Gamma/\{\pm1\}
 \hookrightarrow SO^+(L,F).
\end{equation}

The kernel before dividing by the center is \(\Gamma\cap\mathbb Q^\times=\{\pm1\}\). No action on the larger lattice \(\mathbb Z^3\), and no finite-index equality with its full integral orthogonal group, is asserted.

There are only finitely many \(\Gamma\)-conjugacy classes in \(\mathscr S\). Choose a compact fundamental set \(D\subset\mathbb H\) for \(\bar\Gamma\), and fix \(z_0\in\mathbb H\). For each conjugacy class, conjugate a representative so that its invariant axis meets \(D\). If \(z\in D\) lies on that axis, then

\[
 d(z,\gamma z)=\ell_0.
\]

Writing

\[
 R_D=\max_{w\in D}d(z_0,w),
\]

the triangle inequality gives

\[
 d(z_0,\gamma z_0)\le2R_D+\ell_0.
\]

Proper discontinuity leaves only finitely many elements of \(\bar\Gamma\) moving \(z_0\) by at most this amount. Hence only finitely many \(\bar\Gamma\)-classes occur. This also proves finiteness of the corresponding \(\Gamma\)-classes: if two images are conjugate in \(\bar\Gamma\), lifting a conjugator gives

\[
 \delta\gamma_1\delta^{-1}=\pm\gamma_2.
\]

Both sides before the possible sign have reduced trace \(4\); the minus sign would have reduced trace \(-4\), so it is impossible.

Choose representatives \(\gamma_1,\ldots,\gamma_h\). Then

\begin{equation}\label{koide-quat-eq-shell-class-decomposition}
 \boxed{
 \mathscr S=
 \coprod_{j=1}^h
 \{\delta\gamma_j\delta^{-1}:
 [\delta]\in\Gamma/C_\Gamma(\gamma_j)\},
 \qquad
 C_\Gamma(\gamma_j)=\{\pm\gamma_j^k:k\in\mathbb Z\}.}
\end{equation}

This proves finiteness but does not compute \(h\) or a representative list. Those finite arithmetic data are needed by a class-by-class spectral formula whose coefficients contain geodesic periods.

\subsection{Exact reformulation of the Koide chamber count}\label{koide-quat-exact-reformulation-of-the-koide-chamber-count}

For \(\gamma\in\mathscr S\), let

\[
 c(\gamma)=(a(\gamma),q(\gamma),n(\gamma))
\]

be the inverse coordinate map in
\cref{koide-quat-trace-zero-lattice-and-exact-shell-bijection}. Let

\[
 \mathcal C=\{(a,q,n)\in\mathbb R^3:
 a>0,\ q\ge0,\ n\ge0,\ a+4n>q\},
\]

and retain the canonical height

\[
 H(a,q,n)=4a+17n.
\]

The integral bijections in
\cref{koide-count-thm-coordinate-bijection,koide-param-thm:integral-parametrization}
give

\begin{equation}\label{koide-quat-eq-koide-class-count}
 \boxed{
 R(X)=\sum_{j=1}^h
 \#\left\{
 [\delta]\in\Gamma/C_\Gamma(\gamma_j):
 \begin{array}{l}
 c(\delta\gamma_j\delta^{-1})\in\mathcal C,\\
 H(c(\delta\gamma_j\delta^{-1}))\le X
 \end{array}
 \right\}.}
\end{equation}

There is no hidden finite-to-one factor in
\eqref{koide-quat-eq-koide-class-count}: cosets by the centralizer
parametrize the elements of each class exactly once, and
\eqref{koide-quat-eq-shell-bijection} parametrizes shell elements by triples
exactly once.  Put \(C_j=C_\Gamma(\gamma_j)\).  Here
\(\Gamma/C_j\) is the set of left cosets \(\delta C_j\), not the set of
right cosets \(C_j\delta\), and the exact orbit bijection is

\[
 \Phi_j:\Gamma/C_j\longrightarrow
 \{\delta\gamma_j\delta^{-1}:\delta\in\Gamma\},
 \qquad
 \delta C_j\longmapsto\delta\gamma_j\delta^{-1}.
\]

Indeed, right multiplication of \(\delta\) by an element of \(C_j\) does
not change its conjugate, so \(\Phi_j\) is well defined and surjective; and

\[
 \Phi_j(\delta_1C_j)=\Phi_j(\delta_2C_j)
 \iff \delta_2^{-1}\delta_1\in C_j
 \iff \delta_1C_j=\delta_2C_j,
\]

so it is injective.  If \(\pi:\Gamma\to\bar\Gamma\) is the quotient map,
then

\[
 \Gamma/C_j
 \cong
 \bar\Gamma/C_{\bar\Gamma}(\bar\gamma_j)
\]

via \(\delta C_j\mapsto\pi(\delta)C_{\bar\Gamma}(\bar\gamma_j)\).
To prove the required equality of stabilizers, observe first that
\(\pi(C_j)\subseteq C_{\bar\Gamma}(\bar\gamma_j)\).  Conversely, if
\(\pi(\delta)\) centralizes \(\bar\gamma_j\), then
\(\delta\gamma_j\delta^{-1}=\pm\gamma_j\).  Reduced-trace invariance and
\(\operatorname{trd}(\gamma_j)=4\) exclude the minus sign, so
\(\delta\in C_j\).  Thus
\(\pi(C_j)=C_{\bar\Gamma}(\bar\gamma_j)\); together with
\(\ker\pi=\{\pm1\}\subset C_j\), this proves the displayed bijection and
shows that passage to \(PSL_2\) introduces no factor \(2\). Positive trace
selects the lift sign, but the chamber is not invariant under
\(\gamma\mapsto\gamma^{-1}\); inverse classes and their orientations must
still be retained separately.

Under the splitting \eqref{koide-quat-eq-real-splitting},

\begin{equation}\label{koide-quat-eq-frobenius-norm}
 \|\rho(\gamma)\|_{\mathrm F}^2
 =8+9q^2+\frac{37}{4}a^2
   +\frac{105\sqrt2}{2}an+\frac{333}{2}n^2.
\end{equation}

For \(i\in\mathbb H\),

\[
 2\cosh d(i,\rho(\gamma)i)=\|\rho(\gamma)\|_{\mathrm F}^2.
\]

Since \(\sinh(\ell_0/2)=\sqrt3\), the standard Huber parameter is

\begin{equation}\label{koide-quat-eq-huber-parameter}
 X_i(\gamma)^2
 =\frac{\|\rho(\gamma)\|_{\mathrm F}^2-2}{12}.
\end{equation}

The shell equation simplifies \eqref{koide-quat-eq-frobenius-norm} exactly:

\[
 \|\rho(\gamma)\|_{\mathrm F}^2
 =14+\frac{(7a+15\sqrt2\,n)^2}{4},
\]

and hence

\begin{equation}\label{koide-quat-eq-radial-height}
 \boxed{
 X_i(\gamma)^2
 =1+\frac{(7a+15\sqrt2\,n)^2}{48}.}
\end{equation}

Thus, in the positive chamber, the radial cutoff has the exact irrational linear height \(7a+15\sqrt2\,n\). The Koide height \(H=4a+17n\) is a different linear height. The two sizes are comparable:

\[
 a\le\frac H4,\qquad
 n\le\frac H{17},\qquad
 0\le q<a+4n\le\frac{33}{68}H,
\]

and conversely

\[
 H\le\sqrt{305}\sqrt{a^2+q^2+n^2}.
\]

Norm equivalence in the trace-zero space then compares \(H\) coarsely with \(X_i\). Such comparability identifies the growth scale but does not transfer a main coefficient or a remainder between the sharply bounded regions.

\subsection{Consequence for the open square-root discrepancy}\label{koide-quat-consequence-for-the-open-square-root-discrepancy}

The quaternion construction proves:

\begin{enumerate}
\def\labelenumi{\arabic{enumi}.}
\item
  the ternary shell is exactly a finite union of primitive hyperbolic conjugacy classes in one explicit cocompact arithmetic Fuchsian group;
\item
  \(\log(2+\sqrt3)=\ell_0/2\), so its occurrence in the main constant is tied to the primitive translation length;
\item
  the conjectural exponent \(1/2\) is the same exponent proposed for the standard hyperbolic lattice-point problem in conjugacy classes, giving a precise analytic analogy but not an equivalence of the two cutoff problems.
\end{enumerate}

The third item classifies the remaining analytic problem; it does not prove the
error estimate. Chatzakos--Petridis~\cite[Theorem~1.1]{ChatzakosPetridis_2016}
give

\[
 E(\mathcal H,X;z)=O(X^{2/3})
\]

for the standard radial count after the main expression, including every small-eigenvalue term, is subtracted. Their Conjecture 5.7 proposes

\[
 E(\mathcal H,X;z)=O_\varepsilon(X^{1/2+\varepsilon}).
\]

The cutoff in \eqref{koide-quat-eq-koide-class-count} uses both a sharp chamber and \(4a+17n\), not the radial height \(7a+15\sqrt2\,n\) in \eqref{koide-quat-eq-radial-height}. Their symmetric difference can have main-order size, so even the conjectured radial bound would not by itself prove the Koide bound. Therefore the status of

\[
 R(X)-C_0X=O_\varepsilon(X^{1/2+\varepsilon})
\]

remains \textbf{OPEN}. Closing it requires a spectral or automorphic argument for the exact finite sum \eqref{koide-quat-eq-koide-class-count}, with its exact sector and height.

\subsection{Missing mathematics before an error theorem can be claimed}\label{koide-quat-missing-mathematics-before-an-error-theorem-can-be-claimed}

The bridge does not supply the following results.

\begin{enumerate}
\def\labelenumi{\arabic{enumi}.}
\item
  \textbf{Class data.} Determine \(h\), representatives \(\gamma_j\), and orientation/inversion conventions. Finiteness does not compute the periods in a spectral expansion.
\item
  \textbf{Exact nonmaximal order.} The order has reduced discriminant \(108\), is nonmaximal at the ramified primes, and is not an Eichler order in the standard sense. A result for maximal or specified Eichler orders must be checked at this order or transferred through finite-index groups with all coset multiplicities retained. Index and discriminant alone do not determine the precise local order type required by a trace formula.
\item
  \textbf{Central quotient.} Track the passage from \(\Gamma\subset SL_2(\mathbb R)\) to \(\bar\Gamma\subset PSL_2(\mathbb R)\), including \(\pm1\), inverse classes, and possible orbifold torsion.
\item
  \textbf{Exact cutoff.} Prove a trace formula or effective orbital theorem for all four inequalities defining \(\mathcal C\) and for \(4a+17n\le X\). The radial Huber cutoff is different.
\item
  \textbf{Sharp-boundary control.} Remove smoothing with an explicit estimate at the required scale. The two Pell faces contribute only \(O(\log X)\), as already proved canonically, but this alone does not control every smoothing strip adjacent to all chamber walls.
\item
  \textbf{Exceptional spectrum.} Once an appropriate sector spectral expansion is constructed, one must determine whether it contains \(X^{s_j}\) terms for \(\lambda_j=s_j(1-s_j)<1/4\), weighted by closed-geodesic periods. To subtract only \(C_0X\) with a \(O(X^{1/2+\varepsilon})\) remainder, every relevant exceptional term must be shown to vanish or cancel, or it must be displayed and subtracted. No such vanishing is proved here.
\item
  \textbf{Pointwise square-root estimate.} After secondary terms and the exact sector are handled, the needed pointwise \(1/2+\varepsilon\) bound remains at the conjectural strength of the classical problem. Mean-square results do not give it for every \(X\).
\item
  \textbf{Main coefficient.} A spectral derivation must reproduce \[
  C_0=\frac{9}{2\pi^2}\log(2+\sqrt3)
      \log\!\left(\frac3{2\sqrt2}\right)
  \] using the same lattice, chamber, strict walls, height, and Haar measure. The existing local-density proof fixes this constant but does not calibrate a new trace formula automatically.
\item
  \textbf{Transfer between the two exact heights.} Formula~\eqref{koide-quat-eq-radial-height} already gives the exact standard radial parameter. What is missing is a remainder-level transfer between its height \(7a+15\sqrt2\,n\) and the canonical height \(4a+17n\), together with the chamber walls. Coarse norm comparability does not preserve a discrepancy estimate.
\end{enumerate}

These are concrete mathematical tasks. None invalidates the quaternion bridge; together they explain why it does not close the open claim.

\subsection{Reproducibility and references}\label{koide-quat-reproducibility-and-references}

The canonical quaternion replay is
\path{supporting_materials/workbench/evidence/koide_exact_checkpoint_20260904/quaternion/verify_quaternion_bridge.py}
(SHA-256
\path{6675ba2fc4fa8c2b10e14558e412178b6b5b579c9c96afca78e1db9bcdf5c6b4}).
Its recorded successful output is
\path{supporting_materials/workbench/evidence/koide_exact_checkpoint_20260904/quaternion/verification.log}
(SHA-256
\path{66853c6191a000c7c827aacd1e0e33dcf34dbe3dd61d746f66fa4d805f34c91e}),
and its local manifest is
\path{supporting_materials/workbench/evidence/koide_exact_checkpoint_20260904/quaternion/SHA256SUMS.txt}
(SHA-256
\path{f9d044db0d4b003e7aef7649b3394e6b299c09a6b2eb7adca095f8f9100ba623}).
The checkpoint-wide manifest is
\path{supporting_materials/workbench/evidence/koide_exact_checkpoint_20260904/SHA256SUMS}
(SHA-256
\path{8d108eb928e5d07093badde169f443a4a38b6fcf3eb59d5192bdd33ecef8d7a4}).
The replay checks the basis multiplication table; the trace, norm, and
coordinate formulas for
\((r,s,t,u)\in\{-4,-3,\ldots,4\}^4\); the shell identities for
\((a,q,n)\in\{-24,-23,\ldots,24\}^2\times\{-8,-7,\ldots,8\}\), containing
\(48\) shell points; the
inverse identity on those points; the radial-height identity on that box;
and the Hilbert-symbol formulas at \(2\) and \(3\).  These finite ranges are
consistency checks, not proofs outside the stated boxes.  The exact canonical
sources used for the counting statements imported into this bridge are:

\begin{itemize}
\item
  \path{supporting_materials/workbench/research/koide_counting_repair.tex},
  SHA-256
  \path{965b0a3dca8bfe4f22b9e9208368fa04697070fd25bd06bd67f2481d17cb9d9d};
\item
  \path{supporting_materials/workbench/research/koide_exact_parametrization.tex},
  SHA-256
  \path{a457709dcce0fce0774eb9326b84b00d5d8144c2087f480da704974e499bcdd3};
\item
  \path{supporting_materials/workbench/research/koide_exact_census.tex},
  SHA-256
  \path{18ed5968b7e67d9a0b23d8675b40c37ba46170ad29efdbb9ce10ed632c99600a}.
\end{itemize}

The bounded coordinate-bijection replay, whose scope does not include the
quaternion multiplication table, Hilbert symbols, centralizer, conjugacy
finiteness, or spectral conclusions, is
\path{supporting_materials/workbench/reproduction/t1_koide_exact/verify_koide_narain_root_density.py}
(SHA-256
\path{e4bf2982b26b07add6057b5d2d5c771a52ef52c629b5ad5f0aab24e440e41986}).
Its stored height-\(200000\) output is
\path{supporting_materials/workbench/reproduction/t1_koide_exact/koide_verification_200000.json}
(SHA-256
\path{b907f0ee5bdd3ad80622b669722d218b79a11a5baab0d4571e141936100706b7});
the exact invocation and scope are stated in
\path{supporting_materials/workbench/reproduction/t1_koide_exact/README.md}
(SHA-256
\path{5e9cdcb9e67954c37d8fde69e13549172aa03245674b358e4a1d25f8d9d3bff0}).
The quaternion identities and deductions in this section are proved in the
displayed algebraic calculations themselves, not inferred from that bounded
replay.

References:

\begin{itemize}
\item
  John Voight, \emph{Quaternion Algebras}, Graduate Texts in Mathematics 288,
  especially Section 10.1, Lemma 38.1.2, and Proposition 38.1.3;
  \href{https://doi.org/10.1007/978-3-030-56694-4}{DOI}.
\item
  Dimitrios Chatzakos and Yiannis N. Petridis,
  \emph{The hyperbolic lattice point problem in conjugacy classes},
  arXiv:1504.01307v2, Theorem 1.1 and Conjecture 5.7;
  \href{https://arxiv.org/abs/1504.01307}{arXiv}.
\item
  The exact Koide chamber, height, main term, and Pell-face estimates are in
  Sections~\ref{koide-count-sec-main}--\ref{koide-count-sec-census-1e8}
  of this workbench.
\end{itemize}

% Integrated 4 September 2026; exact in-project provenance is recorded below.
% The exact fixed level, lattice, chamber, height, theorem hypotheses, and locators are retained.
\section{Primary-source theorem audit of effective counting for the fixed Koide ternary quadric sector}\label{koide-lit-primary-source-theorem-audit-of-effective-counting-for-the-fixed-koide-ternary-quadric-sector}

\textbf{Date:} 4 September 2026\\
\textbf{Exact counting function:}

\[
 R(X)=\#\left\{(a,q,n)\in\mathbb Z^3:
 \begin{array}{l}
 a^2-3q^2-18n^2=-2,\\
 a>0,\ q\geq0,\ n\geq0,\ a+4n>q,\\
 4a+17n\leq X
 \end{array}\right\}.
\]

The leading constant proved in
\cref{koide-count-thm-linear-asymptotic} is

\[
 C_0=\frac{9}{2\pi^2}\log(2+\sqrt3)
       \log\!\left(\frac{3}{2\sqrt2}\right).
\]

The canonical project-relative mathematical sources used by this audit are:

\begin{itemize}
\item
  \path{supporting_materials/workbench/research/koide_counting_repair.tex},
  which proves the coordinate bijection, local densities, invariant-volume
  factor, Gorodnik--Oh--Shah hypotheses, and calibrated constant, SHA-256
  \path{965b0a3dca8bfe4f22b9e9208368fa04697070fd25bd06bd67f2481d17cb9d9d};
\item
  \path{supporting_materials/workbench/research/koide_exact_parametrization.tex},
  which proves the exact parametrization and the Benoist--Oh power-saving
  application, SHA-256
  \path{a457709dcce0fce0774eb9326b84b00d5d8144c2087f480da704974e499bcdd3};
\item
  \path{supporting_materials/workbench/research/koide_quaternion_bridge.tex},
  which proves the exact order, shell, centralizer, conjugacy-class, and
  radial-height statements used below, SHA-256
  \path{93759f7ba99870c84c9144a373a53d2e9a8bc3ae909518e47991373ace0b47bb};
\item
  \path{supporting_materials/workbench/research/koide_exact_census.tex},
  which records the separately bounded finite census, SHA-256
  \path{18ed5968b7e67d9a0b23d8675b40c37ba46170ad29efdbb9ce10ed632c99600a}.
\end{itemize}

The primary publications themselves are identified by exact versions,
dates where version-sensitive, locators, and DOI or arXiv links in
\cref{koide-lit-primary-references-and-exact-locators}.  They are not
claimed to be stored in a separate canonical evidence manifest.

\subsection{Determination}\label{koide-lit-determination}

Among the source versions enumerated in
\cref{koide-lit-primary-references-and-exact-locators}, no theorem proves,
for this exact fixed level, full integral lattice, sharp chamber, and sharp
polyhedral height, the unconditional estimate

\[
 R(X)-C_0X=O_\varepsilon(X^{1/2+\varepsilon}).
\]

The strongest exact-domain conclusion obtained from the general theorems
enumerated there is

\[
 \boxed{R(X)=C_0X+O(X^{1-\eta})\quad\text{for some unspecified }\eta>0.}
\]

It follows by applying Benoist--Oh, Definition 13.1 and Theorem 13.2, orbit by orbit after directly verifying effective well-roundedness for the present family. Benoist--Oh do not give a numerical lower bound for \(\eta\), and their proof does not imply \(\eta\geq 1/2\).

The superficially stronger exponents in nearby problems do not answer the question:

\begin{itemize}

\item
  Good's sharp hyperbolic conjugacy-class theorem, as formulated and reproved by Chatzakos--Petridis, gives \(O(X^{2/3})\) only for a different radial cutoff and only after every small-eigenvalue term \(X^{s_j}\) has been placed in the main expression. The same paper states \(O_\varepsilon(X^{1/2+\varepsilon})\) as Conjecture 5.7.
\item
  Liu--Sarnak obtain a smooth radial error \(O_\varepsilon(X^{1/2+\theta+\varepsilon})\), hence \(O_\varepsilon(X^{39/64+\varepsilon})\) from \(\theta=7/64\), and square-root strength only under the generalized Selberg/Ramanujan-at-infinity conjecture for the relevant Jacquet--Langlands transfers. Their paper's global affine-sieve setup assumes \(t\det(f)\) square-free, but the analytic Proposition 3.2 itself is formulated for a variable cocompact lattice with a spectral gap. After the exact quaternionic identification it applies class by class to this lattice. What it still does not supply is a sharp cutoff, the coordinate chamber, or the nonradial height \(4a+17n\).
\item
  Gorodnik--Nevo's published sharp polar-sector theorem loses powers through set thickening. For \(\dim G=3\), it gives \(O(X^{19/20+\epsilon})\) even with tempered spectrum; the Kim--Sarnak complementary-series threshold, which does apply to the exact congruence arithmetic lattice by Jacquet--Langlands, gives the still weaker \(O(X^{39/40+\epsilon})\). This remains a comparison theorem because Gorodnik--Nevo's fixed polar-sector family has not been identified with the exact polyhedron, including all integral orbits.
\item
  Square-root-scale statements in Chatzakos--Petridis, Chatzakos, Voskou, and Lekkas--Petridis are mean-square, averaged, Omega, or conjectural statements. They are not pointwise sharp-sector upper bounds.
\item
  Recent split-stabilizer theorems assume \(-m\det(F)\) is a square and have a \(T\log T\)-type main term. Here \(-(-2)\det(F)=108\) is not a square, so they concern the opposite arithmetic regime.
\end{itemize}

The two independent missing inputs are therefore (i) a sharp-boundary spectral or arithmetic estimate for the actual height and chamber, and (ii) control or elimination of any complementary-spectrum secondary terms for the actual sector. The Pell faces \(q=0\) and \(n=0\), which contribute only \(O(\log X)\), do not solve the moving outer-boundary problem.

\subsection{Exact arithmetic and geometric classification}\label{koide-lit-exact-arithmetic-and-geometric-classification}

\subsubsection{The rational group and stabilizer}\label{koide-lit-the-rational-group-and-stabilizer}

Let

\[
 F(a,q,n)=a^2-3q^2-18n^2,
 \qquad \mathscr X=\{F=-2\}.
\]

The matrix of \(F\) is \(\operatorname{diag}(1,-3,-18)\), so \(\det(F)=54\) and

\[
 -(-2)\det(F)=108\notin\mathbb Q^{\times2}.
\]

The homogeneous form \(F\) is anisotropic over \(\mathbb Q\). Indeed, suppose a primitive integral triple satisfies

\[
 a^2-3q^2-18n^2=0.
\]

Modulo \(3\), one gets \(3\mid a\); write \(a=3a_1\). Dividing by \(3\) and reducing modulo \(3\) gives \(3\mid q\); write \(q=3q_1\). The resulting equation is

\[
 a_1^2=3q_1^2+2n^2.
\]

Modulo \(3\), this says \(a_1^2+n^2\equiv0\pmod3\), whence \(3\mid a_1,n\). Thus \(3\mid a,q,n\), contradicting primitivity.

Put \(G=\operatorname{Spin}(F)\) and \(x_0=(1,1,0)\). The orthogonal complement of \(x_0\) is

\[
 x_0^\perp=\{(3s,s,t):s,t\in\mathbb Q\},
 \qquad F(3s,s,t)=6(s^2-3t^2).
\]

Consequently the connected stabilizer \(H=\operatorname{Stab}_G(x_0)^\circ\) is the norm-one torus

\[
 \operatorname{Res}^{1}_{\mathbb Q(\sqrt3)/\mathbb Q}\mathbb G_m.
\]

It is split over \(\mathbb R\), anisotropic over \(\mathbb Q\), and has no
nontrivial \(\mathbb Q\)-character.  Under the quaternionic identification
in \cref{koide-quat-centralizer-optimal-quadratic-order-and-primitivity},
\(x_0\) corresponds to \(v_0=(I+J)/2\), with \(v_0^2=3\), and

\[
 \mathcal O\cap\mathbb Q(v_0)=\mathbb Z[v_0],
 \qquad
 H(\mathbb Q)\cap\Gamma
 =\{\pm(2+v_0)^k:k\in\mathbb Z\}.
\]

Thus \((H(\mathbb Q)\cap\Gamma)/\{\pm1\}\) is infinite cyclic.  Under
\(H(\mathbb R)^\circ\cong\mathbb R_{>0}\), the logarithm sends its positive
part to \(\log(2+\sqrt3)\mathbb Z\subset\mathbb R\); the quotient is compact,
and in particular \((H\cap\Gamma)\backslash H(\mathbb R)\) has finite
volume. Since \(F\) is \(\mathbb Q\)-anisotropic, the arithmetic Fuchsian
quotient attached to \(G\) is cocompact as well. There is no ambient
hyperbolic cusp.

This is the finite-volume stabilizer regime associated with a closed geodesic.
It is not the rationally split regime with divergent stabilizer in which one
obtains \(T\log T\) growth.

\subsubsection{The exact height really is a norm-sector height}\label{koide-lit-the-exact-height-really-is-a-norm-sector-height}

Define

\[
 \|(a,q,n)\|_*=max\{|4a+17n|,|q|,|a|\}.
\]

This is a genuine norm because it is the pullback of the sup norm under the invertible linear map

\[
 (a,q,n)\longmapsto(4a+17n,q,a),
\]

whose determinant is \(-17\). On the target chamber,

\[
 0<a<4a+17n,
 \qquad 0\leq q<a+4n<4a+17n,
\]

and therefore \(\|(a,q,n)\|_*=4a+17n\). Thus the outer cutoff is not merely comparable to a norm: it is exactly the restriction of a fixed ambient norm.

The additional wall \(a+4n>q\) removes only a bounded real set. If \(q\geq a+4n\), then

\[
 a^2+2=3q^2+18n^2
 \geq3(a+4n)^2+18n^2,
\]

so \(2a^2+24an+66n^2\leq2\). Its only integral point in the positive orthant on \(\mathscr X\) is \((1,1,0)\). The unbounded faces \(q=0\) and \(n=0\) are Pell families and contribute \(O(\log X)\) integral points, but their real boundary strips still have to be included in a sharp-domain argument.

\subsubsection{Invariant volume with a strong enough secondary estimate}\label{koide-lit-invariant-volume-with-a-strong-enough-secondary-estimate}

Use

\[
 \sqrt3q=r\cos\theta,
 \qquad 3\sqrt2n=r\sin\theta,
 \qquad 0\leq\theta\leq\frac\pi2.
\]

On \(F=-2\), \(r=\sqrt{a^2+2}\), and the gauge measure is

\[
 dm_\infty=\frac1{6\sqrt6}\,da\,d\theta.
\]

Writing \(c=17/(3\sqrt2)\), the outer boundary is

\[
 4a+c\sqrt{a^2+2}\sin\theta=X.
\]

If \(A_X(\theta)\) is its positive solution, then uniformly in \(\theta\),

\[
 A_X(\theta)=\frac{X}{4+c\sin\theta}+O(X^{-1}).
\]

Indeed \(A_X(\theta)\asymp X\) uniformly and \(\sqrt{A^2+2}=A+O(A^{-1})\). Once \(X\) contains the bounded wall excision, it follows that

\[
 m_\infty(B_X)=I_\infty X+O(1),
 \qquad
 I_\infty=\frac1{2\sqrt3}\log\!\left(\frac3{2\sqrt2}\right).
\]

This \(O(1)\) refinement matters: a theorem whose main term is the exact invariant volume can be converted to the displayed \(C_0X\) main term without spending a power of \(X\).

\subsubsection{The full integral locus is not assumed to be one orbit}\label{koide-lit-the-full-integral-locus-is-not-assumed-to-be-one-orbit}

The exact lattice must not be silently replaced by a convenient single orbit or by a commensurable maximal-order lattice. Borel--Harish-Chandra finiteness, in the form quoted by Benoist--Oh on journal p.~1892, says that \(G(\mathbb Z)\) has finitely many orbits on \(\mathscr X(\mathbb Z)\). Thus

\[
 \mathscr X(\mathbb Z)=\bigsqcup_{j=1}^{h}x_j\Gamma
\]

for the integral spin group \(\Gamma\) attached to the chosen lattice and finitely many integral representatives \(x_j\). An effective theorem for one orbit must be applied to every \(x_j\Gamma\), and the finitely many main-measure scalars must be summed. The full-lattice scalar is calibrated from Borovoi's local/Brauer--Manin density formula in \cref{koide-count-prop-finite-sign-balance,koide-count-thm-linear-asymptotic}; it is not obtained by declaring that all integral points form the orbit of \(x_0\).

\subsubsection{Exact quaternionic lattice and its congruence status}\label{koide-lit-exact-quaternionic-lattice-and-its-congruence-status}

For the spectral comparisons below, the precise arithmetic lattice can be retained. Put

\[
 B=(18,-6)_{\mathbb Q}
   =\mathbb Q\langle I,J\rangle/(I^2-18,\ J^2+6,\ IJ+JI),
 \qquad K=IJ,
\]

and

\[
 \mathcal O=\mathbb Z1\oplus
  \mathbb Z\frac{I+J}{2}\oplus
  \mathbb Z\frac{I-J}{2}\oplus
  \mathbb Z\frac K2.
\]

The complete integral multiplication table is
\eqref{koide-quat-eq-multiplication-table}, and its reduced-trace Gram
determinant is \(-108^2\) by
\eqref{koide-quat-eq-order-discriminant-determinant}; hence
\(\operatorname{discrd}(\mathcal O)=108\).  The Hilbert-symbol calculations
\eqref{koide-quat-eq-hilbert-three} and
\eqref{koide-quat-eq-hilbert-two}, together with the unramified-place and
real-place calculation summarized in
\eqref{koide-quat-eq-ramification-discriminant}, show that the ramification
set is exactly \(\{2,3\}\).  Thus \(B\) is the indefinite quaternion
division algebra of discriminant \(6\). The map

\[
 (a,q,n)\longmapsto
 \gamma(a,q,n)=2+\frac{qI+aJ+nK}{2}
\]

is the bijection \eqref{koide-quat-eq-shell-bijection} from
\(F(a,q,n)=-2\) to the reduced-trace-\(4\), reduced-norm-one locus in
\(\mathcal O\). Indeed, the parity condition \(a\equiv q\pmod2\), which is
exactly what membership in \(\mathcal O\) requires, follows automatically
from \(F=-2\) modulo \(2\), and

\[
 \operatorname{nrd}(\gamma)=4+\frac32F(a,q,n)=1.
\]

Thus the acting Fuchsian lattice is exactly \(\Gamma=\mathcal O^1\), not a substituted maximal-order lattice. If \(\mathcal M\supset\mathcal O\) is maximal, the reduced-discriminant index formula gives

\[
 [\mathcal M:\mathcal O]=\frac{108}{6}=18,
 \qquad 18\mathcal M\subset\mathcal O,
\]

and consequently

\[
 \mathcal M^1(18\mathcal M)
 \subset\mathcal O^1\subset\mathcal M^1.
\]

Here
\(\mathcal M^1(18\mathcal M)=\{u\in\mathcal M^1:u-1\in18\mathcal M\}\).
Therefore \(\mathcal O^1\) is a congruence arithmetic Fuchsian group in the
exact sense used by Chu--Li. The number \(108\) is the order's reduced
discriminant, not an Eichler level: the order is nonmaximal at the ramified
primes \(2,3\). This non-Eichler feature changes the finite compact-open
subgroup and hence which automorphic representations have fixed vectors and
the dimensions of those fixed-vector spaces; it does not change the
Jacquet--Langlands/Kim--Sarnak archimedean bound.

\subsection{Result matrix}\label{koide-lit-result-matrix}

\begin{longtable}[]{@{}
  >{\raggedright\arraybackslash}p{0.31\linewidth}
  >{\raggedright\arraybackslash}p{0.31\linewidth}
  >{\raggedright\arraybackslash}p{0.31\linewidth}@{}}
\toprule\noalign{}
\begin{minipage}[b]{\linewidth}\raggedright
Source and exact locator
\end{minipage} & \begin{minipage}[b]{\linewidth}\raggedright
What is actually proved
\end{minipage} & \begin{minipage}[b]{\linewidth}\raggedright
Status for the target
\end{minipage} \\
\midrule\noalign{}
\endhead
\bottomrule\noalign{}
\endlastfoot
Duke--Rudnick--Sarnak (1993), Theorem 1.2, p.~145; Example 1.5, pp.~146--147 & One arithmetic orbit on an affine symmetric variety: count \(\sim\) invariant volume; ternary anisotropic stabilizer gives linear rather than \(T\log T\) growth & Correct qualitative regime, no rate and no exact sharp-sector theorem \\
Eskin--McMullen (1993), Theorem 1.4, author PDF p.~3 & Every well-rounded sequence has count asymptotic to volume & Qualitative only; no exponent \\
Gorodnik--Oh--Shah (2009), definitions before Theorem 1.6 and Theorem 1.6, p.~5; Theorem 1.18, pp.~8--9 & Any norm and an admissible generic cone: full integral count \(\sim\) volume & Direct qualitative theorem for the exact norm and cone, whose hypotheses are proved in \cref{koide-count-lem-genuine-norm,koide-count-lem-gos-hypotheses}; no error term \\
Gorodnik--Oh--Shah (2010), Theorem 1.13, author PDF pp.~4--5 & Well-rounded polar sectors and their asymptotic & Qualitative only \\
Benoist--Oh (2012), Theorem 1.3, p.~1894 & Standard product/Euclidean height ball: relative \(O(T^{-\delta})\) for some \(\delta>0\) & Theorem 1.3 alone is not the exact polyhedron \\
Benoist--Oh (2012), Definition 13.1, p.~1931, and Theorem 13.2, pp.~1931--1932 & Any independently verified effectively well-rounded family: \(\operatorname{vol}(B)(1+O(\operatorname{vol}(B)^{-\delta_0}))\) & Applies to the exact family after the proof in \cref{koide-lit-definition-13.1-for-the-actual-b_x}; yields only unspecified \(O(X^{1-\eta})\) \\
Gorodnik--Nevo (2012), Definition 3.1, p.~13; Theorem 8.1 and Corollary 8.3, pp.~34--35 & Sharp fixed polar sectors: \(O(\operatorname{vol}^{1-\zeta+\epsilon})\), \(\zeta=(2n_e(p))^{-1}(1+3\dim G)^{-1}\) & Explicit but much weaker; target polyhedron and finite orbit union require an additional reduction \\
Liu--Sarnak (2010), equations (2.3)--(2.8), pp.~4--5; Proposition 3.2, pp.~10--11; equation (3.65), p.~16 & Smooth radial count for a cocompact lattice containing the fixed cyclic stabilizer: \(O_\epsilon(T^{1/2+\theta+\epsilon})\) & Exact lattice and finite class sum are available; unconditional \(39/64\), conditional \(1/2\), but fixed-width smoothing and the radial height do not match the sharp target \\
Good (1983), Theorem 4, p.~116, as explicitly restated and reproved in Chatzakos--Petridis (2016), Theorem 1.1, p.~3 & Sharp radial hyperbolic conjugacy-class count: \(O(X^{2/3})\) after all \(X^{s_j}\) small-spectrum terms & Closest sharp spectral model, not the exact polyhedron; square root is CP Conjecture 5.7 \\
Bourgain--Kontorovich--Sarnak (2010), Theorem 1.1, p.~2 & Fixed angular Fourier modes in a sharp group norm ball: \(O_{n,k}(T^{4/3})\) after \(T^{2s_j}\) terms & With affine height \(X\asymp T^2\), again \(X^{2/3}\); constants are not uniform enough to recover a sharp interval for free \\
Browning--Gorodnik (2017), Theorem 2.1, p.~11; Corollaries 2.5--2.6, pp.~24--25 & Full Euclidean balls, finite adelic/orbit sums, and congruence conditions: unspecified \(O(\operatorname{vol}^{1-\varrho})\) & Correct nonsquare homogeneous regime, wrong height geometry and no explicit exponent \\
Cao--Huang (2022), Theorem 1.5, p.~4; Theorem 3.4, pp.~21--22 & Arithmetic purity for anisotropic nonsquare ternary quadrics; effective input for Euclidean balls has an unspecified power & Qualitative full-lattice theorem, not a square-root sector result \\
Cao--Huang--Zhang (arXiv:2509.23320), Setting 2.1 and Theorem 2.1 & Product Euclidean/max height balls with congruence-uniform \(O(\mathbf N(\mathfrak l)^\kappa\operatorname{vol}^{1-\delta})\) & Current extension, but \(\delta\) is unspecified and the growing real factor is the standard height ball, not this cone \\
Oh--Shah (2014), Theorem 1.2, p.~3; Theorem 6.1, pp.~13--14 & Divergent-geodesic/split-stabilizer \(T\log T\) count; smooth error \(O(T)\), sharp relative logarithmic error & Excluded by the finite-volume anisotropic stabilizer \\
Xu--Zhang (2023), main theorem; Huang (arXiv:2412.03350), Theorem 1.1; Zhang (arXiv:2603.21080), Theorem 1.1 & Split-stabilizer results assuming \(-m\det(F)\) is a square, with smooth or radial \(T\log T\)-type expansions & \(108\) is nonsquare; wrong main-term regime \\
Huang (arXiv:2412.05903), Theorems 1.1--1.2 & Smooth compact real weight while the level varies as \(m_0p_0^{2h}\); error of logarithmic relative strength & Varying level and smooth weight, not the fixed-level sharp count \\
Li--Liu (arXiv:2501.15033v2), Theorems 1.1--1.2, pp.~2--3; Lemma 2.3, pp.~5--6 & Almost-prime results using a smooth radial orbit count with \(O_\epsilon(d^{1+\epsilon}T^{1/2+\theta+\epsilon})\) & Current use of the Liu--Sarnak input; globally assumes \(t\det(f)\) square-free and supplies no sharp-sector count \\
Voskou (arXiv:2407.03134v3), Theorem 3 and Conjecture 1, p.~3 & Four sharp orientation sectors in a closed-geodesic double-coset problem: \(O(X^{2/3})\); square root conjectured & Current sharp comparator, not the target \\
Lekkas--Petridis (arXiv:2509.12902, version dated 24 Aug.~2026), Theorem 1.2, Theorem 1.5, Conjecture 1.6 & Sharp closed-geodesic double-coset \(O(X^{2/3})\), square-root mean square, pointwise square-root conjecture & Version-specific confirmation that mean square does not furnish the requested pointwise bound \\
\end{longtable}

\subsection{Benoist--Oh: exact applicability and exact limitation}\label{koide-lit-benoistoh-exact-applicability-and-exact-limitation}

Among the general effective theorems enumerated in
\cref{koide-lit-result-matrix}, this is the one whose stated hypotheses can
be verified for the actual family without changing its height or chamber.

\subsubsection{What the theorem assumes}\label{koide-lit-what-the-theorem-assumes}

Benoist--Oh work with a connected semisimple simply connected almost \(K\)-simple group \(G\), a finite-kernel rational representation, a symmetric subgroup \(H\), and the hypothesis that \(H^\circ\) has no nontrivial \(K\)-character; see journal pp.~1892--1893. For a congruence subgroup \(\Gamma_S\), Theorem 12.5, journal p.~1929, supplies the effective equidistribution of the translates of \((\Gamma_S\cap H_S)\backslash H_S\) used in Theorem 13.2. Example (A), journal p.~1896, explicitly treats ternary quadratic forms provided the homogeneous form does not represent zero over \(\mathbb Q\).

All of these hypotheses hold here:

\begin{itemize}

\item
  \(G=\operatorname{Spin}(F)\) is simply connected and absolutely almost simple of type \(A_1\), and \(G(\mathbb R)\) is noncompact.
\item
  The standard three-dimensional orthogonal representation has finite kernel.
\item
  The stabilizer of a nonisotropic vector is symmetric.
\item
  The mod-\(3\) descent in
  \cref{koide-lit-the-rational-group-and-stabilizer} proves that \(F\) does
  not represent zero over \(\mathbb Q\).
\item
  \(H^\circ\) is the \(\mathbb Q\)-anisotropic norm-one torus of \(\mathbb Q(\sqrt3)\), so it has no nontrivial \(\mathbb Q\)-character and its arithmetic quotient is compact.
\item
  With \(S=\{\infty\}\), the full integral spin group is itself a congruence subgroup in Benoist--Oh's sense.
\end{itemize}

There is no omitted wavefront or properness condition: symmetry of \(H\) is exactly what invokes their effective wavefront argument in Lemma 12.3 and Theorem 12.5. Cocompactness makes escape in the automorphic quotient simpler, but does not remove Definition 13.1(3), which still has to be checked on the growing affine set.

Here ``congruence'' is an exact property of the chosen integral group, not an inference from mere commensurability. Under the quaternionic realization of \(G\), the finite-integrality conditions define a compact-open subgroup of \(G(\mathbb A_f)\); equivalently, after placing the corresponding order \(\mathcal O\) in a maximal order \(\mathcal O_{\max}\), some principal subgroup \(\{u\in\mathcal O_{\max}^1:u\equiv1\pmod{N\mathcal O_{\max}}\}\) lies in the exact integral spin group. This distinction matters for spectral-gap transfer: an arbitrary finite-index noncongruence subgroup need not inherit the Kim--Sarnak gap.

\subsubsection{\texorpdfstring{Definition 13.1 for the actual \(B_X\)}{Definition 13.1 for the actual B\_X}}\label{koide-lit-definition-13.1-for-the-actual-b_x}

Let \(B_X\subset\mathscr X(\mathbb R)\) be the exact chamber cut out by the five linear conditions and \(4a+17n\leq X\). In Benoist--Oh's notation, \(B_{X,\epsilon}^+=B_XB_\epsilon\) and \(B_{X,\epsilon}^-=\bigcap_{u\in B_\epsilon}B_Xu\).

Condition (1) is void because there is no finite-place factor.

For condition (2), if \(u\) lies in an \(\epsilon\)-ball about the identity, every fixed linear functional \(L\) defining the chamber or height satisfies

\[
 |L(xu)-L(x)|\ll\epsilon\|x\|_2.
\]

On the unbounded portion of this chamber, \(\|x\|_2\asymp 1+a\asymp1+(4a+17n)\). Hence a membership change under \(u\) is confined to:

\begin{enumerate}
\def\labelenumi{\arabic{enumi}.}

\item
  an angular strip of width \(O(\epsilon)\) at \(q=0\) or \(n=0\);
\item
  an \(O(\epsilon X)\)-thick radial shell at \(4a+17n=X\); or
\item
  an \(O(\epsilon)\)-neighborhood of one of the compact boundary pieces \(a=0\) or \(a+4n=q\).
\end{enumerate}

Since \(dm_\infty=(6\sqrt6)^{-1}da\,d\theta\), their total measure is \(O(\epsilon X)\). Thus

\[
 m_\infty(B_{X,\epsilon}^+\setminus B_{X,\epsilon}^-)
 \ll \epsilon\,m_\infty(B_X),
\]

so Definition 13.1(2) holds with \(\kappa=1\).

For condition (3), Benoist--Oh's Euclidean height is comparable to \(1+a\) throughout a fixed group-neighborhood enlargement of the chamber. For every \(k>0\), uniformly in \(0<\epsilon<1\),

\[
 \int_{B_{X,\epsilon}^+}H_\infty(x)^{-k}\,dm_\infty(x)
 \ll_k\int_0^{O(X)}(1+a)^{-k}\,da.
\]

The right side is \(O(X^{1-k})\) for \(0<k<1\), \(O(\log X)\) for \(k=1\), and \(O(1)\) for \(k>1\). As \(m_\infty(B_X)\asymp X\), this is \(O(m_\infty(B_X)^{1-\delta(k)})\) for some \(\delta(k)>0\) in every case. Definition 13.1(3) is therefore satisfied.

Proposition 14.2 of Benoist--Oh verifies only their standard product height balls. It should not be cited as though it automatically covers this polyhedron. The preceding direct verification is the necessary bridge to Theorem 13.2.

\subsubsection{Finite-orbit passage and main-term normalization}\label{koide-lit-finite-orbit-passage-and-main-term-normalization}

Benoist--Oh state on journal p.~1892 that \(G(\mathcal O_S)\) has finitely many orbits on \(\mathscr X(\mathcal O_S)\). Apply Theorem 13.2 to each exact integral orbit. If \(\eta_j>0\) is the exponent for the \(j\)-th orbit, take \(\eta=\min_j\eta_j>0\). The finite sum has the form

\[
 R(X)=\kappa_{\mathrm{orb}}\,m_\infty(B_X)+O(X^{1-\eta}).
\]

Here \(\kappa_{\mathrm{orb}}\) is the sum of the orbit-volume normalization
constants.  The comparison with Borovoi's Theorem 1.4.1 and Proposition
1.5.1 is carried out in
\cref{koide-count-thm-linear-asymptotic}; specifically,
\eqref{koide-count-eq-singular-series},
\eqref{koide-count-eq-archimedean-asymptotic}, and
\eqref{koide-count-eq-final-constant} give

\[
 \kappa_{\mathrm{orb}}=\mathfrak S(\mathscr X)
 =\frac{9\sqrt3}{\pi^2}\log(2+\sqrt3),
 \qquad
 \kappa_{\mathrm{orb}} I_\infty=C_0.
\]

Together with \(m_\infty(B_X)=I_\infty X+O(1)\), this yields

\[
 R(X)=C_0X+O(X^{1-\eta}).
\]

This counts the actual integral lattice. No assertion that \(\mathscr X(\mathbb Z)\) is one orbit, and no replacement by a maximal order, is involved.

\subsubsection{Why the exponent remains unspecified}\label{koide-lit-why-the-exponent-remains-unspecified}

The proof of Benoist--Oh Theorem 13.2 balances an effective-equidistribution/Sobolev error and a boundary-thickening error:

\[
 c_1\epsilon^{-p}v_X^{1-\delta}+c_2\epsilon^\kappa v_X,
 \qquad v_X=\operatorname{vol}(B_X).
\]

The resulting saving is of the form

\[
 \eta=\frac{\delta\kappa}{p+\kappa}>0.
\]

Neither the Sobolev exponent \(p\) nor the effective equidistribution exponent \(\delta\) is made numerically strong enough in the theorem to show \(\eta\geq1/2\). Thus Benoist--Oh supplies a genuine exact-family power saving, but not the requested square-root bound.

\subsection{The sharp hyperbolic frontier}\label{koide-lit-the-sharp-hyperbolic-frontier}

\subsubsection{Good and Chatzakos--Petridis}\label{koide-lit-good-and-chatzakospetridis}

For a cocompact or cofinite Fuchsian group and a fixed hyperbolic conjugacy class \(\mathcal H\), Chatzakos--Petridis define, in equations (1.4)--(1.6),

\[
 M(\mathcal H,X;z)
 =\sum_{1/2<s_j\leq1}
 A(s_j)\widehat u_j u_j(z)X^{s_j},
\]

where \(\lambda_j=s_j(1-s_j)<1/4\) and \(\widehat u_j\) is the period of \(u_j\) along the invariant closed geodesic. Their Theorem 1.1, author manuscript p.~3, restates Good's Theorem 4, p.~116, and proves

\[
 N(\mathcal H,X;z)-M(\mathcal H,X;z)=O(X^{2/3}).
\]

The \(2/3\) arises from the sharp majorant/minorant error

\[
 O(XY^{-1/2}+Y+X^{1/2}\log X)
\]

after choosing \(Y=X^{2/3}\). Theorems 1.2--1.3 are mean-square results, not pointwise improvements. Conjecture 5.7, author manuscript p.~22, explicitly asks for

\[
 O_\varepsilon(X^{1/2+\varepsilon}).
\]

There are two independent failures of direct transfer to \(R(X)\):

\begin{enumerate}
\def\labelenumi{\arabic{enumi}.}

\item
  the Huber cutoff is radial, whereas \(4a+17n\leq X\) plus the chamber is a nonradial sharp polyhedron; norm comparability does not preserve the leading constant or error because the symmetric difference can have volume of order \(X\);
\item
  even in the radial problem, the theorem subtracts every nonzero small-eigenvalue contribution, not only the linear term.
\end{enumerate}

For the exact Koide sector no spectral expansion establishing actual \(X^{s_j}\) terms has yet been produced. They are therefore a necessary obstruction to check, not terms that this audit claims are already known to occur. A proof of \(R(X)-C_0X=O(X^{1/2+\varepsilon})\) would have to show absence or vanishing of all relevant complementary-spectrum coefficients, or explicitly subtract them and then prove they vanish in the target formulation.

\subsubsection{\texorpdfstring{Current sharp refinements remain at \(2/3\)}{Current sharp refinements remain at 2/3}}\label{koide-lit-current-sharp-refinements-remain-at-23}

Voskou, Theorem 3 of arXiv:2407.03134v3, proves \(O(X^{2/3})\) for four prescribed orientation/side sectors in a closed-geodesic-to-closed-geodesic problem, with all small-spectrum period terms included in the main expression. Conjecture 1 on the same page proposes \(O_\varepsilon(X^{1/2+\varepsilon})\), while Theorem 4 is a mean-square result.

Lekkas--Petridis, in the 24 August 2026 version of arXiv:2509.12902, Theorem 1.2, likewise prove

\[
 N(X,l)=\sum_{1/2<s_j\leq1}
 \frac2\pi\gamma_1(s_j)|\widehat u_j|^2X^{s_j}
 +O(X^{2/3})
\]

for a closed-geodesic double-coset count. Their Theorem 1.5 gives

\[
 \frac1X\int_X^{2X}|E(x,l)|^2\,dx\ll X\log^2X,
\]

and Conjecture 1.6 proposes the pointwise \(X^{1/2+\epsilon}\) bound. Their introduction states that Good's \(2/3\) has not been improved for any of the nine point/geodesic/cusp double-coset configurations. This is strong current evidence about the sharp frontier, but neither paper treats the exact Koide height.

\subsubsection{Fixed angular harmonics do not give a sharp interval for free}\label{koide-lit-fixed-angular-harmonics-do-not-give-a-sharp-interval-for-free}

Bourgain--Kontorovich--Sarnak, Theorem 1.1, prove for fixed left and right \(K\)-Fourier modes \(n,k\)

\[
 \sum_{\|\gamma\|<T}e^{2in\theta_1(\gamma)}e^{2ik\theta_2(\gamma)}
 =\mathbf1_{n=0}\mathbf1_{k=0}c_0T^2
  +\sum_jc_j(n,k)T^{2s_j}+O_{n,k}(T^{4/3}).
\]

On the adjoint orbit, affine height is quadratic in the group norm, so \(T^{4/3}\) again corresponds to \(X^{2/3}\). Recovering a sharp angular interval requires summing many Fourier modes, while the theorem's implied constants depend on \(n,k\). It is not a uniform desmoothing theorem for this chamber.

\subsection{Smooth spectral counting}\label{koide-lit-smooth-spectral-counting}

Liu--Sarnak's paper is the closest source with a square-root-plus-gap exponent. Its global arithmetic setup and the hypotheses of its analytic proposition must be separated.

The introduction, author PDF p.~2, fixes an indefinite anisotropic integral ternary form \(f\), a nonzero level \(t\), and assumes \(t\det(f)\) is square-free. Thus the paper's affine-sieve theorem, as globally stated, does not cover \(t\det(f)=-108\). Equations (2.3)--(2.4), author PDF p.~4, fix one arithmetic orbit \(\mathcal O\) and a smooth \(K\)-invariant radial weight \(F_T\) satisfying

\[
 F_T=1\quad(|x|\leq T/c_0),
 \qquad
 F_T=0\quad(|x|\geq c_0T),
\]

with \(c_0>1\) fixed. In Fermi coordinates, Proposition 3.2 uses a smooth even function equal to one for \(|\rho|\leq R-1\), zero for \(|\rho|\geq R+1\), and with uniformly bounded derivatives. Unlike the global sieve statement, Proposition 3.2 is formulated for any variable subgroup \(\Gamma<\mathrm{SL}_2(\mathbb R)\) satisfying its uniform local condition (3.1), the spectral gap (2.7), and containing the fixed primitive cyclic group \(\Gamma_\xi\). It proves

\[
 H_R(z)=\frac{X(R)}{\operatorname{Area}(\Gamma\backslash\mathbb H)}
 +O_\varepsilon(e^{(1/2+\theta+\varepsilon)R}),
 \qquad X(R)\asymp e^R.
\]

Equation (3.60) has \(T=\sinh R\), and equation (3.65) transfers this to the one-sheeted affine orbit. Thus the smooth affine error is

\[
 O_\varepsilon(T^{1/2+\theta+\varepsilon}).
\]

Equation (2.7) assumes \(\lambda_1\geq1/4-\theta^2\). The paper invokes Jacquet--Langlands and Kim--Sarnak to take \(\theta=7/64\), giving

\[
 O_\varepsilon(T^{39/64+\varepsilon})
\]

unconditionally in its congruence-arithmetic setting; \(\theta=0\), and
hence square root, is conditional on the generalized
Selberg/Ramanujan-at-infinity conjecture for the relevant
Jacquet--Langlands transfers.

The numerical input itself is Kim--Sarnak, Appendix 2, Proposition 2, printed p.~176. It applies to an arbitrary cuspidal automorphic representation of \(\mathrm{GL}_2/\mathbb Q\): for a spherical real component, \(|\Re\mu_{j,\infty}|\leq7/64\), independently of finite conductor. Jacquet--Langlands preserves the real component because the quaternion algebra is split at infinity. Chu--Li, proof of Corollary 1.2, equation (5.2), printed p.~5126, records the resulting

\[
 \lambda_1\geq\frac{975}{4096}
\]

for every congruence arithmetic Fuchsian group. Thus the nonmaximal or
non-Eichler structure of the exact finite order does not invalidate
\(s\leq39/64\).  It changes the finite compact-open subgroup and hence which
automorphic representations have fixed vectors and the dimensions of those
fixed-vector spaces; it does not change this archimedean bound. This
validates the spectral-gap input for the exact lattice. It does not provide
sharp-boundary control or force the periods of possible
complementary-spectrum eigenfunctions to vanish.

For the exact order lattice \(\Gamma=\mathcal O^1\), each primitive trace-four conjugacy class supplies the cyclic subgroup in Proposition 3.2, and there are only finitely many such classes. Hence the proposition and equation (3.65) do give an unconditional \(39/64\) estimate after summing the corresponding \textbf{smooth standard-radial} counts. The square-free hypothesis is not an obstruction to this isolated analytic application.

This still does not desmooth to the target. A transition of fixed width in \(R\) is a fixed multiplicative annulus in \(T\), whose invariant volume is of order \(T\), not \(T^{1/2}\). Shrinking the transition makes the derivative/Sobolev constants in Proposition 3.2 grow, and the paper gives no uniform estimate that balances this loss at square-root scale. It also has no sharp coordinate chamber, and its radial height is not \(4a+17n\). These geometric failures remain even after the exact lattice and finite class sum have been supplied.

The same distinction applies to any informal statement that a tempered spectral component has square-root size. Duke--Rudnick--Sarnak's discussion on p.~155 estimates individual tempered frequencies; it does not control the complete spectral sum evaluated at a point with a shrinking approximation to a sharp characteristic function.

\subsection{General effective sector and full-lattice theorems}\label{koide-lit-general-effective-sector-and-full-lattice-theorems}

\subsubsection{Gorodnik--Nevo's explicit generic exponent}\label{koide-lit-gorodniknevos-explicit-generic-exponent}

Gorodnik--Nevo Definition 3.1 defines

\[
 n_e(2)=1,
 \qquad
 n_e(p)=\text{the least even integer at least }p/2\quad(p>2).
\]

Their Section 8 considers

\[
 S_t(\Phi,\Psi)
 =\{g:\log\|gv_0\|\leq t\}\cap\Phi A^+\Psi.
\]

Theorem 8.1, author PDF p.~34, assumes \(L^{p+}\) matrix coefficients and Lipschitz-well-rounded fixed angular sets; Corollary 8.3, p.~35, passes to an orbit when \((\Gamma\cap H)\backslash H\) is compact. The exponent is

\[
 \zeta=\frac{1}{2n_e(p)(1+3\dim G)}.
\]

For \(\dim G=3\),

\[
 \zeta=\frac1{20n_e(p)}.
\]

Even tempered spectrum has \(p=2\), \(n_e(p)=1\), and hence absolute error \(O(X^{19/20+\epsilon})\). For the exact congruence arithmetic lattice, Jacquet--Langlands and Kim--Sarnak give \(s\leq39/64\). The corresponding matrix coefficients are \(L^{p+}\) for \(p=64/25\), so \(n_e(p)=2\) and the theorem gives only \(O(X^{39/40+\epsilon})\).

This numerical calculation is an upper benchmark for the published generic polar-sector theorem, not a claim that Corollary 8.3 literally counts \(R(X)\). A fixed polar sector contains every radial parameter in fixed polar directions. An arbitrary affine linear cone and its finite-radius polyhedral height need an exact polar decomposition or a quantitatively controlled symmetric-difference argument; all integral orbits must then be included. None of those identifications appears in the theorem statement.

Gorodnik--Nevo themselves note after Corollary 5.3, author PDF p.~23, that their sharp-set counting mechanism remains worse than square-root volume error even for tempered spectrum, whereas square-root quality is available there only for a smooth weighted problem.

\subsubsection{Browning--Gorodnik and subsequent full-lattice work}\label{koide-lit-browninggorodnik-and-subsequent-full-lattice-work}

Browning--Gorodnik Theorem 2.1 gives, for a Euclidean ball \(B_H\) in one congruence orbit,

\[
 |\Gamma_\ell y\cap B_H|
 =\frac{m_{Z_{y,\ell}}(Z_{y,\ell})}{m_{X_\ell}(X_\ell)}m_Y(B_H)
 +O\!\left(\ell^{\dim H+\dim G}m_Y(B_H)^{1-\varrho}\right)
\]

for some unspecified \(\varrho>0\). Corollaries 2.5--2.6 sum adelic and integral orbits and retain the Brauer/Kottwitz density, at the cost of \(\ell^{\dim H+2\dim G}\). This is a useful exact-lattice precedent, but its archimedean set is the full norm ball, not the present sector, and no square-root exponent is asserted.

Cao--Huang Theorem 1.5 proves arithmetic purity of the Hardy--Littlewood property for an anisotropic ternary affine quadric with \(-m\det(q)\) nonsquare. Its Theorem 3.4 invokes Nevo--Sarnak, Gorodnik--Nevo, and Browning--Gorodnik to furnish an unspecified positive exponent, denoted \(\beta_X\) there with \(X\) meaning the affine variety; write it here as \(\beta_{\mathscr X}>0\) to distinguish it from the height variable. It is not a sharp-sector remainder theorem.

The current Cao--Huang--Zhang preprint, arXiv:2509.23320, Setting 2.1 and Theorem 2.1, gives a congruence-uniform power saving

\[
 O\!\left(\mathbf N(\mathfrak l)^\kappa
 \operatorname{vol}(\mathbb B_S(N))^{1-\delta}\right)
\]

for standard product Euclidean/max height balls and a compact piecewise-smooth factor. The exponent \(\delta\) is unspecified, and the unbounded real factor in the target is neither compact nor their standard Euclidean ball. This preprint extends the general effective framework but does not improve the conclusion of the direct Benoist--Oh argument in \cref{koide-lit-definition-13.1-for-the-actual-b_x,koide-lit-finite-orbit-passage-and-main-term-normalization}.

\subsection{Results in different arithmetic regimes}\label{koide-lit-results-in-different-arithmetic-regimes}

\subsubsection{Square discriminant and divergent stabilizer}\label{koide-lit-square-discriminant-and-divergent-stabilizer}

Oh--Shah Theorem 1.2 assumes

\[
 \operatorname{vol}((H\cap\Gamma)\backslash H)=\infty
\]

and obtains a smoothed \(cT\log T+O(T)\) count. Its sharp version has only a logarithmic relative error. Theorem 6.1 allows an open cone and arbitrary norm, but retains the logarithmic main-factor characteristic of its divergent-stabilizer regime; that regime is excluded here.

Xu--Zhang, \emph{Quarterly Journal of Mathematics} 74 (2023), 659--685, assumes \(-m\det(f)\) is a square over the number field. Huang, arXiv:2412.03350, Theorem 1.1, assumes the same split condition and a fixed smooth compact weight; its principal error is \(B(\log B)^{1983/1984}\) in a \(B\log B\) count. Zhang, arXiv:2603.21080, Section 1.1 and Theorem 1.1, again assumes \(-m\det(q)\in(k^\times)^2\) and proves a smooth \(K_0\)-invariant expansion with an error \(O(T^{l_1+l_2-\delta})\), \(0<\delta<1/2\); Remark 1.2 obtains an authentic count for \(\Gamma=\mathrm{SL}_2(\mathcal O_k)\). These are important current results, but all lie in the rationally split torus regime excluded by \(108\notin\mathbb Q^{\times2}\), and none uses the target polyhedral height.

\subsubsection{Varying level, isotropic forms, and special rotational quadrics}\label{koide-lit-varying-level-isotropic-forms-and-special-rotational-quadrics}

Huang, arXiv:2412.05903, studies \(F(x)=m_0N\) with \(N=p_0^{2h}\to\infty\), a smooth compactly supported real weight \(w(x/\sqrt N)\), and prescribed local conditions. Theorems 1.1--1.2 have a \(\sqrt N/(\log N)^A\) error. This is a Linnik-type varying-level theorem, not a fixed affine shell.

Li--Liu, arXiv:2501.15033v2 (18 April 2025), Theorems 1.1--1.2, prove new almost-prime conclusions for anisotropic ternary equations. Their global assumption remains \(t\det(f)\) square-free. The quantitative counting input in Lemma 2.3, pp.~5--6, is

\[
 O_\varepsilon(d^{1+\varepsilon}T^{1/2+\theta+\varepsilon})
\]

for the same kind of smooth radial orbit weight inherited from Liu--Sarnak. It is a current application of the \(7/64\) gap, not a new unsmoothed sector theorem.

Lowry-Duda's thesis, arXiv:1704.02376, Theorem 5.5.4, treats the special rotational form \(x_1^2+x_2^2=x_3^2+h\) in a full Euclidean ball. Written in ordinary radius \(T\), its three-dimensional sharp error is \(O(T^{21/22+\epsilon})\). The homogeneous form is \(\mathbb Q\)-isotropic, and neither the lattice nor the height sector is the target one.

Results on values of irrational ternary forms, shrinking sets while the discriminant varies, or positive-definite ternary representation numbers answer different problems. Lindelöf or generalized Lindelöf hypotheses occurring there do not yield the present fixed-level pointwise sector bound.

\subsection{Spectral-gap, GRH, smoothing, and averaging taxonomy}\label{koide-lit-spectral-gap-grh-smoothing-and-averaging-taxonomy}

\begin{longtable}[]{@{}
  >{\raggedright\arraybackslash}p{0.31\linewidth}
  >{\raggedright\arraybackslash}p{0.31\linewidth}
  >{\raggedright\arraybackslash}p{0.31\linewidth}@{}}
\toprule\noalign{}
\begin{minipage}[b]{\linewidth}\raggedright
Classification
\end{minipage} & \begin{minipage}[b]{\linewidth}\raggedright
Source result
\end{minipage} & \begin{minipage}[b]{\linewidth}\raggedright
Why it does not prove the target
\end{minipage} \\
\midrule\noalign{}
\endhead
\bottomrule\noalign{}
\endlastfoot
Unconditional, exact sharp family & Benoist--Oh Theorem 13.2 plus the direct Definition 13.1 verification & Gives only \(O(X^{1-\eta})\) with \(\eta>0\) unspecified \\
Unconditional, exact norm/cone, full lattice & Gorodnik--Oh--Shah Theorem 1.6, with the exact hypotheses proved in \cref{koide-count-lem-genuine-norm,koide-count-lem-gos-hypotheses} & Qualitative \(C_0X+o(X)\) only \\
Unconditional, sharp radial closed-geodesic count & Good / Chatzakos--Petridis & \(O(X^{2/3})\), different height, and main expression includes small-spectrum terms \\
Unconditional, smooth radial count & Liu--Sarnak Proposition 3.2 with \(\theta=7/64\), summed over the exact finite conjugacy-class decomposition & \(39/64\) for the exact lattice, but fixed multiplicative smoothing and a different radial height \\
Selberg/Ramanujan conditional & Liu--Sarnak with \(\theta=0\) & Square-root only for the smooth radial problem; still not the exact count \\
GRH/Lindelöf conditional & None among the source versions in \cref{koide-lit-primary-references-and-exact-locators} proves this fixed sharp-sector estimate & The occurrences listed there concern varying levels/forms, shrinking-set equidistribution, or lower/Omega results \\
Smoothed & Liu--Sarnak; Oh--Shah; Huang I; Zhang 2026 & Smoothing or split arithmetic changes the theorem; no square-root desmoothing supplied \\
Averaged/mean square & Chatzakos--Petridis Theorems 1.2--1.3; Voskou Theorem 4; Lekkas--Petridis Theorem 1.5 & Does not imply a uniform pointwise estimate at every \(X\) \\
Omega/lower-bound & Chatzakos (2019), including a cocompact geodesic-average \(\Omega(X^{1/2}\log\log\log X)\) & Evidence for scale, not an upper bound \\
Isotropic or split & Oh--Shah; Xu--Zhang; Huang I; Zhang 2026; Lowry-Duda & Different stabilizer, quotient, main order, form, or level \\
Varying level & Huang II & Counts \(F=m_0p_0^{2h}\), not \(F=-2\) \\
\end{longtable}

In particular, ``assuming GRH'' is not a substitute for the missing theorem. The conditional hypothesis that occurs naturally in the closest smooth spectral result is the Selberg/Ramanujan eigenvalue conjecture, and even it only reaches square root before the sharp-boundary unsmoothing problem is addressed.

\subsection{Citation-ready conclusion}\label{koide-lit-citation-ready-conclusion}

Subject to the bounded-search qualification in
\cref{koide-lit-audit-limitations}, the following paragraph records the
conclusion supported by the enumerated source versions:

\begin{quote}
Among the source versions enumerated in \cref{koide-lit-primary-references-and-exact-locators}, none implies the unconditional sharp estimate \(R(X)-C_0X=O_\varepsilon(X^{1/2+\varepsilon})\) for the fixed ternary quadric \(a^2-3q^2-18n^2=-2\) in the chamber \(a>0\), \(q,n\geq0\), \(a+4n>q\), with height \(4a+17n\leq X\). Gorodnik--Oh--Shah, Theorem 1.6, applies to the exact admissible generic cone and an arbitrary norm---with those exact hypotheses proved in \cref{koide-count-lem-genuine-norm,koide-count-lem-gos-hypotheses}---but gives only a qualitative asymptotic. Benoist--Oh, Definition 13.1 and Theorem 13.2, applies after the direct effective-well-roundedness verification in \cref{koide-lit-definition-13.1-for-the-actual-b_x} and, after summing the finitely many integral spin-group orbits, yields \(R(X)=C_0X+O(X^{1-\eta})\) for some unspecified \(\eta>0\), not \(\eta=1/2\). In the closest sharp closed-geodesic problem, Good's theorem, as restated and reproved by Chatzakos--Petridis, Theorem 1.1, gives \(O(X^{2/3})\) only after all small-eigenvalue terms \(X^{s_j}\) are included; Chatzakos--Petridis, Conjecture 5.7, states the \(O_\varepsilon(X^{1/2+\varepsilon})\) improvement as a conjecture. Liu--Sarnak, Proposition 3.2, applies class by class to the exact congruence quaternionic lattice and proves the smooth radial bound \(O_\varepsilon(X^{1/2+\theta+\varepsilon})\), unconditionally with \(\theta=7/64\) and conditionally with \(\theta=0\), but its fixed-width smoothing and radial height do not match the exact sharp count. Thus complementary-spectrum terms and sharp nonradial boundary control remain separate unresolved inputs.
\end{quote}

\subsection{Primary references and exact locators}\label{koide-lit-primary-references-and-exact-locators}

\begin{enumerate}
\def\labelenumi{\arabic{enumi}.}

\item
  Yves Benoist and Hee Oh, ``Effective equidistribution of \(S\)-integral points on symmetric varieties,'' \emph{Annales de l'Institut Fourier} \textbf{62} (2012), 1889--1942. \href{https://doi.org/10.5802/aif.2738}{DOI}, \href{https://www.numdam.org/item/AIF_2012__62_5_1889_0/}{official article/PDF}, \href{https://arxiv.org/abs/0706.1621}{arXiv}. Standing setting pp.~1892--1893; finite orbit statement p.~1892; Theorem 1.3 p.~1894; Theorem 12.5 p.~1929; Definition 13.1 p.~1931; Theorem 13.2 pp.~1931--1932; Proposition 14.2 p.~1933; ternary Example (A) p.~1896.
\item
  Alexander Gorodnik, Hee Oh, and Nimish Shah, ``Integral points on symmetric varieties and Satake compactifications,'' \emph{American Journal of Mathematics} \textbf{131} (2009), 1--57. \href{https://doi.org/10.1353/ajm.0.0034}{DOI}, \href{https://arxiv.org/abs/math/0610497}{arXiv}, \href{https://gauss.math.yale.edu/~ho2/doc/satake_ajm.pdf}{author PDF}. Admissible/generic definitions and Theorem 1.6 pp.~4--5; Theorem 1.18 pp.~8--9.
\item
  Alexander Gorodnik, Hee Oh, and Nimish Shah, ``Strong wavefront lemma and counting lattice points in sectors,'' \emph{Israel Journal of Mathematics} \textbf{176} (2010), 419--444. \href{https://doi.org/10.1007/s11856-010-0035-8}{DOI}, \href{https://arxiv.org/abs/0708.3188}{arXiv}. Theorem 1.13, author PDF pp.~4--5.
\item
  Alexander Gorodnik and Amos Nevo, ``Counting lattice points,'' \emph{Journal für die reine und angewandte Mathematik} \textbf{663} (2012), 127--176. \href{https://doi.org/10.1515/CRELLE.2011.096}{DOI}, \href{https://arxiv.org/abs/0903.1515}{arXiv}, \href{https://www.math.uzh.ch/gorodnik/papers/GN2-final.pdf}{author PDF}. Theorem 1.9 pp.~4--5; Definition 3.1 p.~13; Theorem 4.5 pp.~17--18; Theorem 8.1 and Corollary 8.3 pp.~34--35; Proposition 8.4 p.~35.
\item
  William Duke, Zeév Rudnick, and Peter Sarnak, ``Density of integer points on affine homogeneous varieties,'' \emph{Duke Mathematical Journal} \textbf{71} (1993), 143--179. \href{https://doi.org/10.1215/S0012-7094-93-07107-4}{DOI}, \href{https://www.math.ucla.edu/~wdduke/preprints/density.pdf}{author PDF}. Theorem 1.2 p.~145; Example 1.5 pp.~146--147; discussion of individual tempered frequencies p.~155; Theorem 3.5 p.~161.
\item
  Alex Eskin and Curtis McMullen, ``Mixing, counting, and equidistribution in Lie groups,'' \emph{Duke Mathematical Journal} \textbf{71} (1993), 181--209. \href{https://doi.org/10.1215/S0012-7094-93-07108-6}{DOI}, \href{https://people.math.harvard.edu/~ctm/papers/home/text/papers/mixing/mixing.pdf}{author PDF}. Theorem 1.4, author PDF p.~3; Theorems 6.1 and 6.4 pp.~30--31; Open Question 1 p.~35.
\item
  Jianya Liu and Peter Sarnak, ``Integral points on quadrics in three variables whose coordinates have few prime factors,'' \emph{Israel Journal of Mathematics} \textbf{178} (2010), 393--426. \href{https://doi.org/10.1007/s11856-010-0069-y}{DOI}, \href{https://web.math.princeton.edu/facultypapers/sarnak/few-final.pdf}{author PDF}. Arithmetic hypothesis p.~2; equations (2.3)--(2.8) pp.~4--5; Proposition 3.2 pp.~10--11; equations (3.60) and (3.65) pp.~15--16.
\item
  Anton Good, \emph{Local Analysis of Selberg's Trace Formula}, Lecture Notes in Mathematics 1040, Springer, 1983. \href{https://doi.org/10.1007/BFb0073074}{DOI}. Theorem 4, p.~116. The theorem statement and a new proof were checked through Chatzakos--Petridis (2016) because the monograph text was not directly available in this audit.
\item
  Dimitrios Chatzakos and Yiannis N. Petridis, ``The hyperbolic lattice point problem in conjugacy classes,'' \emph{Forum Mathematicum} \textbf{28} (2016), 981--1003. \href{https://doi.org/10.1515/forum-2015-0102}{DOI}, \href{https://arxiv.org/abs/1504.01307}{arXiv}, \href{https://discovery.ucl.ac.uk/1477487/7/Petridis_forum-2015-0102.pdf}{author manuscript}. Equations (1.4)--(1.6) and Theorem 1.1 p.~3; Theorems 1.2--1.3 pp.~3--4; Conjecture 5.7 p.~22.
\item
  Dimitrios Chatzakos, ``Mean value results and \(\Omega\)-results for the hyperbolic lattice point problem in conjugacy classes,'' \emph{Revista Matemática Iberoamericana} \textbf{35} (2019), 1123--1152. \href{https://doi.org/10.4171/RMI/1080}{DOI}, \href{https://arxiv.org/abs/1610.01462}{arXiv}. Mean-value and Omega theorems, not a pointwise upper bound.
\item
  Jean Bourgain, Alex Kontorovich, and Peter Sarnak, ``Sector estimates for hyperbolic isometries,'' \emph{Geometric and Functional Analysis} \textbf{20} (2010), 1175--1200. \href{https://doi.org/10.1007/s00039-010-0092-5}{DOI}, \href{https://arxiv.org/abs/1001.4541}{arXiv}. Theorem 1.1, p.~2.
\item
  Hee Oh and Nimish Shah, ``Limits of translates of divergent geodesics and integral points on one-sheeted hyperboloids,'' \emph{Israel Journal of Mathematics} \textbf{199} (2014), 915--931. \href{https://doi.org/10.1007/s11856-013-0063-2}{DOI}, \href{https://arxiv.org/abs/1104.4988}{arXiv}. Theorem 1.2 p.~3; Corollary 1.4 pp.~3--4; Theorem 6.1 pp.~13--14.
\item
  Tim D. Browning and Alexander Gorodnik, ``Power-free values of polynomials on symmetric varieties,'' \emph{Proceedings of the London Mathematical Society} \textbf{114} (2017), 1044--1080. \href{https://doi.org/10.1112/plms.12030}{DOI}, \href{https://arxiv.org/abs/1606.06342}{arXiv}. Theorem 2.1 p.~11; Corollaries 2.5--2.6 pp.~24--25.
\item
  Yang Cao and Zhizhong Huang, ``Arithmetic purity of the Hardy--Littlewood property and geometric sieve for affine quadrics,'' \emph{Advances in Mathematics} \textbf{398} (2022), 108236. \href{https://doi.org/10.1016/j.aim.2022.108236}{DOI}, \href{https://arxiv.org/abs/2003.07287}{arXiv}. Theorem 1.5 p.~4; Proposition 2.6 pp.~12--13; Theorem 3.4 pp.~21--22.
\item
  Yang Cao, Zhizhong Huang, and Runlin Zhang, ``Effective equidistribution, arithmetic purity of strong approximation, and geometric sieve for affine quadrics,'' \href{https://arxiv.org/abs/2509.23320}{arXiv:2509.23320}. Setting 2.1 and Theorem 2.1; arXiv HTML retrieved 4 September 2026.
\item
  Mikhail Borovoi, ``On representations of integers by indefinite ternary quadratic forms,'' \emph{Journal of Number Theory} \textbf{90} (2001), 281--293. \href{https://doi.org/10.1006/jnth.2001.2662}{DOI}, \href{https://arxiv.org/abs/math/0006141}{arXiv}, \href{https://www.math.tau.ac.il/~borovoi/papers/reprfin.pdf}{author PDF}. Theorem 1.4.1 and Proposition 1.5.1; qualitative local/Brauer--Manin coefficient, no error exponent.
\item
  Fei Xu and Runlin Zhang, ``Counting integral points on indefinite ternary quadratic equations over number fields,'' \emph{Quarterly Journal of Mathematics} \textbf{74} (2023), 659--685. \href{https://doi.org/10.1093/qmath/haac039}{DOI}, \href{https://arxiv.org/abs/2103.10707}{arXiv}. Square-discriminant hypothesis in the main theorem.
\item
  Zhizhong Huang, ``Quantitative strong approximation for ternary quadratic
  forms I.''

  \href{https://arxiv.org/abs/2412.03350}{arXiv:2412.03350}; Theorem 1.1;
  square-discriminant, smooth fixed-weight \(B\log B\) count.
\item
  Zhizhong Huang, ``Quantitative strong approximation for ternary quadratic
  forms II.''

  \href{https://arxiv.org/abs/2412.05903}{arXiv:2412.05903}; Theorems 1.1--1.2;
  varying level \(m_0p_0^{2h}\) and smooth compact weight.
\item
  Runlin Zhang, ``Effective count of integer points on ternary affine quadrics and effective equidistribution,'' \href{https://arxiv.org/abs/2603.21080}{arXiv:2603.21080}. Section 1.1 and Theorem 1.1; split-stabilizer smooth count; Remark 1.2 on the authentic radial count.
\item
  Marios Voskou, ``Refined Counting of Geodesic Segments in the Hyperbolic
  Plane.''

  \href{https://arxiv.org/abs/2407.03134}{arXiv:2407.03134v3}; Theorems 3--4
  and Conjecture 1, pp.~3--4.
\item
  Dimitrios Lekkas and Yiannis N. Petridis, ``A relative trace formula and counting geodesic arcs in the hyperbolic plane,'' \href{https://arxiv.org/abs/2509.12902}{arXiv:2509.12902}, version dated 24 August 2026. Theorem 1.2, Theorem 1.5, and Conjecture 1.6.
\item
  Henry H. Kim, ``Functoriality for the exterior square of \(\mathrm{GL}_4\) and the symmetric fourth of \(\mathrm{GL}_2\),'' with Appendix 2 by Henry H. Kim and Peter Sarnak, ``Refined estimates towards the Ramanujan and Selberg conjectures,'' \emph{Journal of the American Mathematical Society} \textbf{16} (2003), 139--183. \href{https://doi.org/10.1090/S0894-0347-02-00410-1}{DOI}, \href{https://www.ams.org/journals/jams/2003-16-01/S0894-0347-02-00410-1/S0894-0347-02-00410-1.pdf}{official AMS PDF}. Appendix 2, Proposition 2, printed p.~176: \(|\Re\mu_{j,\infty}|\leq7/64\) for a spherical real component of any cuspidal \(\mathrm{GL}_2/\mathbb Q\) representation, and the displayed consequence \(\lambda_1\geq975/4096\).
\item
  Michelle Chu and Han Li, ``Small generators of cocompact arithmetic Fuchsian groups,'' \emph{Proceedings of the American Mathematical Society} \textbf{144} (2016), 5121--5127. \href{https://doi.org/10.1090/proc/13177}{DOI}, \href{https://arxiv.org/abs/1603.00926}{arXiv}. Section 3.1, p.~5122, defines congruence arithmetic Fuchsian groups; proof of Corollary 1.2, equation (5.2), p.~5126, records \(\lambda_1\geq975/4096\) for every such group via Jacquet--Langlands and Kim--Sarnak.
\item
  Hervé Jacquet and Robert P. Langlands, \emph{Automorphic Forms on \(\mathrm{GL}(2)\)}, Lecture Notes in Mathematics 114, Springer, 1970. \href{https://doi.org/10.1007/BFb0058988}{DOI}, \href{https://publications.ias.edu/sites/default/files/automorphic-forms-on-gl2_rpl_9.pdf}{IAS-hosted text}. Chapter III, especially §16, supplies the global quaternionic transfer; at split places, including the real place here, the local component is preserved.
\item
  John Voight, \emph{Quaternion Algebras}, Graduate Texts in Mathematics 288, Springer, 2021. \href{https://doi.org/10.1007/978-3-030-56694-4}{DOI}, \href{https://jvoight.github.io/quat-book.pdf}{open author version}. Proposition 15.5.2 (maximal superorders); §23.1.3 (Eichler orders at ramified primes); equations (27.6.6)--(27.6.7) (compact-open finite adelic order/unit group); Theorem 28.2.10 (strong approximation for indefinite \(B^1\)); Definition 38.7.2 (principal congruence groups).
\item
  Jiamin Li and Jianya Liu, ``Anisotropic quadratic equations in three
  variables.''

  \href{https://arxiv.org/abs/2501.15033}{arXiv:2501.15033v2};
  Theorems 1.1--1.2, pp.~2--3; Lemma 2.3, pp.~5--6. The global form-level
  hypothesis is \(t\det(f)\) square-free; the quantitative input is a smooth
  radial, congruence-orbit estimate with
  \[
   O_\varepsilon\!\left(
     d^{1+\varepsilon}T^{1/2+\theta+\varepsilon}
   \right).
  \]
\item
  David Lowry-Duda, \emph{On Some Variants of the Gauss Circle Problem}, \href{https://arxiv.org/abs/1704.02376}{arXiv:1704.02376}. Theorem 5.5.4, thesis pp.~97--98.
\end{enumerate}

\subsection{Audit limitations}\label{koide-lit-audit-limitations}

The claim-by-claim source versions and locators used in this bounded audit
are enumerated in
\cref{koide-lit-primary-references-and-exact-locators}.  The canonical tree
does not preserve byte-identical snapshots or retrieval hashes for those
external publications, so this file does not assert a reproducible local
source corpus.  Good's 1983 book was the sole exception to direct full-text
access: its theorem locator and statement were checked through the explicit
restatement and independent proof in Chatzakos--Petridis (2016) and the
version-dated Lekkas--Petridis treatment.

The statement that no theorem was located is a result about the source
versions enumerated here, not a proof that no unpublished or differently
indexed result exists.  It does establish that none of the enumerated works
may be cited as a proof of the exact unconditional square-root discrepancy.



\section{Interpretation and limits}

Five logically distinct facts should not be merged.  First, the Koide equation is a null-cone equation and has an exact determinant and phase model.  Second, the defect-one integral shell admits Lorentzian reflection dynamics, unique reduced roots, exact ternary coordinates, and explicit Pell families on its chamber faces.  Third, its HNF ideal data and positive-unit powers give an exact parametrization, inverse, multiplicity-one statement, sharp floor formula, and shifted Hecke-phase identity.  Fourth, homogeneous dynamics and arithmetic invariant measures give the explicit main coefficient, while direct effective well-roundedness and Benoist--Oh give an unspecified power saving.  Fifth, the exact quaternion shell bijection produces finitely many primitive hyperbolic conjugacy classes, but the radial Huber height is not the unchanged Koide height.  The smooth radial $39/64$ and sharp radial $2/3$ theorems therefore do not prove the open pointwise five-wall square-root estimate.  None of these statements explains why measured charged-lepton masses should obey the Koide relation.

The $S^6$ occurrence is still weaker: Proposition~\ref{prop:S6-dictionary} identifies coefficient spaces, not global moduli.  Corollary~\ref{cor:ambient-only} shows that a point of the Koide cone gives a singular real period matrix for every choice of real parts, whereas \eqref{eq:radial-determinant-asymptotic} places the recorded positive radial cusp matrices on the nonzero-determinant side for all sufficiently large $T$.  The full Koide cone is the discriminant hypersurface of the ambient coefficient model; its nonnegative-mass image lies in the rank-one boundary of the positive-semidefinite cone.  Neither is a family of smooth fibres supplied by the manuscript.

\bibliographystyle{apalike}
\bibliography{references}
\end{document}
